\chapter{Observational Problems}

\section*{Observation --- planetarium}

\probhead{15.1}{Earth}{2011}{Katowice, Poland}{OBS-PLAN 1}{25}
% subject: Chapter 1 (formerly Problem 1.42)
% source id: 2011_OBS-PLAN_1   tag: nova/Moon/comet marking, horizon check
The sky shown in the planetarium is stationary, and the observer is on the
surface of the Earth. Visible on the sky are: a comet, the Moon and a nova of
about 2nd magnitude. From the 11th minute, a grid representing horizontal
coordinates will be projected on the sky, and will remain on until the end of
the question. You are given a map of the sky for epoch J2000.0, using a polar
projection with a linear scale in declination, covering stars down to about
5th magnitude.

Note: Azimuth is counted from $0^\circ$ to $360^\circ$ starting at S through
W, N, E.

\begin{description}
\item[(A)] On the map of the sky, mark (with a cross) and label the nova
  (mark it ``N'') and the Moon (mark it with a Moon symbol) and draw the
  shape and position of the comet.
\item[(B)] In the table below, circle only those objects which are above the
  astronomical horizon. Note: you will lose 1 point for every incorrect
  answer.

  \begin{center}
  \begin{tabular}{lll}
  M20 -- Trifid Nebula & $o$ Cet -- Mira & $\delta$ CMa -- Wezen \\
  $\alpha$ Cyg -- Deneb & M57 -- Ring Nebula & $\beta$ Per -- Algol \\
  $\delta$ Cep -- Alrediph & $\alpha$ Boo -- Arcturus &
    M44 -- Praesepe (Beehive Cluster) \\
  \end{tabular}
  \end{center}
\item[(C)] When the coordinate grid is visible, mark on the map the northern
  part of the local meridian (from the zenith to the horizon) and the
  ecliptic north pole (with a cross and marked ``P'').
\item[(D)] For the displayed sky, give the:
  geographical latitude of the observer $\varphi$;
  Local Sidereal Time $\theta$;
  time of year, by circling the calendar month:
  Jan, Feb, Mar, Apr, May, Jun, Jul, Aug, Sep, Oct, Nov, Dec.
\item[(E)] Give the names of the objects whose approximate horizontal
  coordinates are:
  azimuth $A_1 = 45^\circ$ and altitude $h_1 = 58^\circ$;
  azimuth $A_2 = 278^\circ$ and altitude $h_2 = 20^\circ$.
  (If you can, use Bayer designations, IAU abbreviations and Messier numbers
  or English or Latin names.)
\item[(F)] Give the horizontal coordinates (azimuth, altitude) of:
  Sirius ($\alpha$ CMa): $A_3$, $h_3$;
  the Andromeda Galaxy (M31): $A_4$, $h_4$.
\item[(G)] Give the equatorial coordinates ($\alpha$, $\delta$) of the star
  marked on the sky with a red arrow.
\end{description}

\ioaafig{2011_OBS-PLAN_1-s2.pdf}

\probhead{15.2}{Mars}{2011}{Katowice, Poland}{OBS-PLAN 2}{25}
% subject: Chapter 1 (formerly Problem 1.43)
% source id: 2011_OBS-PLAN_2   tag: areographic latitude, culmination altitudes
Four consecutive days on the surface of Mars will be shown in the
planetarium. There is a Martian base visible on the horizon. During the
Martian daytime the sky will be slightly brightened. The moons of Mars and
the other planets will not be displayed. The local meridian will be
continuously visible on the sky. You are given a map of the sky for epoch
J2000.0, using a polar projection with a linear scale in declination,
covering stars down to about 5th magnitude.

Note: Azimuth is counted from $0^\circ$ to $360^\circ$ starting at S through
W, N, E.

\begin{description}
\item[(H)] Give the areographic (Martian) latitude of the observer
  $\varphi$.
\item[(I)] Give the altitudes of upper ($h_u$) and lower ($h_l$) culmination
  of:
  Pollux ($\beta$ Gem) and Deneb ($\alpha$ Cyg).
\item[(J)] Give the areocentric (Martian) declination of:
  Regulus ($\alpha$ Leo) and Toliman ($\alpha$ Cen).
\item[(K)] Sketch diagrams to illustrate your working in questions (I) and
  (J) above.
  [Two blank diagrams are provided, each showing the zenith, nadir, and the
  S--N horizon line through the observer.]
\item[(L)] On the map of the sky, mark (with a cross) and label (``M'') the
  Martian celestial North Pole.
\item[(M)] Give the azimuth $A$ of the observer as seen from the Martian
  base.
\item[(N)] Estimate the location of the base on Mars, and circle the
  appropriate description:
  \begin{description}
  \item[(a)] near the Equator
  \item[(b)] near the northern Tropic circle
  \item[(c)] near the northern Arctic circle
  \item[(d)] near the North Pole
  \end{description}
\item[(O)] A time axis is given showing the Martian year and the seasons in
  the northern hemisphere (spring, summer, autumn, winter). Mark the date
  represented by the planetarium display on the axis.
\end{description}

\ioaafig{2011_OBS-PLAN_2-s3.pdf}

\probhead{15.3}{Planetarium part I}{2014}{Suceava, Romania}{OBS-PLAN 1}{}
% subject: Chapter 1 (formerly Problem 1.44)
% source id: 2014_OBS-PLAN_1   tag: planetarium sky identification
The inside part of the planetarium round consists of 3 questions and takes
30 minutes.

\begin{description}
\item[Question 1.] The sky projected in the dome corresponds to Suceava
  (longitude $26^\circ 15'$), at 18:00 UT, on a certain day of a certain
  month. Two arcs of circle will now be projected. The arcs are segmented.
  Each segment represents an interval of some degrees. This number is not
  the same for each arc.
  \begin{description}
  \item[(a)] Identify each arc by circling the correct name (Equator /
    Meridian / Ecliptic) and give the angular size of each segment (in
    degrees).
    [Answer table: first arc, second arc --- name and segment size]
  \item[(b)] Estimate the local sidereal time of the sky you see in the
    dome.
  \item[(c)] Determine the month to which the projected sky would correspond
    at the given time. Fill in the box the number of the month (1 to 12).
  \end{description}
\item[Question 2.] The assistant will use a small red arrow pointer to point
  at objects in the sky; each object will be pointed at for 2 minutes.
  Locations of three Messier objects will be pointed at one by one. For each
  Messier object pointed at, fill in the boxes the Messier catalogue number
  and the number which indicates its type, as follows: 1 for galaxy, 2 for
  nebula, 3 for open cluster, 4 for globular cluster. Also, for each object,
  fill in the appropriate box the IAU abbreviation of the constellation
  where the object is located, using the provided table of constellations.
  [Answer table: 1st, 2nd, 3rd Messier object --- number, type,
  constellation]
\item[Question 3.] Three stars will be pointed at successively; each star
  will be pointed at for 2 minutes. Fill in the appropriate box the name of
  the star (or Bayer designation) and the number which indicates its type
  (1 for single, 2 for double). Also, for each star, fill in the appropriate
  box the IAU abbreviation of the constellation where the star is located,
  using the provided table of constellations.
  [Answer table: 1st, 2nd, 3rd star --- name, type, constellation]
\end{description}

\probhead{15.4}{Three missing bright stars in the Magelang sky}{2015}{Magelang, Indonesia}{OBS-PLAN 1}{60}
% subject: Chapter 1 (formerly Problem 1.45)
% source id: 2015_OBS-PLAN_1   tag: planetarium missing-star identification
The projected image in the planetarium is a simulated image of the Magelang
sky for July 28, 2015 at 21.00 (UTC+7). There are three stars missing in the
image. Mark the position of the three missing bright stars in Answer Sheet 1
by a cross ($\times$) and write their names beside the cross marks using
their IAU designations. For guidance you can also include a few stars that
you see in the planetarium and mark them with circles (o). The list of the
bright stars which are possibly missing is attached:

\begin{center}
\begin{tabular}{lll}
$\alpha$ Pav & $\alpha$ Lyr & $\alpha$ Aql \\
$\alpha$ Vir & $\alpha$ CrB & $\alpha$ Gru \\
$\alpha$ Ser & $\alpha$ Cen & $\alpha$ Cyg \\
$\alpha$ Oph & $\alpha$ Cru & $\alpha$ Boo \\
$\alpha$ PsA & $\alpha$ Sco & $\beta$ Aqr \\
$\beta$ Cru & $\beta$ Oph & $\beta$ Peg \\
$\gamma$ Sco & $\gamma$ Cyg & $\delta$ Pav \\
$\varepsilon$ Peg & $\varepsilon$ Boo & \\
\end{tabular}
\end{center}

(Duration: 6 minutes)

\ioaafig{2015_OBS-PLAN_1-s2.pdf}

\probhead{15.5}{Position of a bright supernova}{2015}{Magelang, Indonesia}{OBS-PLAN 2}{50}
% subject: Chapter 1 (formerly Problem 1.46)
% source id: 2015_OBS-PLAN_2   tag: planetarium position marking
Assume that a bright supernova appears in the sky on September 24, 2015 and
you watch it at 22.00 (UTC+7). The planetarium will now show the sky for
this date with the supernova. Mark the position of the supernova by a cross
($\times$) in Answer Sheet 2. You can mark other stars by circles (o) for
guidance. (Duration: 3 minutes)

\ioaafig{2015_OBS-PLAN_2-s2.pdf}

\probhead{15.6}{Rise and set times of a pointed star}{2015}{Magelang, Indonesia}{OBS-PLAN 3}{40}
% subject: Chapter 1 (formerly Problem 1.47)
% source id: 2015_OBS-PLAN_3   tag: rise/set time estimation
The planetarium will now display the Magelang sky for January 3, 2016, at
23.00 (UTC+7). The examiner will point to a star with an arrow. Estimate the
times at which this star will rise and set from Magelang on this date up to
20 minutes accuracy. (Duration: 6 minutes)

\probhead{15.7}{Eight historical supernovae}{2016}{Bhubaneswar, India}{OBS-PLAN OP1}{40}
% subject: Chapter 1 (formerly Problem 1.48)
% source id: 2016_OBS-PLAN_OP1   tag: identify supernova positions in projected sky
Eight well known historical supernovae will appear in the projected sky one
at a time (not necessarily in chronological order). You have to identify the
appropriate map (Map 1 / Map 2) where a particular supernova belongs and
mark it in the corresponding map with a `+' sign and write codes `S1' to
`S8' beside it.

Each supernova code will be projected on the dome for 10 seconds, followed
by appearance of the supernova for 60 seconds and then 20 seconds for you to
mark the answers.

\begin{description}
\item[(OP1.1)] For S1, S2, S3, S4 and S5, the projected sky corresponds to
  the sky as seen from Rio de Janeiro on the midnight of 21st May.
\item[(OP1.2)] For S6, S7 and S8, the projected sky corresponds to the sky
  as seen from Beijing on the midnight of 20th November. There will be a gap
  of two minutes after S5 for change over and adaptation to the new sky.
\end{description}

\textit{[Answer maps Map 1 and Map 2 were distributed separately and are not in the available source papers]}

\probhead{15.8}{Sky of another planet}{2016}{Bhubaneswar, India}{OBS-PLAN OP2}{10}
% subject: Chapter 1 (formerly Problem 1.49)
% source id: 2016_OBS-PLAN_OP2   tag: non-terrestrial sky, celestial sphere
We are now projecting the sky of another planet. The sky will be slowly
rotated for 5 minutes. Identify the visible celestial pole of this planet
and mark it with a `+' sign and label it as `P' on the appropriate map
(Map 1 / Map 2).

\textit{[Answer maps Map 1 and Map 2 were distributed separately and are not in the available source papers]}

\probhead{15.9}{Sky above Keszthely at midnight}{2019}{Keszthely, Hungary}{OBS-PLAN 1}{}
% subject: Chapter 1 (formerly Problem 1.50)
% source id: 2019_OBS-PLAN_1   tag: planetarium sky identification
This is the sky above Keszthely at midnight. The projected sky does not show
any Solar System objects.

\begin{description}
\item[(1.1)] There are 3 novae on the projected sky at 2nd magnitude. Mark
  their positions by circles on the star chart. (Please circle only 3 stars.
  If there are more than 3 circles, then each one in a wrong location will
  result in 1 point deduction.)
\item[(1.2)] The Messier objects have been removed from the star chart given
  to you. Mark all the Messier list globular clusters present in the
  projected sky on the star chart using crosses (X) and write the Messier
  number of each object near the cross marks.
\item[(1.3)] The projected sky corresponds to the second half of which month
  (at midnight CEST) in Keszthely? Circle the correct month.
  JAN / FEB / MAR / APR / MAY / JUN / JUL / AUG / SEP / OCT / NOV / DEC
\item[(1.4)] What is the local sidereal time? (To an accuracy of 15
  minutes.)
\item[(1.5)] List six zodiacal constellations which are partly or entirely
  visible. (Use the official IAU abbreviations or IAU designation. Every
  constellation named which is not visible in the projected sky will result
  in 1 point deduction.)
\end{description}

\ioaafig{2019_OBS-PLAN_1-s2.pdf}

\probhead{15.10}{Planetarium projected image 2}{2019}{Keszthely, Hungary}{OBS-PLAN 2}{}
% subject: Chapter 1 (formerly Problem 1.51)
% source id: 2019_OBS-PLAN_2   tag: planetarium sky identification
We are standing somewhere on the Earth. The projected sky does not show any
Solar System objects.

\begin{description}
\item[(2.1)] Determine the geographical latitude of this observing site. In
  which hemisphere is the site situated? N / S (Circle the right one.)
\item[(2.2)] Determine the azimuth of the 3 brightest stars on the projected
  sky. Azimuth is measured from North towards the East. Write the name of
  these stars in English or using their Bayer designation and their
  azimuths.
\item[(2.3)] Yellow $\times$ signs show the position of 3 comets. Which
  comet is closest to the ecliptic? (Circle the number: 1 / 2 / 3.)
\item[(2.4)] List nine constellations that contain circumpolar stars seen
  from the given observing site. (Use the official IAU abbreviations or IAU
  designation.)
\item[(2.5)] Mintaka ($\delta$ Orionis) is setting at this moment. How many
  hours earlier did it rise? (To an accuracy of 15 minutes.)
\end{description}

\probhead{15.11}{Planetarium projected image 3}{2019}{Keszthely, Hungary}{OBS-PLAN 3}{}
% subject: Chapter 1 (formerly Problem 1.52)
% source id: 2019_OBS-PLAN_3   tag: planetarium sky identification
For this view, we are now standing on the Moon. At this instant viewed, the
Earth is centrally eclipsing the Sun (see the red circle on the sky).
Consequently the Moon is at one of its nodes now. Assume the longitudinal
and latitudinal librations are exactly $0^\circ$ at this moment.

\begin{description}
\item[(3.1)] At the time of this observation, which season is it in Hungary?
  (Circle the correct answer.) Spring / Summer / Autumn / Winter
\item[(3.2)] There is a yellow circle on the projected sky (next to the red
  circle), which denotes minor planet Juno, which is at a distance of
  exactly 3 au from the Sun at this moment. Estimate its distance to the
  Moon at this instant. (Rounded to the nearest integer in units of million
  km.) Assume all orbits to be circular.
\item[(3.3)] Approximately how much time (in Earth days) after the projected
  event will \ldots
  \begin{description}
  \item[\rm\ldots] the Sun set at your observing site?
  \item[\rm\ldots] the Earth set at your observing site?
  \end{description}
\item[(3.4)] Determine the Lunar (Selenographic) coordinates of this
  observing site (as defined in the provided lunar map). What is the name of
  the large surface lunar area where your observing site is situated? Do not
  use your national language, please use the official IAU nomenclature. (See
  the lunar map.)
\item[(3.5)] Estimate the distance from the observing site to the Apollo-11
  landing site ($0.6875^\circ$\,N, $23.4333^\circ$\,E).
\end{description}

\ioaafig{2019_OBS-PLAN_3-s1.pdf}

\probhead{15.12}{Simulated Sky}{2021}{Bogot\'a (remote)}{OBS-PLAN 1}{75}
% subject: Chapter 1 (formerly Problem 1.53)
% source id: 2021_OBS-PLAN_1   tag: simulated sky, 17 questions
17 questions in 45 minutes. On the projection screen in front of you, a
series of images will be projected. Look at each image and answer the
following questions.

\begin{description}
\item[(1.1)] (Image 1, 2 minutes) Identify the drawn lines on the image.
  [Answer table: for line numbers 1, 2, 3, tick one of --- Celestial
  Equator, Ecliptic, Local Meridian, Local Vertical, Zero RA, Galactic
  Equator.] \hfill[6.0 pt]
\item[(1.2)] (Image 2, 2 minutes) This image was taken from the town of Fada
  in Chad. What is the latitude of this place? Select only one choice:
  \begin{description}
  \item[A.] $10^\circ$
  \item[B.] $17^\circ$
  \item[C.] $22^\circ$
  \item[D.] $33^\circ$
  \end{description}
  \hfill[3.0 pt]
\item[(1.3)] (Image 3, 3 minutes) Two of the following stars are transiting
  the meridian in the image. Tick the boxes next to those stars:
  Arcturus, Mizar, Spica, Regulus. \hfill[4.0 pt]
\item[(1.4)] (No image; 8 minutes for questions 4--6) If the age of the Moon
  is 27.5 days, in which direction can we observe it and at what time?
  Select only one choice.
  \begin{description}
  \item[A.] Before sunrise to the east.
  \item[B.] After sunset to the west.
  \item[C.] Close to zenith at sunrise.
  \item[D.] At 9 am rising in the East.
  \end{description}
  \hfill[2.0 pt]
\item[(1.5)] Which of the Right Ascensions ($RA$) listed below will
  approximately coincide with the local meridian at Pune, India
  ($74^\circ$E) on April 10 at 19:15 local time?
  \begin{description}
  \item[A.] 6 h
  \item[B.] 7 h
  \item[C.] 8 h
  \item[D.] 9 h
  \item[E.] 10 h
  \end{description}
  \hfill[5.0 pt]
\item[(1.6)] At what time during this night (April 10) would we find Regulus
  [$\alpha$ Leo, $RA = 10^{\mathrm{h}}08^{\mathrm{m}}22.3^{\mathrm{s}}$,
  $Dec = +11^\circ 58' 02''$] at its maximum altitude as seen from Pune,
  India? Select only one choice.
  \begin{description}
  \item[A.] 18:00h
  \item[B.] 19:00h
  \item[C.] 20:00h
  \item[D.] 21:00h
  \item[E.] 22:00h
  \end{description}
  \hfill[4.0 pt]
\end{description}

For questions 7--10 the projected sky (Image 4) corresponds to latitude
$-30^\circ$ (South) and longitude $-74^\circ$ (West).

\begin{description}
\item[(1.7)] (Image 4, 2 minutes) Select the constellation pack that crosses
  the celestial equator on the image.
  \begin{description}
  \item[A.] Ori, Tau, Cet, Eri, Psc, Aqr, Aql, Ser
  \item[B.] Mic, Tau, Per, Cae, Mon, Aql, And, Ser
  \item[C.] Gru, Tau, Cha, Ara, Aqr, Aql, And, Aps
  \item[D.] For, Tri, Cet, Psc, Sct, Aql, And, Ser
  \end{description}
  \hfill[3.0 pt]
\item[(1.8)] (Image 4, 2 minutes) Choose the name of the stars indicated by
  the arrows (tick the correct cell).
  [Answer table: for star numbers 1, 2, 3, tick one of --- Achernar,
  Alnair, Altair, Alnath, Ankaa, Canopus.] \hfill[6.0 pt]
\item[(1.9)] (Image 4, 3 minutes) Are the following Messier objects present
  in the projected sky? Select YES or NO for each:
  M15 (Peg), M16 (Aql), M27 (Vul), M6 (Sco), M15 (Cas), M27 (CMi).
  \hfill[6.0 pt]
\item[(1.10)] (Image 4, 3 minutes) Choose the name of the selected stars
  (tick the correct cell). If the same object is listed more than once,
  choose only the option with the correct constellation.
  [Answer table: for star numbers 1, 2, 3, tick one of --- Miaplacidus
  (Car), Miaplacidus (Lac), Sham (Sge), Sham (For), The Persian (Gru), The
  Persian (Ind).] \hfill[3.0 pt]
\item[(1.11)] (Image 5, for questions 11 and 12, 4 minutes) From the list
  below select the option that contains the constellations that are along
  the line of the ecliptic:
  \begin{description}
  \item[A.] Psc, Cap, Sgr, Oph
  \item[B.] Lib, Sco, Oph, Sgr
  \item[C.] Cap, Sge, Lib, Vir
  \item[D.] Sgr, Oph, Sco, Vir
  \item[E.] CrA, Sco, Oph, Leo
  \end{description}
  \hfill[4.0 pt]
\item[(1.12)] Which constellation or part of constellation is encircled in
  the image?
  \begin{description}
  \item[A.] Sagittarius (Sgr)
  \item[B.] Corvus (Crv)
  \item[C.] Telescopium (Tel)
  \item[D.] Corona Australis (CrA)
  \end{description}
  \hfill[3.0 pt]
\item[(1.13)] (Image 6, 2 minutes) You are making an observation through a
  telescope pointed towards the centre of the galaxy. You find an
  interesting object as shown in image 6. What is this object?
  \begin{description}
  \item[A.] M42 -- Orion Nebula
  \item[B.] M31 -- Andromeda Galaxy
  \item[C.] NGC 2024 -- Flame Nebula
  \item[D.] M8 -- Lagoon Nebula
  \item[E.] M20 -- Trifid Nebula
  \end{description}
  \hfill[3.0 pt]
\item[(1.14)] (Image 7, 2 minutes) In the Messier Astronomical Catalogue
  there is a mysterious object: a galaxy. This object in the image cannot be
  easily identified. Some astronomers had said that it was a duplicate of
  another Messier object, but here we have it as an original and awesome
  galaxy. Which object is this?
  \begin{description}
  \item[A.] M31 -- Andromeda Galaxy
  \item[B.] M33 -- Triangulum Galaxy
  \item[C.] M51 -- Whirlpool Galaxy
  \item[D.] M101 -- Pinwheel Galaxy
  \item[E.] M102 -- Spindle Galaxy
  \end{description}
  \hfill[5.0 pt]
\item[(1.15)] (Image 8, 2 minutes) Estimate the magnitude of Shaula in the
  constellation Scorpius.
  \begin{description}
  \item[A.] 1
  \item[B.] 1.5
  \item[C.] 2
  \item[D.] 2.5
  \item[E.] 3
  \end{description}
  % note: the last option is misprinted "C." in the original paper
  \hfill[2.0 pt]
\item[(1.16)] (Image 9, 4 minutes) Several stars have been indicated in
  image 9. Tick the stars indicated in the table below.
  [Answer table: for star numbers 1, 2, 3, 4, tick one of --- Acubens,
  Adhara, Aludra, Alzirr, Arneb, Gomeisa, $\alpha$ Mon, Mebsuta, Mirzam,
  Wasat.] \hfill[8.0 pt]
\item[(1.17)] (Image 10, 4 minutes) A Messier object has been indicated in
  image 10. Tick the correct cell below to indicate what it is.
  [Answer table: for object numbers 1, 2, 3, 4, tick one of --- M41, M42,
  M46, M45, M48, M50, M64, M93.] \hfill[8.0 pt]
\end{description}

\probhead{15.13}{Planetarium round}{2023}{Katowice, Poland}{OBS-PLAN 1}{20}
% subject: Chapter 1 (formerly Problem 1.54)
% source id: 2023_OBS-PLAN_1   tag: planetarium sky identification
% Task 1 (`Knowledge of the sky') of the three-task planetarium round.
The projector will display the sky as seen from near the equator
($0^\circ$\,N, $19^\circ$\,E). The rotation of the sky will be stopped for
about 2 minutes for part (a), then it will start to rotate for parts (b) and
(c). The objects for parts (b) and (c) will be displayed simultaneously.
(Projection time 10 minutes.)

\begin{description}
\item[(a)] A meteor shower will be visible in the sky. Determine the
  constellation of the radiant and estimate its right ascension and
  declination coordinates. \hfill(3 points)
\item[(b)] Identify which of the following variable stars visible in the sky
  are in low (write `DIM') or high (write `BRIGHT') brightness states. The
  mean magnitude as shown in the atlas and the magnitude range are given for
  each star.

  \begin{center}
  \begin{tabular}{lccc}
  Name & atlas mag. & mag. range & DIM / BRIGHT \\
  \hline
  $\gamma$ Cas (Cih) & 2 & 1.6 -- 3.0 & \\
  $\delta$ Cep & 4 & 3.5 -- 4.4 & \\
  $\mu$ Cep (Erakis) & 4 & 3.4 -- 5.1 & \\
  $\beta$ Per (Algol) & 2 & 2.2 -- 3.4 & \\
  $o$ Cet (Mira) & 3.5 & 2.0 -- 10.1 & \\
  $\chi$ Cyg & 4.5 & 3.3 -- 14.1 & \\
  L$_2$ Pup & 4.5 & 2.6 -- 6 & \\
  $\delta$ Sco (Dschubba) & 2 & 1.6 -- 2.3 & \\
  \end{tabular}
  \end{center}
  \hfill(8 points)
\item[(c)] Identify the constellations whose borders are marked and give
  their IAU abbreviations. \hfill(9 points)
\end{description}

\probhead{15.14}{Retrograde Mars}{2023}{Katowice, Poland}{OBS-PLAN 2}{20}
% subject: Chapter 3 (formerly Problem 3.38)
% source id: 2023_OBS-PLAN_2   tag: Mars retrograde loop, heliocentric geometry
The projector will display Mars moving relative to the background stars over
one season of visibility (1.5 years) starting from the heliacal rising, chosen
so that Mars will be at maximum ecliptic latitude at opposition.

The ecliptic will also be displayed, marked with the positions of the Sun
during the year and the current date. The Sun will always be below the
horizon. Synodic period of Mars $= 780$ days. (Projection time: 10 minutes.)

\begin{description}
\item[(a)] Record the following quantities: \hfill(8 points)
  \begin{description}
  \item[i.] the dates of quadrature (when the elongation of Mars is $90^\circ$)
  \item[ii.] the date of the beginning of retrograde motion and the date of
    the end of retrograde motion
  \item[iii.] the date of opposition
  \item[iv.] the ecliptic latitude at opposition
  \item[v.] the width in ecliptic longitude of the loop made by the planet
  \end{description}
\end{description}

Based on your observations and assuming the orbits of Earth and Mars are
circular,
\begin{description}
\item[(b)] On the answer sheet, mark the positions of the Sun, Earth and Mars
  at the moments of opposition and quadrature in the heliocentric system and
  determine the radius of the orbit of Mars in a.u.\ geometrically, without
  using Kepler's Laws. Show your method in the answer sheet. \hfill(9 points)
\item[(c)] Derive the inclination of the orbit of Mars to the ecliptic.
  \hfill(3 points)
\end{description}

\probhead{15.15}{TRAPPIST-1}{2023}{Katowice, Poland}{OBS-PLAN 3}{35}
% subject: Chapter 9 (formerly Problem 9.12)
% source id: 2023_OBS-PLAN_3   tag: planetarium exoplanet transits, resonances
Aliens have found out that Earth's astronomers discovered planets in the
TRAPPIST-1 system by observing numerous transits. They have used their flying
saucer (similar to the one you were in for the observation round) to take you
to the 5th planet (designated~f) of TRAPPIST-1, and have asked you to show
them the methods Earthlings use to uncover the parameters of the system. A
clock displaying time in Earth hours will be visible. The whole presentation
lasts 520\,h (1\,s represents 1\,h). (Projection time 10 minutes)

Based on your observations (you can use the space on the last sheet for
observing notes),
\begin{description}
\item[(a)] determine the following quantities for the planet you are on (use
  Earth hours for the times): \hfill(7 points)
  \begin{description}
  \item[i.] length of the sidereal day [h]
  \item[ii.] orbital period [h]
  \item[iii.] length of the `solar' day [h]
  \item[iv.] circular orbit --- YES / NO
  \item[v.] obliquity (axial tilt)
  \end{description}
\item[(b)] and the following quantities for each planet b, c, d and e:
  synodic period [h] and maximum elongation [$^\circ$]. \hfill(16 points)
\item[(c)] calculate the orbital period in hours and the semi-major axis in
  tau (where 1 tau $=$ ``TRAPPIST-1f astronomical unit'' $=$ the semi-major
  axis of the orbit of TRAPPIST-1f) of each planet b, c, d, e.
  \hfill(8 points)
\item[(d)] The term `gravitational resonance' is used to describe the
  phenomenon when ratio % sic: source omits "the"
  of the orbital periods of two planets in a system
  is close to the ratio of two integers. Find which pair(s) of planets
  correspond to each of the listed resonances if any: 3:2, 8:5, 5:3, 8:3,
  4:1, 6:1. \hfill(4 points)
\end{description}

\probhead{15.16}{Planetarium problem 1}{2024}{Rio de Janeiro, Brazil}{OBS-PLAN 1}{}
% subject: Chapter 1 (formerly Problem 1.55)
% source id: 2024_OBS-PLAN_1   tag: planetarium identification
The Moon and planets are shown in the projected sky, if they are above the
horizon. The Milky Way and other diffuse objects are not shown. In the
planetarium's projection, azimuth starts at the North Cardinal Point, which
is assigned 0 degrees; the East Cardinal Point has an azimuth of 90 degrees,
and so on.

\begin{description}
\item[Question 1 (4 points)] Write the names of the zodiacal constellations
  (in either Latin or the standard IAU abbreviation) that are partially or
  completely above the horizon.
\item[Question 2 (4 points)] There is a bright star (brighter than 1.3 mag)
  located at each of the Alt-Az coordinates given below. Give their Bayer
  designation.
  \begin{description}
  \item[(a)] Star 1: Azimuth: $254^\circ 23'$; Altitude: $68^\circ 22'$
  \item[(b)] Star 2: Azimuth: $113^\circ 24'$; Altitude: $24^\circ 30'$
  \end{description}
\item[Question 3 (2 points)]~
  \begin{description}
  \item[(a)] The alpha star of a boreal constellation (i.e. northern
    hemisphere constellation) lies almost on the local meridian. Give the
    Bayer designation of this star.
  \item[(b)] Give the Messier Catalogue number, or New General Catalogue
    designation, of a nebula that is located almost at the intersection of
    the Ecliptic and the local meridian.
  \end{description}
\end{description}

\probhead{15.17}{Planetarium problem 2}{2024}{Rio de Janeiro, Brazil}{OBS-PLAN 2}{}
% subject: Chapter 1 (formerly Problem 1.56)
% source id: 2024_OBS-PLAN_2   tag: planetarium identification
The Moon and planets are shown in the projected sky, if they are above the
horizon. The Milky Way and other diffuse objects are not shown. In the
planetarium's projection, azimuth starts at the North Cardinal Point, which
is assigned 0 degrees; the East Cardinal Point has an azimuth of 90 degrees,
and so on.

\begin{description}
\item[Question 1 (4 points)]~
  \begin{description}
  \item[(a)] In which constellation is the Moon located? Give the Latin name
    or the IAU abbreviation.
  \item[(b)] Given that, for this projected sky, the Sun has set 1 hour and
    26 minutes ago, identify the current month.
  \end{description}
\item[Question 2 (4 points)] Which of the following stars are circumpolar?
  Circle your answers.

  \begin{center}
  \begin{tabular}{cccc}
  $\alpha$ Car & $\gamma$ Cru & $\beta$ Hyi & $\beta$ Pav \\
  $\beta$ Crv & $\beta$ Car & $\alpha$ Eri & $\upsilon$ Oct \\
  $\alpha$ Sco & $\beta$ Lup & $\gamma$ Vel & $\alpha$ Cnc \\
  \end{tabular}
  \end{center}
\item[Question 3 (2 points)] A certain artificial satellite was seen passing
  through this sky with the following coordinates:
  \begin{itemize}
  \item At rise: Azimuth: $98^\circ 47'$; Height: $7^\circ 31'$
  \item At maximum brightness: Azimuth: $2^\circ 27'$; Height:
    $54^\circ 47'$
  \item At disappearance: Azimuth: $311^\circ 19'$; Height: $26^\circ 47'$
  \end{itemize}
  Using the azimuth references given above, write the Bayer designation of
  one bright star, for each case, very close to the position of the
  satellite, when it is at:
  \begin{description}
  \item[(a)] Maximum brightness
  \item[(b)] Disappearance
  \end{description}
\end{description}

\probhead{15.18}{Hypothesized Planet in the Solar System}{2025}{Mumbai, India}{OBS-PLAN OP01}{25}
% subject: Chapter 3 (formerly Problem 3.36)
% source id: 2025_OBS-PLAN_OP01   tag: orbit inclinations from planetarium traces
For this question use sky map Map-OP01. This is an angle preserving
projection of the sky.

The sky projected on the planetarium dome is as seen from the solar system
barycentre. The Earth's orbit as seen from this position has been marked in
green both on the projected sky and in sky map Map-OP01. The first 3 minutes
are for dark adaptation and familiarization of the sky.

After first 3 minutes, four Trans-Neptunian Objects (TNOs) labelled O1, O2,
O3 and O4 will begin to appear on the dome. The brightness of the TNOs has
been enhanced to make them visible. No other solar system object is visible
in the projected sky. Time has been sped up to make their movements
noticeable.

We will assume that the orbital plane of a hypothesized new planet has an
inclination equal to the mean of the inclinations of the orbits of these 4
TNOs.

\ioaafig{2025_OBS-PLAN_OP01-s3.pdf}

\begin{description}
\item[(OP01.1)] Arrange the 4 TNOs in ascending order of the lengths of their
  semi-major axes and write their names in the appropriate boxes in the
  Summary Answersheet. Show measurements in the working sheet to justify your
  answer. \hfill[8]
\item[(OP01.2)] Trace at least 50\% of the visible trajectories for each of
  the 4 TNOs on the map Map-OP01. From these traces find the inclination
  angles, $i_1, i_2, i_3, i_4$, of the 4 TNOs, respectively, with respect to
  the Ecliptic. Hence obtain the best estimate of the inclination,
  $i_{\mathrm{hp}}$, of the plane of the orbit of the hypothesized planet.
  \hfill[17]
\end{description}

\probhead{15.19}{Observer on a Ringed Planet}{2025}{Mumbai, India}{OBS-PLAN OP02}{25}
% subject: Chapter 3 (formerly Problem 3.37)
% source id: 2025_OBS-PLAN_OP02   tag: sky as seen from a ringed planet
For this question use sky map Map-OP02.

The sky projected on the planetarium dome is as seen by an observer on some
hypothetical ringed planet. The observer is located very near to the equator
of the planet.

Assume that the ring lies in the equatorial plane of the planet and the
distance of the planet from its parent star is more than 15\,au. Note that
except for the 10 moons of this planet no other object from this planetary
system (other planets, their satellites, asteroids and even the parent star)
are visible on the projected sky.

At the start of the question the sky will make 3 full rotations at the rate
of 1 rotation every minute. After that the sky will stop rotating at the
observer's midnight. This sky view will remain on the dome for next 7
minutes.

Solve the following questions once the sky stops rotating.

\ioaafig{2025_OBS-PLAN_OP02-f1.pdf}
% exam; not part of the question paper PDF.

\begin{description}
\item[(OP02.1)] Mark the positions of any 4 (only 4) out of the 10 visible
  moons of the planet with a $\otimes$ sign and label them as ``M'' in the
  sky map. \hfill[8]
\item[(OP02.2)] The figure shown below is as seen from the top of the north
  pole of the planet (grey circle) and the horizontal lines are parallel
  light rays from its parent star.

% FIGURE NEEDED: diagram from page 2 of the 2025 planetarium paper — grey
% circle of radius R_P = 6.00 x 10^4 km labelled P, with parallel light rays
% from the parent star drawn as horizontal lines.

  Using the view of the ring on the dome, draw the inner and outer boundaries
  of the ring on this diagram, reproduced in the Summary Answersheet. Your
  sketch must be to scale. \hfill[12]
\item[(OP02.3)] Use the diagram above to determine the width of the ring,
  $w_{\mathrm{ring}}$ (in km), given that the radius of the planet is
  $6.00 \times 10^4$\,km. \hfill[5]
\end{description}

\probhead{15.20}{The Sky over Hue}{2026}{Hanoi, Vietnam}{OBS-PLAN OP01}{45}
% subject: Chapter 1
% source id: 2026_OBS-PLAN_OP01   tag: constellations, planets, missing stars, apparent solar time
This is the simulated sky in February 10, 2025 as seen from Quan Tu\d{o}ng
\DJ\`ai, an ancient observatory located in Hu\'{\^e}, Vietnam. The Sun is not
shown in the simulation.
% editorial: the horn diacritics of the second word of the observatory's name
% (u-horn, o-horn with dot below) cannot be set in T1 encoding; approximated.

\begin{description}
\item[(a)] (5 points) In which constellation is the Sun located? Write its
  Latin name or IAU abbreviation.
\item[(b)] (5 points) Which constellation is at the West cardinal point? Write
  its Latin name or IAU abbreviation.
\item[(c)] (14 points) Check the boxes next to the constellations hosting the
  visible planets (points will be deducted for incorrect or excess answers).
  Underline the name of the checked constellation (in the answer sheet) where
  Mercury is located.

  \begin{center}
  \begin{tabular}{l@{\ }c@{\qquad}l@{\ }c@{\qquad}l@{\ }c@{\qquad}l@{\ }c}
  Aries/Ari & $\square$ & Taurus/Tau & $\square$ & Gemini/Gem & $\square$ & Cancer/Cnc & $\square$ \\
  Leo & $\square$ & Virgo/Vir & $\square$ & Libra/Lib & $\square$ & Ophiuchus/Oph & $\square$ \\
  Scorpius/Sco & $\square$ & Sagittarius/Sgr & $\square$ & Capricornus/Cap & $\square$ & Aquarius/Aqr & $\square$ \\
  Pisces/Psc & $\square$ & Orion/Ori & $\square$ & Auriga/Aur & $\square$ & Cetus/Cet & $\square$ \\
  \end{tabular}
  \end{center}

\item[(d)] (16 points) Check the boxes next to 7 stars (in the answer sheet)
  that are missing from this simulated sky (points deducted for incorrect or
  excess answers).

  \begin{center}
  \begin{tabular}{l@{\ }c@{\qquad}l@{\ }c@{\qquad}l@{\ }c}
  Hamal/$\alpha$ Ari & $\square$ & Menkalinan/$\beta$ Aur & $\square$ & Mirach/$\beta$ And & $\square$ \\
  Alhena/$\gamma$ Gem & $\square$ & Diphda/$\beta$ Cet & $\square$ & Schedar/$\alpha$ Cas & $\square$ \\
  Menkar/$\alpha$ Cet & $\square$ & Sadalmelik/$\alpha$ Aqr & $\square$ & Alderamin/$\alpha$ Cep & $\square$ \\
  Scheat/$\beta$ Peg & $\square$ & Algol/$\beta$ Per & $\square$ & Sadr/$\gamma$ Cyg & $\square$ \\
  Ankaa/$\alpha$ Phe & $\square$ & Rigel/ $\beta$ Ori & $\square$ & Fomalhaut/$\alpha$ PsA & $\square$ \\
  \end{tabular}
  \end{center}

\item[(e)] (5 points) Estimate the local apparent solar time.
\end{description}

\probhead{15.21}{Comet Hale-Bopp}{2026}{Hanoi, Vietnam}{OBS-PLAN OP02}{15}
% subject: Chapter 1
% source id: 2026_OBS-PLAN_OP02   tag: comet path dates, ecliptic crossing, elongation
Comet Hale-Bopp (C/1995 O1) appears in this simulated sky from February 15,
1997 to June 15, 1997.

\begin{description}
\item[(a)] (5 points) Determine the date (DD/MM) on which the comet nucleus is
  closest to Betelgeuse ($\alpha$ Ori).
\item[(b)] (5 points) Determine the date (DD/MM) on which the comet nucleus
  crosses the ecliptic.
\item[(c)] (5 points) Estimate the elongation $\varepsilon$ of the comet
  nucleus from the Sun when it crosses the ecliptic.
\end{description}

\probhead{15.22}{The Sky from the Far Side of the Moon}{2026}{Hanoi, Vietnam}{OBS-PLAN OP03}{15}
% subject: Chapter 1
% source id: 2026_OBS-PLAN_OP03   tag: Sun and Earth constellations, Earth altitude from a lunar site
This is the simulated sky above a site ($45.4^\circ$S, $180^\circ$E) on the
far side of the Moon, not so far from Von K\'arm\'an Crater, where Chang'e-4
became the first spacecraft to soft-land there in 2019. Now the Moon is in its
first quarter at zero libration as viewed from Earth.

\begin{description}
\item[(a)] (5 points) In which constellation is the Sun located? Write its
  Latin name or IAU abbreviation.
\item[(b)] (5 points) In which constellation is the Earth located? Write its
  Latin name or IAU abbreviation.
\item[(c)] (5 points) What is the altitude $h$ of the Earth from this site?
\end{description}

\clearpage

\section*{Observation --- written and sky map}

\probhead{15.23}{Pointer: equator, equinox, Pegasus, Aries, Aldebaran offset}{2007}{Chiang Mai, Thailand}{OBS-SKY 1}{10}
% subject: Chapter 1 (formerly Problem 1.57)
% source id: 2007_OBS-SKY_1   tag: celestial-object pointer tasks
Use only the provided celestial object pointer to point at the target
specified in the questions. The marker (examiner) at the observation station
will mark the answers directly in the question sheet.

\begin{description}
\item[(1.1)] Move the pointer along the celestial equator. \hfill(1 point)
\item[(1.2)] Aim the pointer at the vernal equinox. \hfill(1 point)
\item[(1.3)] In the constellation of Pegasus and its vicinity there is an
  obvious square of bright stars (Great Square of Pegasus); aim the pointer
  at the brightest star of the square. \hfill(2 points)
\item[(1.4)] Aim the pointer at the star named alpha-Arietis
  ($\alpha$-Ari). \hfill(2 points)
\item[(1.5)] Start from the star named Aldebaran ($\alpha$-Tauri) in the
  constellation Taurus, turn the pointer 35 degrees northward followed by 6
  degrees westward (in equatorial coordinates). Then, aim the pointer at the
  brightest star in the field of view. \hfill(4 points)
\end{description}

\probhead{15.24}{Binoculars: Hyades limiting magnitude, M31 shape}{2007}{Chiang Mai, Thailand}{OBS-SKY 2}{10}
% subject: Chapter 5 (formerly Problem 5.31)
% source id: 2007_OBS-SKY_2   tag: binocular observation, limiting magnitude
Use the provided binoculars to observe the objects and then answer the
questions by writing or drawing directly in the question sheets.

\begin{description}
\item[(2.1)] The open star cluster ``Hyades'' in constellation Taurus is one
  of the nearest clusters to us, being only 151 light years away. From the
  provided chart with brightness of some stars indicated by the apparent
  magnitude in parentheses, please estimate the apparent magnitude of the
  star Gamma-Tauri ($\gamma$-Tauri) to the nearest first decimal digit.
  \hfill(5 points)

\ioaafig{2007_OBS-SKY_2-s1.pdf}

\item[(2.2)] Observe the Andromeda Galaxy (M31) then draw the approximate
  shape and size of the galaxy that you see through the binoculars in the
  frame below with correct orientation (in equatorial coordinates). The
  field of view of the binoculars is 6.8 degrees. \hfill(5 points)

\ioaafig{2007_OBS-SKY_2-s2.pdf}
\end{description}

\probhead{15.25}{Missing bright stars in Perseus and Andromeda}{2009}{Tehran, Iran}{OBS-SKY 1}{40}
% subject: Chapter 1 (formerly Problem 1.58)
% source id: 2009_OBS-SKY_1   tag: naked-eye star identification on chart
Figure 1 (frame size $100^\circ \times 70^\circ$) shows a part of the sky,
for 22 October 2009 at 21:00 local time. Four bright stars in the Perseus
and Andromeda constellations are missing in this chart. Find these missing
stars by looking at the sky. Then, draw a cross on the location of each
missing bright star in these two constellations on the chart (i.e.\
Figure 1). Use the numbers in Table 1-1 to indicate these crosses.

Note: Polaris is indicated by the ``N'' symbol in Figure 1.

\begin{center}
\begin{tabular}{ccl}
\multicolumn{3}{c}{Table 1-1}\\
Number & Common Name & Bayer Name \\
\hline
1 & Mirfak & Alpha Persei \\
2 & Alpheratz & Alpha Andromedae \\
3 & -- & Epsilon Persei \\
4 & Menkib & Xi Persei \\
5 & -- & Gamma Persei \\
6 & Algol & Beta Persei \\
7 & Almach & Gamma Andromedae \\
8 & -- & Delta Andromedae \\
9 & -- & 51 Andromedae \\
10 & Mirach & Beta Andromedae \\
11 & Atik & Zeta Persei \\
\end{tabular}
\end{center}

\ioaafig{2009_OBS-SKY_1-s2.pdf}

\probhead{15.26}{Naked-eye magnitude estimates near Cepheus}{2009}{Tehran, Iran}{OBS-SKY 2}{80}
% subject: Chapter 6 (formerly Problem 6.31)
% source id: 2009_OBS-SKY_2   tag: naked-eye magnitude comparison
\begin{description}
\item[2.1:] Figure 2 shows a part of the sky which contains \textbf{Cephei
  constellation}, for 22 October 2009 at 22:00 local time. Five bright stars
  in Cephei constellation are identified by numbers (1, 2, \dots, 5) and
  common names. Estimate the angular distances (in units of degrees) between
  two pairs of stars shown in table 2-1 and complete this table with your
  answers. \textbf{(40 Points)}

  \medskip
  \noindent Table 2-1: Angular Distance

  \begin{center}
  \begin{tabular}{lc}
  \hline
  Pairs of stars & Angular Distance (degrees) \\
  \hline
  1 (Errai) and 2 (Alfirk) & \\
  1 (Errai) and 3 (Alderamin) & \\
  \hline
  \end{tabular}
  \end{center}

\item[2.2:] Use table 2-2 and figure 2, then Estimate the ``apparent visual
  magnitude'' of stars 2 (Alfirak) and 3 (Alderamin) and complete Table 2-3.
  % sic: "Alfirak" as printed in the source; the star is called Alfirk
  % elsewhere in the same problem
  \textbf{(40 Points)}

  \medskip
  \noindent Table 2-2

  \begin{center}
  \begin{tabular}{cc}
  \hline
  Star Name & Apparent Visual Magnitude \\
  \hline
  Polaris & 1.95 \\
  Altais & 3.05 \\
  Segin & 3.34 \\
  \hline
  \multicolumn{2}{c}{All of these stars, are marked in the figure 2} \\
  \hline
  \end{tabular}
  \end{center}

  \medskip
  \noindent Table 2-3: Magnitude Estimation

  \begin{center}
  \begin{tabular}{ccc}
  \hline
  Star Number & Star Name & Apparent Visual Magnitude \\
  \hline
  2 & Alfirk & \\
  3 & Alderamin & \\
  \hline
  \end{tabular}
  \end{center}
\end{description}

\ioaafig{2009_OBS-SKY_2-f1.pdf}
% ($IOAA_PAPERS_DIR/observation/2009_OBS.pdf, page 7): sky chart around
% Cepheus with stars 1 (Errai), 2 (Alfirk), 3 (Alderamin), 5 (Iota Cephei),
% plus Polaris, Altais and Segin marked.

\probhead{15.27}{Laser pointing: Antares, Vega, Altair, Peacock}{2012}{Rio de Janeiro, Brazil}{OBS-SKY 5}{}
% subject: Chapter 1 (formerly Problem 1.59)
% source id: 2012_OBS-SKY_5   tag: naked-eye star and constellation pointing
Use your green laser pointer to spot the stars Antares, Vega, Altair and
Peacock. Also point to the constellation Corona Australis.

\probhead{15.28}{Star chart part II}{2014}{Suceava, Romania}{OBS-SKY 2}{}
% subject: Chapter 1 (formerly Problem 1.60)
% source id: 2014_OBS-SKY_2   tag: star chart marking
In this part you will use the sky-map found in the envelope. The map
represents the sky in Suceava (latitude $47^\circ 39'$ North, longitude
$26^\circ 15'$ East) at 19:00 UT on the day of the test. The observer who
made the sky-map was at a very high altitude above Suceava; the Zenith point
is in the centre of the chart. Use a pencil for marking and drawing lines on
the sky-map. Solve questions 1 to 4 on one copy of the map and questions 5
to 8 on the second copy of the map. You have 30 minutes to finish this part.

\begin{description}
\item[(1)] (2p) Draw on the map the horizon for an observer located on the
  ground in Suceava.
\item[(2)] (8p) Draw the celestial equator, the ecliptic, the galactic
  equator and the local meridian on the map with continuous lines.
\item[(3)] (9p) Mark the cardinal points (as N for north, E for east, S for
  south and W for west). Mark all the visible planets (except Uranus and
  Neptune) of the Solar System on the map and number them as 1, 2, \ldots,
  6 in the order of increasing orbital radius (skip number 3 for the Earth).
  Note that planets are not currently shown on the map.
\item[(4)] (4p) Identify and mark the four brightest stars in the visual
  band above the horizon line. Number the stars starting from 1 --- the
  brightest, and continue with the fainter ones till number 4 for the
  faintest. Fill in a table with the Bayer names of the four identified
  stars.
\item[(5)] (6p) Draw on the map approximate figures of any 15 constellations
  which lie completely above the horizon. Each constellation you mark should
  be identified on the map with the IAU abbreviation, using the provided
  table of constellations.
\item[(6)] (5p) Mark on the map the positions of the following objects:
  \begin{description}
  \item[(a)] the Messier objects M31, M27, M13;
  \item[(b)] $\beta$ Cygni, $\delta$ Ursae Minoris.
  \end{description}
\item[(7)] (10p) Estimate the sidereal time of the map; write the value in
  the box.
\item[(8)] (6p) Estimate the equatorial coordinates (right ascension
  $\alpha$ and declination $\delta$) of the star Altair ($\alpha$ Aquilae).
\end{description}

\textit{[The envelope sky map was distributed separately and is not in the available source papers]}

\probhead{15.29}{Nova Sagittarii field sketch}{2015}{Magelang, Indonesia}{OBS-SKY 1}{}
% subject: Chapter 1 (formerly Problem 1.61)
% source id: 2015_OBS-SKY_1   tag: naked-eye chart drawing
In the blank square given in the answer sheet (field of view $45^\circ$,
North direction marked), you should assume that Nova Sagittarii 2 is in the
centre (marked by a green star):
\begin{center}
RA (J2000): $18^{\mathrm{h}} 36^{\mathrm{m}} 56.84^{\mathrm{s}}$ \qquad
Dec (J2000): $-28^\circ 55' 39.8''$
\end{center}

\begin{description}
\item[(1)] Draw the declination and right ascension lines passing through
  Nova Sagittarii. Also draw parallel gridlines at a convenient interval and
  label them accordingly. (The field of view is $45^\circ$.)
  \hfill(10 points)
\item[(2)] Mark with a filled circle the five brightest stars in the field
  of view on the answer sheet, along with their names. You can add several
  other stars for your orientation. \hfill(25 points)
\item[(3)] Draw an unfilled circle representing the Moon's position tomorrow
  night (on 29 July 2015). \hfill(10 points)
\item[(4)] Mark the position of any five Messier objects in the field of
  view by a cross ($\times$), and write down their names beside their marks.
  \hfill(25 points)
\item[(5)] Draw a line representing the galactic plane. \hfill(10 points)
\item[(6)] Mark the positions of Nunki ($\sigma$ Sgr) and Shaula
  ($\lambda$ Sco) by a star ($*$). Draw a line connecting the two stars.
  Estimate the angular distance between these stars and write it on the line
  connecting both stars. \hfill(20 points)
\end{description}

\ioaafig{2015_OBS-SKY_1-s4.pdf}

\probhead{15.30}{Mark five stars on the map}{2016}{Bhubaneswar, India}{OBS-SKY OM1}{20}
% subject: Chapter 1 (formerly Problem 1.62)
% source id: 2016_OBS-SKY_OM1   tag: star chart marking
Mark any 5 (five) of the following stars on the map by putting a circle (O)
around the appropriate star and writing its code next to it. If you mark
more than 5 stars, only the first 5 in serial order will be considered.

\begin{center}
\begin{tabular}{cll}
Code & Name & Bayer Name \\
\hline
S1 & Caph & $\beta$ Cas \\
S2 & Asellus Australis & $\delta$ Cnc \\
S3 & Acrux & $\alpha$ Cru \\
S4 & Alphard & $\alpha$ Hya \\
S5 & Sheliak & $\beta$ Lyr \\
S6 & Albireo & $\beta$ Cyg \\
S7 & Rasalhague & $\alpha$ Oph \\
S8 & Kaus Australis & $\epsilon$ Sgr \\
\end{tabular}
\end{center}

\textit{[The full-sky answer map was distributed separately and is not in the available source papers]}

\probhead{15.31}{Sky map identification}{2016}{Bhubaneswar, India}{OBS-SKY OM2}{15}
% subject: Chapter 1 (formerly Problem 1.63)
% source id: 2016_OBS-SKY_OM2   tag: star chart identification
Mark the location of any 3 (three) of the following galaxies on the map by
putting a `+' sign at the appropriate place in the map and writing its code
next to it. If you mark more than 3 galaxies, only the first 3 in serial order
will be considered.

\begin{center}
\begin{tabular}{ccc}
Code & Name & M number \\
\hline
G1 & Triangulum Galaxy & M\,33 \\
G2 & Whirlpool Galaxy & M\,51 \\
G3 & Southern Pinwheel Galaxy & M\,83 \\
G4 & Virgo A & M\,87 \\
G5 & Sombrero Galaxy & M\,104 \\
\end{tabular}
\end{center}

\textit{[The full-sky answer map was distributed separately and is not in the available source papers]}

\probhead{15.32}{Sky map identification}{2016}{Bhubaneswar, India}{OBS-SKY OM3}{5}
% subject: Chapter 1 (formerly Problem 1.64)
% source id: 2016_OBS-SKY_OM3   tag: star chart identification
Draw the ecliptic on the map and label it as `E'.

\textit{[The full-sky answer map was distributed separately and is not in the available source papers]}

\probhead{15.33}{Sky map identification}{2016}{Bhubaneswar, India}{OBS-SKY OM4}{5}
% subject: Chapter 1 (formerly Problem 1.65)
% source id: 2016_OBS-SKY_OM4   tag: star chart identification
Show the position of the Autumnal Equinox (descending node of the ecliptic) on
the map by a `+' sign and label it as `A'.

\textit{[The full-sky answer map was distributed separately and is not in the available source papers]}

\probhead{15.34}{Sky map identification}{2016}{Bhubaneswar, India}{OBS-SKY OM5}{5}
% subject: Chapter 1 (formerly Problem 1.66)
% source id: 2016_OBS-SKY_OM5   tag: star chart identification
Draw the local meridian for Bhubaneswar on Winter Solstice day (22nd December)
at local midnight and label it as `M'.

\textit{[The full-sky answer map was distributed separately and is not in the available source papers]}

\probhead{15.35}{Day observation stations (7 problems)}{2017}{Phuket, Thailand}{OBS-SKY 1}{}
% subject: Chapter 1 (formerly Problem 1.67)
% source id: 2017_OBS-SKY_1   tag: daytime observation stations
This part of the exam consists of 7 problems located in separate exam
stations. You have 5 minutes to complete each problem. (The number of problems
was to be reduced to 5 if the night observation was successful.)

\begin{description}
\item[(P1)] \textbf{Naked eye observation from real sky with panoramic
  360-degree image.} Estimate the LST (Local Sidereal Time) at the time the
  image was taken, rounded to the nearest hour.

  Included: a panoramic 360-degree image of the sky at night at an unknown
  location; computer screen; keypad to pan around the image; coordinates of
  bright stars.
  \textit{[Interactive panoramic image was shown on an on-site computer station and is not reproduced here]}

  \begin{center}
  \begin{tabular}{llll}
  Name & Bayer Designation & Declination (Dec) & Right Ascension (RA) \\
  \hline
  Rigil Kentaurus & $\alpha$ Cen & $-60^\circ 50' 02.3737''$
    & $14^{\mathrm{h}}39^{\mathrm{m}}36.5^{\mathrm{s}}$ \\
  Arcturus & $\alpha$ Boo & $+19^\circ 10' 56''$
    & $14^{\mathrm{h}}15^{\mathrm{m}}39.7^{\mathrm{s}}$ \\
  Vega & $\alpha$ Lyr & $+38^\circ 47' 01''$
    & $18^{\mathrm{h}}36^{\mathrm{m}}56.3^{\mathrm{s}}$ \\
  Capella & $\alpha$ Aur & $+45^\circ 59' 53''$
    & $05^{\mathrm{h}}16^{\mathrm{m}}41.4^{\mathrm{s}}$ \\
  Altair & $\alpha$ Aql & $+08^\circ 52' 06''$
    & $19^{\mathrm{h}}50^{\mathrm{m}}47.0^{\mathrm{s}}$ \\
  Aldebaran & $\alpha$ Tau & $+16^\circ 30' 33''$
    & $04^{\mathrm{h}}35^{\mathrm{m}}55.2^{\mathrm{s}}$ \\
  Antares & $\alpha$ Sco & $-26^\circ 25' 55''$
    & $16^{\mathrm{h}}29^{\mathrm{m}}24.5^{\mathrm{s}}$ \\
  Spica & $\alpha$ Vir & $-11^\circ 09' 41''$
    & $13^{\mathrm{h}}25^{\mathrm{m}}11.6^{\mathrm{s}}$ \\
  Deneb & $\alpha$ Cyg & $+45^\circ 16' 49''$
    & $20^{\mathrm{h}}41^{\mathrm{m}}25.9^{\mathrm{s}}$ \\
  Dubhe & $\alpha$ UMa & $+61^\circ 45' 04''$
    & $11^{\mathrm{h}}03^{\mathrm{m}}43.7^{\mathrm{s}}$ \\
  Polaris & $\alpha$ UMi & $+89^\circ 15' 51''$
    & $02^{\mathrm{h}}31^{\mathrm{m}}49.1^{\mathrm{s}}$ \\
  Alpheratz & $\alpha$ And & $+29^\circ 05' 26''$
    & $00^{\mathrm{h}}08^{\mathrm{m}}23.3^{\mathrm{s}}$ \\
  Schedar & $\alpha$ Cas & $+56^\circ 32' 14''$
    & $00^{\mathrm{h}}40^{\mathrm{m}}30.4^{\mathrm{s}}$ \\
  \end{tabular}
  \end{center}

\item[(P2)] \textbf{Planet observation with real sky in panoramic 360-degree
  image.} Count the number of planets visible in this image above the horizon
  and name the constellations they are in (with IAU designations, i.e.\ Ursa
  Major or UMa).

  Included: a panoramic 360-degree image of the sky at night at an unknown
  location; computer screen; keypad to pan around the image.
  \textit{[Interactive panoramic image was shown on an on-site computer station and is not reproduced here]}

\item[(P3)] \textbf{Analemma on another planet.} Find the obliquity (axial
  tilt) of the planet.

  Included: a generated analemma (position of a star taken from the surface of
  a planet with interval separated by the mean solar day of the planet over an
  orbital period around a star) of a fictitious planet orbiting around a star.
  The result is graphed on a paper with each major grid representing
  $5^\circ$.
  \ioaafig{2017_OBS-SKY_1-s1.pdf}

\item[(P4)] \textbf{Exposure time from a photograph.} Estimate the exposure
  time of a given ``star trails'' image.

  Included: a ``star trails'' image that was taken by a still camera capturing
  an image over a period of time; ruler.
  \textit{[The star-trails photograph was provided at the exam station and is not in the available source papers]}

\item[(P5)] \textbf{Find True North from Moon shadow.} Draw an arrow pointing
  North in the data sheet.

  Included: simulated positions of the Moon shadows of a pole at certain
  intervals in the span of a day. The observer is located in the Southern
  hemisphere at latitude $27^\circ$S. The Moon's declination that night is
  $+15^\circ$. Ruler, compass (drawing tool), geometry kit.
  \ioaafig{2017_OBS-SKY_1-s2.pdf}
  \ioaafig{2017_OBS-SKY_1-s3.pdf}

\item[(P6)] \textbf{Find latitude from equatorial mount.} Without altering the
  polar alignment, find the observer's latitude based on a previously
  polar-aligned equatorial mount.

  Included: an equatorial-mount telescope that has already been properly
  aligned to a location in the northern celestial pole; bubble level. The
  latitude dial on the mount is covered (you may not use it).

\item[(P7)] \textbf{Precision polar alignment with equatorial mount.} Perform
  a polar alignment on the equatorial mount provided.

  Included: equatorial mount with polar scope (has not been polar aligned);
  date and time (GMT, UTC+0) of the time of performing the polar alignment;
  diagram of the sky's position at the time; a light source substituted for
  Polaris to be used for proper polar alignment (already visible in the polar
  scope); longitude of observer.

  Date and Time: 30 Aug 2017 / 23:30. \quad Longitude: $10^\circ$\,E.
  \textit{[The polar-scope sky diagram was provided at the exam station and is not in the available source papers]}
\end{description}

\probhead{15.36}{Day observation problem 1}{2018}{Beijing, China}{OBS-SKY O1}{40}
% subject: Chapter 1 (formerly Problem 1.68)
% source id: 2018_OBS-SKY_O1   tag: daytime observation
Figure 1 is a whole-sky star chart of Yanqing, Beijing at 20:30 tonight
(UTC+8) with the limiting magnitude $= 5^{\mathrm{m}}$ ($\mathrm{m}$ =
magnitude). Four stars (about $1^{\mathrm{m}}$--$3^{\mathrm{m}}$) and one
planet (brighter than $2^{\mathrm{m}}$) are missing in this chart. In the
chart, the distance from the centre is in proportion to zenith distance.

\ioaafig{2018_OBS-SKY_O1-s2.pdf}

\begin{description}
\item[(1)] (20 points) Draw a cross (X) on the location of each missing star
  and mark ``T'' on the chart, and draw a cross (X) on the location of the
  missing planet and mark ``P'' on the chart.
\item[(2)] (5 points) Please mark the orientation of the star chart with ``N''
  ``E'' ``S'' ``W'' at the edge of the star chart.
\item[(3)] (10 points) On the chart, the celestial equator passes through many
  constellations. Please write down the name of any five of these
  constellations (IAU codes).
\item[(4)] (5 points) Using the star chart, estimate the altitude of Aldebaran
  ($\alpha$ Tau), to the nearest degree.
\end{description}

\probhead{15.37}{Day observation problem 2}{2018}{Beijing, China}{OBS-SKY O2}{20}
% subject: Chapter 1 (formerly Problem 1.69)
% source id: 2018_OBS-SKY_O2   tag: daytime observation
Figure 2 is a star chart of a recent opposition of Jupiter. The grid in the
figure is the ecliptic coordinates. Please estimate the date of this
opposition, to the nearest day.

\ioaafig{2018_OBS-SKY_O2-s1.pdf}

\probhead{15.38}{Day observation problem 3}{2018}{Beijing, China}{OBS-SKY O3}{20}
% subject: Chapter 1 (formerly Problem 1.70)
% source id: 2018_OBS-SKY_O3   tag: daytime observation
Figure 3 is a star chart of a part of the sky on March 21, 2018. The longitude
and latitude of the observation site is $120^\circ$E, $40^\circ$N (UTC+8). The
grid in the figure is an equatorial grid. The thicker vertical line in the
centre is the meridian. Estimate the mean solar time to an accuracy of better
than $0.5^{\mathrm{h}}$.

\ioaafig{2018_OBS-SKY_O3-s1.pdf}

\probhead{15.39}{Day observation problem 4}{2018}{Beijing, China}{OBS-SKY O4}{20}
% subject: Chapter 1 (formerly Problem 1.71)
% source id: 2018_OBS-SKY_O4   tag: daytime observation
Figures 4.1--4.4 are four photos of Messier objects. For each of them, please
write down the Messier catalogue number and name the constellation where it is
located.

\ioaafig{2018_OBS-SKY_O4-s1.pdf}
\ioaafig{2018_OBS-SKY_O4-s2.pdf}
\ioaafig{2018_OBS-SKY_O4-s3.pdf}
\ioaafig{2018_OBS-SKY_O4-s4.pdf}

\probhead{15.40}{Solar Physics: the solar slap}{2021}{Bogot\'a (remote)}{OBS-SKY 1}{75}
% subject: Chapter 4 (formerly Problem 4.10)
% source id: 2021_OBS-SKY_1   tag: CME evolution from solar images
Using the provided series of solar images, select the appropriate images showing
the evolution of a coronal mass ejection (CME) and identify the CME front in
each. From these measurements, characterise the evolution of the CME.

\probhead{15.41}{Daytime observation O5}{2022}{Kutaisi, Georgia}{OBS-SKY O5}{9}
% subject: Chapter 1 (formerly Problem 1.72)
% source id: 2022_OBS-SKY_O5   tag: constellation identification
The image below shows schematic representations of three different types of
reflectors.

\ioaafig{2022_OBS-SKY_O5-s1.pdf}

For each one, choose the correct name from the list below. (3 points each)

(A) Newtonian \quad (B) Pfund \quad (C) Gregorian \quad (D) Herschelian \quad
(E) Keplerian \quad (F) Coud\'e \quad (G) Galilean \quad (H) Cassegrain

\probhead{15.42}{Daytime observation O6}{2022}{Kutaisi, Georgia}{OBS-SKY O6}{21}
% subject: Chapter 1 (formerly Problem 1.73)
% source id: 2022_OBS-SKY_O6   tag: daytime observation
The equatorial constellations are listed below. Fill in the missing ones in
the correct order. Use the IAU codes or official names accepted by the IAU.

\begin{enumerate}
\item Serpens (Ser)
\item Ophiuchus (Oph)
\item \underline{\hspace{6em}}
\item Aquarius (Aqr)
\item Pisces (Psc)
\item Cetus (Cet)
\item \underline{\hspace{6em}}
\item \underline{\hspace{6em}}
\item \underline{\hspace{6em}}
\item Monoceros (Mon)
\item Canis Minor (CMi)
\item \underline{\hspace{6em}}
\item \underline{\hspace{6em}}
\item \underline{\hspace{6em}}
\item Virgo (Vir)
\end{enumerate}

Which of these 15 constellations are also on the ecliptic? If you give more
than the required number of constellations, only the first ones you write (in
order) will be considered.

\probhead{15.43}{Daytime observation O7}{2022}{Kutaisi, Georgia}{OBS-SKY O7}{20}
% subject: Chapter 1 (formerly Problem 1.74)
% source id: 2022_OBS-SKY_O7   tag: daytime observation
Write the numbers in the Messier catalogue, and the constellations (IAU code)
of the given Messier objects (4 points for each object).

\ioaafig[1.0]{2022_OBS-SKY_O7-s1.pdf}
\ioaafig{2022_OBS-SKY_O7-s2.pdf}

\probhead{15.44}{Daytime observation O8}{2022}{Kutaisi, Georgia}{OBS-SKY O8}{30}
% subject: Chapter 1 (formerly Problem 1.75)
% source id: 2022_OBS-SKY_O8   tag: daytime observation
There are three additional and two missing stars on the figure given below.
Circle them and number them as 1, 2, 3 (additional stars), 4, 5 and 6 (missing
stars).
% NOTE: transcribed verbatim -- the original says "two missing stars" yet
% numbers the missing stars 4, 5 and 6.

\ioaafig{2022_OBS-SKY_O8-s2.pdf}

\probhead{15.45}{Daytime observation O9}{2022}{Kutaisi, Georgia}{OBS-SKY O9}{5}
% subject: Chapter 1 (formerly Problem 1.76)
% source id: 2022_OBS-SKY_O9   tag: daytime observation
Calculate the time interval between the two figures given below if the first
figure occurs earlier than the second one, and the time interval is less than
one sidereal day.

\ioaafig{2022_OBS-SKY_O9-s1.pdf}

\probhead{15.46}{Circumference around Corvi on the star chart}{2024}{Rio de Janeiro, Brazil}{OBS-SKY 1}{18}
% subject: Chapter 1 (formerly Problem 1.77)
% source id: 2024_OBS-SKY_1   tag: star chart angular distances
Using the Star Chart shown in Figure 1, draw a circumference with a radius of
$8^\circ$ around $\alpha$ Corvi. You are given that the angular distance
between the stars Kraz ($\beta$ Corvi) and Spica ($\alpha$ Vir) is
$17.2^\circ$. Draw a circle around the stars with apparent magnitudes between
4 and 6 ($4 \le m \le 6$) that are within the circumference. Do not include
$\alpha$ Corvi. You will be penalised for each star you incorrectly identify
or miss in the requested magnitude range.

\ioaafig{2024_OBS-SKY_1-s2.pdf}

\probhead{15.47}{Star chart question 2}{2024}{Rio de Janeiro, Brazil}{OBS-SKY 2}{16}
% subject: Chapter 1 (formerly Problem 1.78)
% source id: 2024_OBS-SKY_2   tag: star chart work
Figure 2 shows an alt-azimuthal projection of the sky, with the zenith at the
centre.

\ioaafig{2024_OBS-SKY_2-s2.pdf}

\begin{description}
\item[(2.1)] (4 points) On the Star Chart (Figure 2), identify and mark the
  Cardinal Points N, E, S, and W.
\item[(2.2)] (6 points) Draw the Celestial Equator, the Local Meridian, and
  the Ecliptic.
\item[(2.3)] (6 points) On the list below, mark an X next to the stars that
  are present on the Star Chart shown in Figure 2.

  \begin{tabular}{l}
  (\ \ ) Procyon ($\alpha$ CMi) \\
  (\ \ ) Denebola ($\beta$ Leo) \\
  (\ \ ) Izar ($\epsilon$ Boo) \\
  (\ \ ) Atria ($\alpha$ TrA) \\
  (\ \ ) Tarf ($\beta$ Cnc) \\
  \end{tabular}
\end{description}

\probhead{15.48}{Star chart question 3}{2024}{Rio de Janeiro, Brazil}{OBS-SKY 3}{16}
% subject: Chapter 1 (formerly Problem 1.79)
% source id: 2024_OBS-SKY_3   tag: star chart work
\ioaafig{2024_OBS-SKY_3-s2.pdf}

\begin{description}
\item[(3.1)] (7 points) Identify the star circled on the map shown in
  Figure 3 (use its Bayer designation) and evaluate the time when the star
  crosses the local meridian, knowing that this map corresponds to 23:00 local
  time.
\item[(3.2)] (4 points) Evaluate at what time $\beta$ Crucis will set below
  the horizon.
\end{description}

\probhead{15.49}{Know your Sky}{2025}{Mumbai, India}{OBS-SKY OM01}{16}
% subject: Chapter 1 (formerly Problem 1.80)
% source id: 2025_OBS-SKY_OM01   tag: sky map object marking
Use the skymap ``Map-OM01'' to answer the questions below.

\ioaafig{2025_OBS-SKY_OM01-s5.pdf}

\begin{description}
\item[(OM01.1)] Mark all the objects listed below with a box ($\square$)
  around the object. Label each of your marked objects with the corresponding
  Object No.

  \begin{center}
  \begin{tabular}{cccc}
  Object No. & Object Name & Object No. & Object Name \\
  \hline
  1 & $\beta$ Aur & 5 & $\delta$ Gem \\
  2 & $\delta$ Cep & 6 & $\beta$ CVn \\
  3 & $\delta$ Cnc & 7 & $\alpha$ Lyn \\
  4 & $\delta$ Cet & 8 & $\beta$ Per \\
  \end{tabular}
  \end{center}

\item[(OM01.2)] Mark the positions of the following 6 galaxies from the
  Messier Catalogue using a plus sign ($+$) and label them with their
  corresponding Messier number.

  M32, M51, M74, M81, M94, M101

\item[(OM01.3)] Draw the Ecliptic on the map and label it as ``E''.

\item[(OM01.4)] A total solar eclipse occurred on 1 August 2008. At a certain
  place on Earth the totality occurred at local noon.
  \begin{description}
  \item[(OM01.4a)] Mark the position of the Sun at the time of the eclipse
    with a cross ($\times$) and label it as ``S''.
  \item[(OM01.4b)] Draw the Moon at the appropriate position on the map as
    seen from the same location on 28 July 2008 at local noon, and label it as
    ``M''. The drawing should be of appropriate shape and orientation, but
    need not be to scale. The bright side of the Moon should be shaded.
  \end{description}
\end{description}

\probhead{15.50}{Know your Grid}{2025}{Mumbai, India}{OBS-SKY OM02}{16}
% subject: Chapter 1 (formerly Problem 1.81)
% source id: 2025_OBS-SKY_OM02   tag: sky map work
Use the skymap ``Map-OM02'' to answer the questions below.

The constellation lines and boundaries (as per IAU standards) of two
constellations denoted as C1 and C2 are shown in the full-sky map. Alternative
depictions of the same constellations according to a few cultures are also
shown on the right panel for your reference, if needed. A certain coordinate
grid is also shown.

\ioaafig{2025_OBS-SKY_OM02-s2.pdf}

\begin{description}
\item[(OM02.1)] Identify the constellations C1 and C2 and write their names
  (of Latin origin) or IAU abbreviations in the table in the Summary
  Answersheet.
\item[(OM02.2)] Three empty red squares and three empty blue circles are shown
  on the map. There is a grid line passing through each of these squares and
  circles.
  \begin{description}
  \item[(OM02.2a)] The lines that pass through the red squares are lines of
    constant Ecliptic latitude ($\beta$) / Ecliptic longitude ($\lambda$) /
    Declination ($\delta$) / Right Ascension ($\alpha$) / Galactic latitude
    ($b$) / Galactic longitude ($l$). Tick ($\checkmark$) the correct option
    in the Summary Answersheet.
  \item[(OM02.2b)] The lines that pass through the blue circles are lines of
    constant Ecliptic latitude ($\beta$) / Ecliptic longitude ($\lambda$) /
    Declination ($\delta$) / Right Ascension ($\alpha$) / Galactic latitude
    ($b$) / Galactic longitude ($l$). Tick ($\checkmark$) the correct option
    in the Summary Answersheet.
  \end{description}
\item[(OM02.3)] Identify the North and South poles of the grid. Label these
  points as ``N'' and ``S'', respectively, on the map ``Map-OM02''.
\item[(OM02.4)] Two of the following are present in the given sky map.
  Identify these by marking with appropriate symbols (shown below) on the
  corresponding entire curve/line.
  \begin{description}
  \item[(i)] Ecliptic (small bars)
  \item[(ii)] Celestial Equator (small circles)
  \item[(iii)] Galactic Equator (small crosses)
  \end{description}
\item[(OM02.5)] Mark the Vernal Equinox (VE) and Autumnal Equinox (AE) on the
  grid with $\otimes$ and write VE and AE beside them, respectively.
\item[(OM02.6)] Indicate the direction of the Sun's annual motion by drawing
  an arrow close to the Vernal Equinox.
\item[(OM02.7)] Write the values, in appropriate units, inside each red square
  and blue circle given on the ``Map-OM02'', of the corresponding grid lines
  passing through them.
\item[(OM02.8)] The location of 4 constellations (apart from C1 and C2) are
  shown on the grid by light-green shaded areas. Consider the following list
  of constellations.

  Aquarius (Aqr), Cygnus (Cyg), Leo (Leo), Orion (Ori), Perseus (Per),
  Sagittarius (Sgr).

  On the map ``Map-OM02'', label the appropriate shaded areas with the IAU
  abbreviations of the constellations that are present in the above list.
  Mark a cross ($\times$) on those shaded areas, if any, which do not appear
  in the above list.
\end{description}

\probhead{15.51}{Know your Time}{2025}{Mumbai, India}{OBS-SKY OM03}{18}
% subject: Chapter 1 (formerly Problem 1.82)
% source id: 2025_OBS-SKY_OM03   tag: sky map work
The given sky map (in Mercator projection) shows the constellation of Orion as
seen from a certain location X (longitude $\lambda_X = 70^\circ$\,E) on
21 March 2025 at 22:00 local time. The point ``W'' marks the cardinal West
point. The altitude and azimuth values are marked on the grid.

\ioaafig{2025_OBS-SKY_OM03-s3.pdf}

\begin{description}
\item[(OM03.1)] What is the approximate latitude ($\phi_X$) of the
  location X?
\item[(OM03.2)] The grid provided on the Summary Answersheet has the same
  projection (Mercator), and the angular scale in both azimuth and altitude
  are identical to that of the grid provided in the question. On this grid
  draw to scale the constellation of Orion as it will be at another location Y
  (with latitude $\phi_Y = 40^\circ$\,S and longitude
  $\lambda_Y = 50^\circ$\,W) on 21 January 2026 at 18:00 local time. An
  approximate outline of the constellation is enough, with the points A--F
  being marked clearly. Identify the cardinal point ``P'' shown on the grid
  (tick ($\checkmark$) the appropriate box in the Summary Answersheet). You
  may make suitable approximations to arrive at your answer.

  You may use the following relations between the Hour Angle ($H$),
  Declination ($\delta$), Altitude ($a$), Azimuth ($A$) and Latitude
  ($\phi$):
  \[ \cos H = \frac{\sin a - \sin\delta \sin\phi}{\cos\delta \cos\phi} \]
  \[ \sin\delta = \sin\phi \sin a + \cos\phi \cos a \cos A \]
\end{description}

\section*{Observation --- telescope}

\probhead{15.52}{CCD imaging and identification of a cluster field}{2008}{Bandung, Indonesia}{OBS-TEL 1}{5}
% subject: Chapter 5 (formerly Problem 5.32)
% source id: 2008_OBS-TEL_1   tag: telescope + CCD photography, object ID
Participants have to identify as many stars as possible in the field of a
celestial photograph using the telescope and CCD provided by the committee.
Participants have to choose one out of the five recommended regions in the
sky listed below. Then, point the telescope to the direction of the selected
sky region. Take three photographs with different exposure times and record
the images of the sky by the CCD camera. Save the observational data.
Transfer the data to the printing facilities to print out the result. Ask
the technical assistant for a help. Choose the best prints-out and use the
image to identify the stars in the field of observations. The following
procedures are,

\begin{enumerate}
\item Choose only one out of the following five recommended regions to the
  directions of (marked by the following bright clusters):
  \begin{description}
  \item[i.] M7 ($\alpha = 17^{\mathrm{h}}53.^{\mathrm{m}}3$,
    $\delta = -34^\circ 46.'2$)
  \item[ii.] M8 ($\alpha = 18^{\mathrm{h}}04.^{\mathrm{m}}2$,
    $\delta = -24^\circ 22.'0$)
  \item[iii.] M20 ($\alpha = 18^{\mathrm{h}}02.^{\mathrm{m}}4$,
    $\delta = -22^\circ 58.'7$)
  \item[iv.] M21 ($\alpha = 18^{\mathrm{h}}04.^{\mathrm{m}}2$,
    $\delta = -22^\circ 29.'5$)
  \item[v.] M23 ($\alpha = 17^{\mathrm{h}}57.^{\mathrm{m}}0$,
    $\delta = -18^\circ 59.'4$)
  \end{description}
  You may not change your choice.
\item Type in your country code, your student code, and your choice of
  region into the answer page provided in the computer.
\item Point the telescope to the chosen cluster by using telescope
  controller. If necessary you may move the telescope slightly to get the
  best position in the frame of the CCD by checking the display of CCDops
  software.
\item Display the region in The Sky map software provided in the computer,
  to confirm that the telescope is pointed to the selected object in the
  sky. You may change field of view of the sky chart.
\item You may invert the background images into white color, as in the chart
  mode. Copy and paste the sky chart from The Sky into the answer page. Use
  ``Ctrl c'' and ``Ctrl v'' buttons from the keyboard, respectively.
\item Type the equatorial coordinates of the centre of that object in the
  answer page as indicated in The Sky.
\item Take three photographs of the chosen object by using the attached CCD
  camera and CCDops software, with various exposure times. Choose exposure
  time in the range between 1 and 120 seconds. Image is automatically
  subtracted with dark frame with the same exposure time.
\item You must invert the background images into white color. Copy and paste
  images from CCDops into the answer page. Use ``Ctrl c'' and ``Ctrl v''
  buttons from the keyboard, respectively.
\item Save your answer page into hard disk in a Microsoft Word file format.
\item Print your answer page which consists of the photographs and the
  corresponding sky chart.
\item Go to the identification room and bring with you the prints-out of the
  sky chart and the photographs. Ask some help from technical assistant, if
  necessary.
\item Choose the best out of the three printed images and identify as many
  objects as possible on it.
\item Use the assigned computer and The Sky software to identify objects.
  Type your identification to your answer page.
\item Type on the answer page the names (or catalogue number), RA, Dec, and
  magnitude of each identified star and put the sequential number on the
  photograph. Make sure that you list the stars on the answer page following
  the same order and number as on the photograph.
\item Estimate the limiting magnitude of the photograph you choose,
  empirically.
\end{enumerate}

\probhead{15.53}{Telescope pointing and field identification}{2009}{Tehran, Iran}{OBS-TEL 3}{80}
% subject: Chapter 5 (formerly Problem 5.33)
% source id: 2009_OBS-TEL_3   tag: telescope pointing, field ID
Note: You have only 13 Minutes to answer all parts of this Question.

Before starting this part, please note: the telescope is pointed by the
examiner towards Caph ($\beta$ Cas). Please note the readings on the grade
circles before moving the telescope (to be used in 3.2).

\begin{description}
\item[(3.1)] Choose one of the 4 recommended stars listed below; write down
  the name of the selected star in table 3-1 and point the telescope to that
  star. Then, notify the examiner to check it. \hfill(40 Points)
  \begin{description}
  \item[1-] Deneb (Alpha Cygni)
  \item[2-] Alfirk (Beta Cephei)
  \item[3-] Algol (Beta Persei)
  \item[4-] Capella (Alpha Aurigae)
  \end{description}
  (You have 6 Minutes to answer 3.1)
\item[(3.2)] The Telescope was parked to Caph in Cassiopeia constellation
  (RA: $0^{\mathrm{h}}9.7^{\mathrm{m}}$; Dec: $59^\circ 12'$). Using the
  clock beside the telescope write down the local time (with the format of
  HH:MM:SS) in the appropriate field in Table B. Then, by using the graded
  circle on the telescope mount, estimate the ``declination'' and the ``hour
  angle'' of the target measured from South, which you chose in part one of
  this question. Then, complete Table 3-2. \hfill(40 Points)

  (You have 7 Minutes to answer 3.2)
\end{description}

\probhead{15.54}{The Little Dolphin}{2011}{Katowice, Poland}{OBS-TEL 1}{25}
% subject: Chapter 1 (formerly Problem 1.83)
% source id: 2011_OBS-TEL_1   tag: asterism location from coordinates
An asterism known as the Little Dolphin lies near a line connecting the stars
$\alpha$ Peg (Markab) and $\beta$ Peg (Scheat). It is marked with a circle on
the large-scale map.
% NOTE: the original says "cicle" here and spells the constellation
% "Delfinus" in the next paragraph; both corrected in this transcription.

The map also shows the constellation of Delphinus, the Dolphin, with the
brightest stars labelled with their Bayer designations ($\alpha$, $\beta$,
$\gamma$, $\delta$ and $\epsilon$).

\ioaafig{2011_OBS-TEL_1-s1.pdf}

The coordinates of $\alpha$ and $\beta$ Peg and the Little Dolphin (in right
ascension order) are:

\begin{center}
\begin{tabular}{lcc}
 & Right Ascension $\alpha$ & Declination $\delta$ \\
\hline
Little Dolphin & $23^{\mathrm{h}}\,02^{\mathrm{m}}$ & $+23.0^\circ$ \\
$\beta$ Peg & $23^{\mathrm{h}}\,04^{\mathrm{m}}$ & $+28.1^\circ$ \\
$\alpha$ Peg & $23^{\mathrm{h}}\,05^{\mathrm{m}}$ & $+15.2^\circ$ \\
\end{tabular}
\end{center}

Based on your observations, make two drawings on the answer sheet.

On Drawing 1:
\begin{itemize}
\item Draw the view of the constellation Delphinus (Del) as seen through the
  finder scope. Include as many stars as you can see in the field of view.
\item With an arrow, mark the apparent direction of motion of the stars across
  the field of view of the finder scope caused by the rotation of the Earth.
\item Label the stars with the Bayer designations given on the map ($\alpha$,
  $\beta$, $\gamma$, $\delta$ and $\epsilon$).
\item Also label the brightest of these 5 stars ``$m_{\max}$''.
\item Also label the faintest of these 5 stars ``$m_{\min}$''.
\end{itemize}

On Drawing 2:
\begin{itemize}
\item Draw the view of the Little Dolphin as seen through the larger
  telescope. Include as many stars as you can see in the field of view.
\item With an arrow, mark the apparent direction of motion of the stars across
  the field of view of the telescope caused by the rotation of the Earth.
\item Label the stars of the Little Dolphin $\alpha'$, $\beta'$, $\gamma'$,
  $\delta'$ and $\epsilon'$ such that they match the labels of the stars in
  the constellation Delphinus as given on the map.
\item Label the brightest of these 5 stars ``$m_{\max}$''.
\end{itemize}

\probhead{15.55}{Determining declination}{2011}{Katowice, Poland}{OBS-TEL 2}{25}
% subject: Chapter 1 (formerly Problem 1.84)
% source id: 2011_OBS-TEL_2   tag: drift timing gives declination
The two pictures on the next page show a small asterism, as seen directly on
the sky and as a mirror image. Three stars are labelled: S1, S2 and Sx. The
position of the asterism is also marked with a rectangle on the larger-scale
map of the sky.

\ioaafig{2011_OBS-TEL_2-s1.pdf}

\begin{itemize}
\item Find this asterism and point your telescope to it.
\item Using the illuminated reticle as a fixed reference point, and the
  stopwatch, measure the time taken for the stars S1, S2 and Sx to move across
  the field. You may rotate the eyepiece so that the cross-hairs of the
  reticle are in the most convenient position for your measurement.
\item Use your measurements and the known declinations of stars S1 and S2 as
  given below to determine the declination of star Sx.
\item On the answer sheet, give your measurements and working, and estimate
  the random error in your result.
\end{itemize}

For each set of measurements you make, draw the view through the eyepiece on
the answer sheet. (Use the blank circular field on the answer sheet.)

Mark the drawing with the compass directions N and E. Draw the reticle and the
tracks of the stars to show the motion which you timed using the stopwatch.

Mark the ends of each timed track and show which time measurement refers to
which track --- for example, for measurement ``T1'' marking the ends
``Start T1'' and ``End T1''.

The angle of the reticle can be easily adjusted by rotating the eyepiece
around its optical axis. If you change the angle of the reticle for a new
measurement, draw a new diagram.

The declinations of the field stars S1 and S2 are:
\[ \mathrm{S1}: \ \delta = +19^\circ 48' 18'' \qquad
   \mathrm{S2}: \ \delta = +20^\circ 06' 10'' \]

Assume that: $\delta(\mathrm{S2}) > \delta(\mathrm{Sx}) >
\delta(\mathrm{S1})$.

\probhead{15.56}{Field of view of the telescope (NGC 6231)}{2012}{Rio de Janeiro, Brazil}{OBS-TEL 1}{}
% subject: Chapter 5 (formerly Problem 5.34)
% source id: 2012_OBS-TEL_1   tag: estimating eyepiece field of view
Estimate the field of view of this telescope, using 10\,mm Pl\"oss eyepiece
and star chart-1, showing nearby region of open cluster NGC 6231. Star chart
1 shows two angular distances. Use them as reference. Express your answer in
arc minutes and tenths of it.

\ioaafig{2012_OBS-TEL_1-s1.pdf}

\probhead{15.57}{Magnitude of a missing star in NGC 6231}{2012}{Rio de Janeiro, Brazil}{OBS-TEL 2}{}
% subject: Chapter 6 (formerly Problem 6.32)
% source id: 2012_OBS-TEL_2   tag: visual magnitude estimation
Use \emph{star chart-2} to estimate the magnitude of the missing star, shown
as a cross, inside NGC 6231. Use the magnitude of other stars as reference.

\textbf{Note}: To avoid confusion between decimal dots and real stars, dots
where supressed. % sic: "where supressed" as printed in the source
So, magnitude 60 corresponds to magnitude 6.0. Give your
answer using one decimal figure and 0.1 precision.

\ioaafig{2012_OBS-TEL_2-f1.pdf}
% ($IOAA_PAPERS_DIR/observation/2012_OBS.pdf, page 3): star chart of
% NGC 6231 with the missing star marked as a cross and magnitudes of
% reference stars printed without decimal points.

\probhead{15.58}{Colours of the binary epsilon Trianguli Australis}{2012}{Rio de Janeiro, Brazil}{OBS-TEL 3}{}
% subject: Chapter 8 (formerly Problem 8.27)
% source id: 2012_OBS-TEL_3   tag: visual colour estimation of a binary
Point your telescope to the binary star $\varepsilon$-Trianguli Australis
using star chart-3 as a guide. That pair's components are magnitude 4.1 and
9.3, separated by $82''$. Choose the best option for the correct color of
each star:

\begin{description}
\item[Brighter:] White/blue (\quad) \qquad Yellow (\quad) \qquad
  Red (\quad)
\item[Dimmer:] White/blue (\quad) \qquad Yellow (\quad) \qquad
  Red (\quad)
\end{description}

\probhead{15.59}{Identify open/globular clusters and nebulae}{2012}{Rio de Janeiro, Brazil}{OBS-TEL 4}{}
% subject: Chapter 11 (formerly Problem 11.16)
% source id: 2012_OBS-TEL_4   tag: visual object classification
Identify objects pointed out by the evaluator as an Open Cluster (OC), Globular
Cluster (GC), Emission Nebula (EN) or Planetary Nebula (PN), for each of four
objects.

\probhead{15.60}{Night observation round}{2014}{Suceava, Romania}{OBS-TEL 1}{}
% subject: Chapter 5 (formerly Problem 5.35)
% source id: 2014_OBS-TEL_1   tag: telescope observation tasks
For the observational test outdoor, a Newtonian telescope on equatorial
mount EQ5 ($D = 200$\,mm, $F = 1000$\,mm) is used. Note: the telescope is
already aligned, but not necessarily calibrated --- do not change the
position of the tripod! The observational round in the field should take a
maximum of 30 minutes.

\begin{enumerate}
\item Name any five constellations which will be at the meridian 2 hours
  from the start of your observation test?
\item Point the telescope to M39. When you finish, ask your assistant to
  verify. Write in the box the number which corresponds to the object type
  (1 - globular cluster, 2 - double cluster, 3 - open cluster, 4 - galaxy,
  5 - nebula).
\item The right ascension and the declination for $\beta$ Aql (Alshain) are
  $\alpha = 19^{\mathrm{h}}55^{\mathrm{m}}$ and $\delta = 6^\circ 26'$. By
  using the telescope find out the right ascension and the declination for
  $\delta$ Cep. Write down the values in the appropriate boxes.
\item Point the telescope to the coordinate
  $\alpha = 2^{\mathrm{h}}22^{\mathrm{m}}$ and $\delta = 57^\circ 10'$. When
  you finish, ask your assistant to verify. Write in the box the number
  which corresponds to the object type (1 - globular cluster, 2 - double
  cluster, 3 - open cluster, 4 - galaxy, 5 - nebula).
\item Estimate UT when the meridian, the ecliptic and the equator are
  intersecting at the same point in this night. You may use the telescope or
  any other method.
\item Estimate the galactic latitude of $\xi$ Dra (Grumium).
\item Estimate the ecliptic latitude of $\varepsilon$ Cyg (Gienah).
\end{enumerate}

\probhead{15.61}{Main telescope: point to M7 and observe}{2015}{Magelang, Indonesia}{OBS-TEL 1}{}
% subject: Chapter 5 (formerly Problem 5.36)
% source id: 2015_OBS-TEL_1   tag: telescope pointing and observation
Main Telescope (Total Duration: 20 minutes)

\begin{enumerate}
\item Point your telescope to NGC 6475 (M7)\\
  RA (J2000): $17^{\mathrm{h}}\,53^{\mathrm{m}}\,54^{\mathrm{s}}$\\
  Dec (J2000): $-34^\circ 49'$\\
  Show the object to your telescope assistant. If you fail to find the
  object in the first 10 minutes, ask the assistant to help, but you will
  lose the mark. \hfill(23 points)
\item Write down the cardinal directions (North, South, East and West) on
  Chart 1 in the answer sheet. \hfill(10 points)
\item Three stars of M7 are missing from the Chart 1. Mark the position of
  the missing stars with crosses (x). \hfill(21 points)
\item The apparent magnitudes of comparison stars A, B and C are 7.6, 7.2,
  and 5.6 respectively. Estimate and write down the apparent magnitudes of
  the missing stars beside respective cross marks. \hfill(21 points)
\item Estimate the FOV of the instrument, using the given stopwatch. Show
  your calculation in the answer sheet. \hfill(25 points)
\end{enumerate}

\ioaafig{2015_OBS-TEL_1-s1.pdf}

\probhead{15.62}{Telescope already set: identify the field}{2016}{Bhubaneswar, India}{OBS-TEL OT1}{10}
% subject: Chapter 5 (formerly Problem 5.37)
% source id: 2016_OBS-TEL_OT1   tag: telescope field identification
When you arrive at your observing station, DO NOT disturb the telescope
before attempting this question.

The telescope is already set to a deep sky object. Identify the object and
tick the correct box in the Summary Answersheet.

\emph{Note:} You can use any technique to identify the object. However, if
you disturb the telescope, you will NOT be helped to bring it back to the
original position.

\probhead{15.63}{Telescope observation task}{2016}{Bhubaneswar, India}{OBS-TEL OT2}{20}
% subject: Chapter 5 (formerly Problem 5.38)
% source id: 2016_OBS-TEL_OT2   tag: telescope observation
\begin{description}
\item[(OT2.1)] Point the telescope to M45. Show the object to the
  examiner. \hfill(5 points)

  \emph{Note:} 1. After 5 minutes, 1 mark will be deducted for a delay of
  every minute (or part thereof) in pointing the telescope.
  2. You have a single chance to be evaluated. If your pointing is incorrect
  the examiner will change the pointing to M45 for the next part of the
  question.
\item[(OT2.2)] Your Summary Answersheet shows the telescopic field of M45.
  In the image, the seven (7) brightest stars of the cluster are replaced by
  a `+' sign. Compare the image with the field you see in the telescope and
  number the `+' marks from 1 to 7 in the order of decreasing brightness
  (brightest is 1 and faintest is 7) of the corresponding stars. \hfill(15 points)
\end{description}

\textit{[M45 answersheet image not in available sources]}

\probhead{15.64}{Telescope observation task}{2016}{Bhubaneswar, India}{OBS-TEL OT3}{20}
% subject: Chapter 5 (formerly Problem 5.39)
% source id: 2016_OBS-TEL_OT3   tag: telescope observation
The examiner will give you a moon filter, an eyepiece with a cross-wire and
a stopwatch. Point the telescope towards the Moon. Attach the filter to the
telescope. On the surface of the Moon, you will see several ``seas'' (maria)
which are nearly circular in shape. Estimate the diameter of Mare
Serenitatis, $D_{\mathrm{MSr}}$, labelled as ``1'' in the figure below, as a
fraction of the lunar diameter, $D_{\mathrm{Moon}}$, by measuring the
telescope drift times, $t_{\mathrm{Moon}}$ and $t_{\mathrm{MSr}}$, for the
Moon and the mare, respectively.

\ioaafig{2016_OBS-TEL_OT3-s1.pdf}

\probhead{15.65}{Night observation with the real sky}{2017}{Phuket, Thailand}{OBS-TEL 1}{}
% subject: Chapter 5 (formerly Problem 5.40)
% source id: 2017_OBS-TEL_1   tag: telescope observation tasks
This part of the exam involves observation of the real sky. You must
complete two tasks using the equipment provided. You have 6 minutes to
complete the first task, and 4 minutes to complete the second task. Once you
complete a task, call out to the proctor to have it graded; once graded,
that answer is considered final and you may not return to it again. Once the
timer has expired any ungraded task will be graded as is.

\textbf{N1: Observation with Equatorial Mount Telescope}

Use the Equatorial Mount Telescope to observe the target given in the star
chart. Write down the Name and make sure the target is focused.

Included:
\begin{itemize}
\item Equatorial Mount Telescope
\item Mount is already polar aligned
\item Eyepiece: 25 mm
\item Finder scope already aligned with the main scope
\item Telescope starts already pointing at ``starting star'' (see map)
\item Star chart with starting and target position marked.
\end{itemize}

Target name: \underline{\hspace{5cm}}

\ioaafig{2017_OBS-TEL_1-s1.pdf}

\textbf{N2: Observation with Dobsonian Telescope}

Use the Dobsonian Telescope provided to observe only one of the following
objects and focus on the selected object.

\begin{center}
\begin{tabular}{cll}
 & Name & Bayer Designation \\
\hline
$\square$ & Menkar & $\alpha$ Cet \\
$\square$ & Markab & $\alpha$ Peg \\
$\square$ & Hamal & $\alpha$ Ari \\
$\square$ & Fomalhaut & $\alpha$ PsA \\
$\square$ & Alpheratz & $\alpha$ And \\
\end{tabular}
\end{center}

Included:
\begin{itemize}
\item Dobsonian Telescope
\item Eyepieces: 12 and 25 mm
\item Finder scope is already aligned with the main scope.
\end{itemize}

\probhead{15.66}{Eyepiece field of view}{2018}{Beijing, China}{OBS-TEL O5}{10}
% subject: Chapter 5 (formerly Problem 5.41)
% source id: 2018_OBS-TEL_O5   tag: field of view measurement
You will use a telescope, equipped with one eyepiece but no finderscope, to
observe five red LED screens at a distance.

The eyepiece's field of view is $45^\circ$. Please estimate the field of
view (with an accuracy of $0.1^\circ$) of the telescope when observing.

\probhead{15.67}{Text observed through the telescope}{2018}{Beijing, China}{OBS-TEL O6}{40}
% subject: Chapter 5 (formerly Problem 5.42)
% source id: 2018_OBS-TEL_O6   tag: telescope resolution/reading
You will use a telescope, equipped with one eyepiece but no finderscope, to
observe five red LED screens at a distance. Some of these screens display
the names of planets, some show Messier numbers, and some show equatorial
coordinates.

Write down the text you observed on the screen.

\probhead{15.68}{Telescope task 1}{2019}{Keszthely, Hungary}{OBS-TEL 1}{5}
% subject: Chapter 5 (formerly Problem 5.43)
% source id: 2019_OBS-TEL_1   tag: telescope observation
\emph{Telescope: 150/750 Newton. Eyepieces: 25 mm, 10 mm, Barlow lens:
2$\times$. The telescope is already polar aligned. Use the 25 mm eyepiece
for this task.}

\textbf{Finderscope alignment} \hfill available time: 5 minutes

The finderscope is NOT aligned at the beginning. Point the telescope to
Saturn and align the finderscope parallel to the main tube.

If the alignment of Saturn is not within the crosshair of the finderscope,
the telescope assistant will correct it --- and you receive only partial or
no points.

\medskip
\emph{Alternative task (used in case of bad sky conditions at low
altitudes):} The finderscope is NOT aligned at the beginning. Point the
telescope to Altair ($\alpha$ Aql) and align the finderscope parallel to the
main tube. If the alignment is not satisfactory, the telescope assistant
will correct it --- and you receive only partial or no points.

\probhead{15.69}{Telescope task 2}{2019}{Keszthely, Hungary}{OBS-TEL 2}{15}
% subject: Chapter 5 (formerly Problem 5.44)
% source id: 2019_OBS-TEL_2   tag: telescope observation
\emph{Telescope: 150/750 Newton. For this task, we recommend using the 10 mm
eyepiece and the 2$\times$ Barlow lens.}

\textbf{Observation of Saturn} \hfill available time: 10 minutes

\begin{itemize}
\item In the upper box, the circle represents the disk of Saturn and the
  horizontal line is the E--W direction on the sky. Pay attention to the
  direction of North (see top right corner).

  Mark the position of Titan by a cross.
\item The smaller box on the bottom right corner of the first box is for
  drawing the rings of Saturn. Again the circle represents the disk of
  Saturn.

  Draw the rings of Saturn in this box with the correct size and
  orientation. Both the outer and inner edges of the ring are necessary, no
  faint ring details or gaps are needed. Keep the orientation of the image
  the same as the orientation in the upper box.
\item Estimate the angular distance (in arcsec) and position angle (in
  degrees) of Titan relative to the center of Saturn. You may do your
  calculations besides the answer.
\end{itemize}

Apparent major axis of the ring: $43''$


\medskip
\emph{Alternative task (used in case of bad sky conditions at low
altitudes): \textbf{Epsilon Lyrae}}

\begin{itemize}
\item Find $\epsilon$ Lyr, and make a drawing of the field of view (with the
  object and other stars) with the 10 mm eyepiece. Label the directions
  North and East by two arrows and mark them as `N' and `E'.
\item Estimate the angular distance between the wide pair
  ($\epsilon_1$--$\epsilon_2$), and estimate the position angle of the same
  pair.
\item Increase the magnification with the 2$\times$ Barlow lens to be able
  to resolve and separate the two close pairs. Estimate the angle (in
  degrees to the nearest integer) subtended by the two close pairs relative
  to each other (the enclosed angle of the two lines going through the two
  narrow pairs). Do not give any PA, only the relative angle of the two
  close pairs. No drawing is needed.
\end{itemize}

If you cannot find $\epsilon$ Lyr, the assistant can point to it for you,
but only after 5 minutes. In this case you will lose the marks for pointing
the telescope to the object. The telescope assistant will check the object at
the end of the 10 min limit. If you are ready sooner, keep the star in the
FOV, and wait for the check.

\probhead{15.70}{Telescope task 3}{2019}{Keszthely, Hungary}{OBS-TEL 3}{10}
% subject: Chapter 5 (formerly Problem 5.45)
% source id: 2019_OBS-TEL_3   tag: telescope observation
\emph{Telescope: 150/750 Newton. Use the 25 mm eyepiece for this task.}

\textbf{M57 --- field of view} \hfill available time: 10 minutes

\begin{itemize}
\item Find the planetary nebula M57 (in constellation Lyra), and put it in
  the centre of the field of view in the main scope.
\item The star chart in the answer sheet shows a part of constellation Lyra.
  In this chart, draw the FOV circle around M57 as accurately as possible.
\end{itemize}

If you cannot find M57, the assistant will help you, but only after 5
minutes. In this case you will lose the marks for pointing to the object.

\ioaafig{2019_OBS-TEL_3-s1.pdf}

\probhead{15.71}{Telescope task 4}{2019}{Keszthely, Hungary}{OBS-TEL 4}{15}
% subject: Chapter 5 (formerly Problem 5.46)
% source id: 2019_OBS-TEL_4   tag: telescope observation
\emph{Telescope: 150/750 Newton. Use the 25 mm eyepiece for this task.}

\textbf{Variable star --- AF Cygni} \hfill available time: 15 minutes

\begin{itemize}
\item Use the given charts of the constellation Cygnus to find the variable
  star AF Cyg.

  The large scale finder chart has normal orientation (N is up, E is to the
  left). The smaller scale chart has `telescope' orientation (S is up, W is
  to the left). Brightnesses of reference stars are given without decimal
  points, e.g.\ `97' means 9.7 magnitude.

  If you do not find AF Cyg, the telescope assistant cannot help you to
  point to it in this task.
\item Estimate the magnitude of AF Cyg by comparing it with the reference
  stars and write it down, with decimal point, at one decimal accuracy
  (i.e.\ 9.7).

  Write the time of your observation in UTC. You may ask the telescope
  assistant for the time in the local time zone (CEST).
\end{itemize}

\textit{[AF Cyg finder charts not in source PDFs]}

\probhead{15.72}{Telescope task 5}{2019}{Keszthely, Hungary}{OBS-TEL 5}{5}
% subject: Chapter 5 (formerly Problem 5.47)
% source id: 2019_OBS-TEL_5   tag: telescope observation
\emph{You are not allowed to use the telescope for this task.}

\textbf{Naked eye brightness estimation} \hfill available time: 5 minutes

\begin{itemize}
\item Estimate the visual magnitude of the two naked-eye stars marked on the
  stellar chart of constellation Ursa Minor:
  \begin{description}
  \item[(a)] $\zeta$ UMi (zeta UMi = Alifa) --- STAR 2
  \item[(b)] $\gamma$ UMi (gamma UMi = Pherkad) --- STAR 1
  \item[(c)] Write your estimate with one decimal accuracy (e.g.\ 8.6).
  \end{description}
\item Estimate the angular distance between $\gamma$ UMi (STAR 1) and
  Polaris in degrees.
\end{itemize}

\textit{[UMi chart not in source PDFs]}

\probhead{15.73}{Night observation O1}{2022}{Kutaisi, Georgia}{OBS-TEL O1}{15}
% subject: Chapter 5 (formerly Problem 5.48)
% source id: 2022_OBS-TEL_O1   tag: telescope observation
\emph{Time available: 5 min}

Your task is to identify the circled objects in the images that you will see
through the telescopes (numbered 1--5). The images are posted in the windows
of the academic building and the telescopes are already pointed to each
image. A single object will be marked on each board. For each object in the
list below, if it was marked on one of the boards, write the corresponding
telescope number in the list. You will get 30 seconds per telescope.

\begin{center}
\begin{tabular}{lc|lc}
Object & Number & Object & Number \\
\hline
Jupiter & & $\beta$ Her / Antilicus & \\
Saturn & & $\alpha$ Oph / Rasalhague & \\
Mars & & $\epsilon$ Peg / Enif & \\
$\alpha$ And / Alpheratz & & $\alpha$ Per / Mirfak & \\
$\alpha$ Aql / Altair & & $\alpha$ Sco / Antares & \\
$\alpha$ Boo / Arcturus & & $\alpha$ Ser / Unukalhai & \\
$\alpha$ CrB / Alphecca & & $\epsilon$ Sgr / Kaus Australis & \\
$\beta$ Dra / Rastaban & & $\beta$ UMi / Kochab & \\
\end{tabular}
\end{center}

\probhead{15.74}{Night observation O2}{2022}{Kutaisi, Georgia}{OBS-TEL O2}{15}
% subject: Chapter 5 (formerly Problem 5.49)
% source id: 2022_OBS-TEL_O2   tag: telescope observation
\emph{Time available: 4 min}

Determine the field of view FOV of the given telescope using the 25 mm
eyepiece. Show your calculations for the method used.

\probhead{15.75}{Night observation O3}{2022}{Kutaisi, Georgia}{OBS-TEL O3}{6}
% subject: Chapter 5 (formerly Problem 5.50)
% source id: 2022_OBS-TEL_O3   tag: telescope observation
\emph{Time available: 2 min}

Using the result from problem O2, determine the focal length of the
telescope. Assume the apparent field of view of the 25 mm eyepiece is 45
degrees.

Show and explain your calculations.

\probhead{15.76}{Night observation O4}{2022}{Kutaisi, Georgia}{OBS-TEL O4}{14}
% subject: Chapter 5 (formerly Problem 5.51)
% source id: 2022_OBS-TEL_O4   tag: telescope observation
\emph{Time available: 4 min}

$\beta$ Cyg (Albireo) is a visual binary star. You will be shown a picture
of this binary system from one of the windows of the building nearby. Aim
the telescope at this picture, and estimate the acute angle between the
horizon and the line going through the system's stars.

\probhead{15.77}{Observation round: the alien saucer screens}{2023}{Katowice, Poland}{OBS-TEL 1}{50}
% subject: Chapter 5 (formerly Problem 5.52)
% source id: 2023_OBS-TEL_1   tag: telescope observation tasks
\textbf{General instructions.} Scientists have discovered a crashed alien
flying saucer. High up inside the hold, they found several screens
transmitting views of the sky and telescopes have been set up to let you see
them clearly from the level of the deck. Use your telescope to observe the
(simulated) targets on the screens and record your results.

There are 5 screens on the opposite side: the central one will display video
for tasks 1 and 2, the other four will display static images for tasks 3 and
4. The two screens closer to the centre will display the (same) image for
task 3, and the two outer screens will display the (same) image for task 4.
Point your telescope at the screens furthest away from you.

You will have a total of 30 minutes to complete the observing tasks, however
tasks 1 and 2 will only be displayed once: just as with real observations
you will only have one opportunity to collect the data. There will be two
clocks visible showing the time remaining in the round.

At the start of the round a clock on the central screen will show the
simulated time at the observer's location. The clock will have the correct
orientation when seen through the telescope. The time will be shown for 3
minutes after which it will disappear; use this to set a start time for your
observations.

Caution: the scale of the field of view is different between the video and
still images.

\medskip
\textbf{Observation 1: `Asteroid occultation'} (15 points)

Calculations based on the orbital elements predict that an asteroid will
occult the star HD 163390 for 21 s, with the maximum occultation (mid-time)
occurring at 23:03:32 UT. However, the ephemeris is not perfect and the
prediction may be wrong by up to 20 s for the time and by 10 s for the
duration.

Based on your observations, find the true mid-time and duration of the
occultation. To identify the star use Map 1 and the following coordinates:
\[ \mathrm{HD\ 163390:}\quad \mathrm{RA} = 17^h58^m05^s, \quad
   \mathrm{DEC} = -18^\circ 50' 46.14'' \]
The map and the sky are in the same epoch.

[Answer table: mid-time of occultation $\pm$ error; duration of occultation
$\pm$ error]

\ioaafig{2023_OBS-TEL_1-s1.pdf}

\medskip
\textbf{Observation 2: `Starlink'} (15 points)

In the same star field as for Question 1, a `train' of Starlink satellites
will appear near the meridian of $17^h59^m$ at around 23:05 UT. Their
passage will last for around three minutes.

You may assume that the centre of the star field is at an altitude of
$20^\circ$ and that the satellites are 400 km above the Earth's surface
moving on circular orbits with equal distances between them. You may also
assume that satellites will move vertically (perpendicular to the horizon).

\begin{description}
\item[(a)] Measure the angular velocity of the satellites as seen by an
  observer on the simulated sky.
\item[(b)] Measure the time interval between the passes of successive
  satellites and mark their path on the sky chart (Map 1).
\item[(c)] Calculate the theoretical angular velocity of the satellites as
  seen by the observer, using the information given in the question.
\item[(d)] Estimate the distance in km between two consecutive satellites.
\end{description}

Constants: $G = 6.674 \times 10^{-11}\,\mathrm{N\,m^2\,kg^{-2}}$;
$M_\oplus = 5.972 \times 10^{24}$\,kg; $R_\oplus = 6378$\,km.

\medskip
\textbf{Observation 3: `Planetary Moons'} (10 points)

The screen will display an image of one of the planets of the Solar System
as seen on August 15, 2023, at 00:00 UT. Identify any five moons and mark
them on the answer sheet (you may use the moon position chart attached below
and the table showing their brightness).

\ioaafig{2023_OBS-TEL_1-s2.pdf}

\begin{center}
\begin{tabular}{cll}
Number & Name & Magnitude \\
\hline
I & Mimas & 13.0 \\
II & Enceladus & 11.8 \\
III & Tethys & 10.4 \\
IV & Dione & 10.6 \\
V & Rhea & 9.9 \\
VI & Titan & 8.5 \\
VII & Hyperion & 14.4 \\
VIII & Japetus & 11.0 \\
\end{tabular}
\end{center}

[Answer sheet: image of the planet on which the positions of any 5 moons are
to be marked with a dot and labelled with their numbers (I, II, \ldots)]

\medskip
\textbf{Observation 4: `Supernova'} (10 points)

The other screen presents the view of a galaxy and a bright (mag $< 11$)
object which was not visible previously. Estimate the right ascension (RA)
and declination (DEC) coordinates of this star and estimate its magnitude.
You may use Map 2, with stellar coordinates and a list of magnitudes.

[Answer table: right ascension, declination, estimated magnitude]

\ioaafig{2023_OBS-TEL_1-s3.pdf}

\probhead{15.78}{Discovery of 'Z'}{2025}{Mumbai, India}{OBS-TEL OT01}{25}
% subject: Chapter 5 (formerly Problem 5.53)
% source id: 2025_OBS-TEL_OT01   tag: sky map + telescope object discovery
You are given a sky map ``Map-OT01'' along with the question paper. This map
shows only stars but no diffuse objects.

\ioaafig{2025_OBS-TEL_OT01-s1.pdf}

Four known stars (S1, S2, S3 and S4), which are present in the map above,
are given in the table below with their Common names, Bayer designations and
Equatorial coordinates.

\begin{center}
\begin{tabular}{cllll}
Sr. No. & Common Name & Bayer Name & RA & Dec \\
\hline
S1 & Alpheratz & $\alpha$ Andromedae & $00^h08^m24^s$ & $29^\circ 5'16''$ \\
S2 & Markab & $\alpha$ Pegasi & $23^h04^m46^s$ & $15^\circ 12'17''$ \\
S3 & Scheat & $\beta$ Pegasi & $23^h03^m47^s$ & $28^\circ 4'58''$ \\
S4 & Algenib & $\gamma$ Pegasi & $00^h13^m14^s$ & $15^\circ 10'59''$ \\
\end{tabular}
\end{center}

Complete the tasks (OT01.1) and (OT01.2) while you are planning the
observations.

\begin{description}
\item[(OT01.1)] Your first task is to mark these 4 stars (with a circle
  around each star) and label them as S1, S2, S3 and S4 on the sky map
  ``Map-OT01'' provided to you. \hfill[6]
\item[(OT01.2)] An astronomer has discovered a new diffuse object `Z' at the
  following coordinates ---
  RA: $21^h36^m10.6^s$, Dec: $-26^\circ 10'24.4''$.

  Mark the position of this diffuse object on the same sky map ``Map-OT01''
  with a $\oplus$ sign and label it as `Z'. Assume that a linear rectangular
  grid for equatorial coordinates is valid in the relevant region of the
  map. \hfill[7]
\end{description}

The following task is to be performed once you reach the telescope station.

On the screen diagonally opposite to your station, initially a welcome
message, followed by a sample sky (sky unrelated to the question) will be
displayed along with a countdown timer. You may use this time to orient the
telescope towards the screen and familiarize yourself with other equipment
provided at the station. At the end of this time a part of the sky given in
the map ``Map-OT01'' will be projected on the screen for the next 6 minutes.
Note that the scale of the projection shown on the screen is different from
the actual scale seen in the sky.

\begin{description}
\item[(OT01.3)] Find the new object `Z' with the telescope using any
  appropriate eyepiece. Then centre the object properly in the field of view
  of the eyepiece with the crosshair, and show it to the examiner at your
  station.

  At the end of 6 minutes, the projection will be blurred for 20 seconds. At
  this point you must step away from the telescope. The projection will be
  restored for the examiner to check the view through the telescope. This
  marks the end of the first task. \hfill[12]
\end{description}

\probhead{15.79}{Lensing Time Delay}{2025}{Mumbai, India}{OBS-TEL OT02}{25}
% subject: Chapter 5 (formerly Problem 5.54)
% source id: 2025_OBS-TEL_OT02   tag: telescope observation
Gravitational lensing can result in multiple images of a background source
if the source, the lensing object and the observer are nearly aligned. These
multiple images take different times to reach the observer, and if the
background source is variable, each image shows the same feature in its
variability after specific time delays. These time delay measurements are
extremely useful to estimate the current expansion rate of the Universe, the
Hubble constant.

We will consider the gravitational lens system shown in the above figure.
The left hand panel shows a galaxy cluster (lens) together with 4 images of
a background quasar formed due to gravitational lensing. The 4 images,
labelled A, B, C and D, have different fluxes as each image is magnified by
a different amount. For any given image the magnification does not change
with time. Light takes the largest time to travel for the image labelled D.

\ioaafig{2025_OBS-TEL_OT02-s1.pdf}

The light coming from this quasar is variable, and astronomers have
monitored this system for more than a decade now. The right hand panel of
the figure shows the light curve for the image D.

On the screen opposite to your station you will see a movie of the
gravitational lens system. This movie is 28 s long and will loop 6 times
with breaks of 1 minute or 2 minutes between runs. Every second on the watch
corresponds to 250.0 days in the actual lens system.

\begin{description}
\item[(OT02.1)] Let the time delays of image D with respect to images A, B
  and C be given as $t_{DA} = t_D - t_A$, $t_{DB} = t_D - t_B$, and
  $t_{DC} = t_D - t_C$, respectively. Find these time delays, taking any
  necessary steps to reduce the uncertainty in your results. \hfill[25]
\end{description}

\probhead{15.80}{Saturn}{2026}{Hanoi, Vietnam}{OBS-TEL OT01}{40}
% subject: Chapter 5
% source id: 2026_OBS-TEL_OT01   tag: telescope pointing, zenith distance, ring axis angle, hour angle
\begin{description}
\item[(a)] (10 points) Point the telescope at Saturn, then call your examiner
  for confirmation and marking after using the highest magnification
  available. If you are unable to point at the target, the examiner will do it
  for you, but you will get no points.
\item[(b)] (5 points) Write down your observing time and estimate the zenith
  distance $z$ of Saturn.
\item[(c)] (11 points) Write down your observing time and estimate the acute
  angle of the rotation axis of Saturn relative to the horizon of the
  observing site.
\item[(d)] (14 points) The declination of Saturn is $+2^\circ$, assume that
  its diurnal circle is great circle, determine its hour angle $HA$ at your
  observing time.
\end{description}

\probhead{15.81}{Messier 29}{2026}{Hanoi, Vietnam}{OBS-TEL OT02}{15}
% subject: Chapter 5
% source id: 2026_OBS-TEL_OT02   tag: star-hopping with the atlas to M29
Using the attached sky atlas page, point the telescope at Messier 29, then
call your examiner for confirmation and marking after using the highest
magnification available.

\ioaafig{2026_OBS-TEL_OT02-f1.pdf}

\probhead{15.82}{Reverberation Mapping}{2026}{Hanoi, Vietnam}{OBS-TEL OT03}{20}
% subject: Chapter 10
% source id: 2026_OBS-TEL_OT03   tag: echo-mapping flash timing, black hole mass
Reverberation mapping (or echo mapping) is a technique used to estimate the
mass of a black hole at the center of an Active Galactic Nucleus by observing
the variations in light emitted from the system. In this setup, a repeating
dot flash represents a sudden flare from the accretion disk, and a stationary
concentric ring appearing after a while represents the light echoing through
the gas region around the central black hole.

\begin{description}
\item[(a)] (5 points) Target Acquisition: Point your telescope toward the top
  floor of the Phenikaa Tower (the highest building on the Phenikaa campus),
  where the central dot and ring flashes will appear. Call your examiner for
  confirmation and marking. (If you cannot find the target, you may ask your
  examiner to point the telescope for you for 0 points in this part.)
\item[(b)] (10 points) Measure the time delay $(\tau)$ between the dot flash
  and the subsequent ring at least 3 times (up to 5 times). Record all
  individual readings, and calculate the mean time delay and its uncertainty.
\item[(c)] (5 points) Assume the gas orbiting the black hole with velocity of
  $3000\,\mathrm{km\,s^{-1}}$ and one measured second corresponding to one
  day. Calculate the black hole's mass $M$ (in $\mathrm{M}_\odot$).
\end{description}

\clearpage
\section*{Solutions to Chapter 15}

\solhead{15.1}{Earth}
% source id: 2011_OBS-PLAN_1
Note: for questions requiring giving a position, correct to $1^\circ$ gets
full marks (usually 1 point), correct to $2^\circ$ gets half marks (usually
0.5 points).

\begin{description}
\item[(A)] The nova, the Moon and the comet are marked on the sky map (see
  the marked answer map). \hfill(2+1+3 points)
  \ioaafig{2011_OBS-PLAN_1-s2.pdf}
\item[(B)] Objects above the astronomical horizon (to be circled):
  $o$ Cet -- Mira; $\delta$ CMa -- Wezen; $\alpha$ Cyg -- Deneb;
  $\beta$ Per -- Algol; $\delta$ Cep -- Alrediph; M44 -- Praesepe (Beehive
  Cluster). Below the horizon (not circled): M20 -- Trifid Nebula;
  M57 -- Ring Nebula; $\alpha$ Boo -- Arcturus.
  \hfill(+0.5 for each correct, $-1$ for each incorrect)
\item[(C)] The northern part of the local meridian and the ecliptic north
  pole are marked on the answer map. \hfill(3+2 points)
\item[(D)] Geographical latitude of the observer:
  $\varphi = 50^\circ \pm 1^\circ$; \hfill(1 point)\\
  Local Sidereal Time:
  $\theta = 5^{\mathrm{h}}25^{\mathrm{m}} \pm 8^{\mathrm{m}}$;
  \hfill(2 points)\\
  time of year: Jan or Feb. \hfill(2 points)
\item[(E)] $A_1 = 45^\circ$, $h_1 = 58^\circ$: M45; \hfill(1 point)\\
  $A_2 = 278^\circ$, $h_2 = 20^\circ$: $\alpha$ Leo. \hfill(1 point)
\item[(F)] Sirius ($\alpha$ CMa):
  $A_3 = 340^\circ \pm 2^\circ$, $h_3 = 20^\circ \pm 1^\circ$;
  \hfill(0.5+0.5 points)\\
  the Andromeda Galaxy (M31):
  $A_4 = 108^\circ \pm 2^\circ$, $h_4 = 41^\circ \pm 2^\circ$.
  \hfill(0.5+0.5 points)
\item[(G)] $\alpha = 3^{\mathrm{h}}24^{\mathrm{m}} \pm 15^{\mathrm{m}}$;
  $\delta = 49^\circ \pm 3^\circ$. \hfill(1+1 points)
\end{description}

\solhead{15.2}{Mars}
% source id: 2011_OBS-PLAN_2
\begin{description}
\item[(H)] Areographic (Martian) latitude of the observer:
  $\varphi = 22.5^\circ$. \hfill(2 points)
\item[(I)] Pollux ($\beta$ Gem):
  $h_u = 61^\circ \pm 0.5^\circ$;
  $h_l = -50^\circ \pm 1^\circ$ (calculated). \hfill(1+2 points)\\
  Deneb ($\alpha$ Cyg):
  $h_u = 32^\circ \pm 0.5^\circ$;
  $h_l = 13^\circ \pm 0.5^\circ$. \hfill(1+1 points)
\item[(J)] Regulus ($\alpha$ Leo):
  $\delta = -22.5^\circ \pm 0.5^\circ$, from culmination
  $h_u = 45^\circ \pm 0.5^\circ$; \hfill(2 points)\\
  Toliman ($\alpha$ Cen):
  $\delta = -48^\circ \pm 0.5^\circ$, from culmination
  $h_u = 19.5^\circ \pm 0.5^\circ$. \hfill(2 points)
\item[(K)] Sketch diagrams on the meridian plane (zenith--nadir circle with
  the S--N horizon line): the celestial pole is drawn at altitude $\varphi$
  above the north point; the star's declination $\delta$ is laid off from
  the celestial equator to give the culmination altitude $h$.
  \hfill(2+2 points)
  \ioaafig{2011_OBS-PLAN_2-s1.pdf}
\item[(L)] The Martian celestial North Pole is marked ``M'' on the answer
  map (in Cygnus, near the Cyg/Cep border). \hfill(2 points)
  \ioaafig{2011_OBS-PLAN_1-s2.pdf}
\item[(M)] $A = 330^\circ - 180^\circ = 150^\circ$. \hfill(2 points)
\item[(N)] (b) near the northern Tropic circle. \hfill(3 points)
\item[(O)] The date is marked on the time axis at the beginning of northern
  spring. \hfill(3 points)
  \ioaafig{2011_OBS-PLAN_2-s2.pdf}
\end{description}

\solhead{15.3}{Planetarium part I}
% source id: 2014_OBS-PLAN_1
\begin{description}
\item[Question 1.]~
  \begin{description}
  \item[(a)] First arc: Meridian, segment size $10^\circ$;
    \hfill(1+2 points)\\
    second arc: Equator, segment size $15^\circ$. \hfill(1+2 points)
  \item[(b)] $\theta_{\mathrm{sidereal}} = 13^{\mathrm{h}}30^{\mathrm{m}}$
    ($\pm 15^{\mathrm{m}}$ full points; $\pm 30^{\mathrm{m}}$ half points).
    \hfill(6 points)
  \item[(c)] Month number: June (6). \hfill(8 points)
  \end{description}
\item[Question 2.] (15 points)

  \begin{center}
  \begin{tabular}{lccc}
  & Messier object & Type & Constellation \\
  \hline
  1st & M101 & 1 (galaxy) & UMa \\
  2nd & M57 & 2 (nebula) & Lyr \\
  3rd & M92 & 4 (globular cluster) & Her \\
  \end{tabular}
  \end{center}

  (2 points for each Messier number, 1 point for each type, 2 points for
  each constellation.)
\item[Question 3.] (15 points)

  \begin{center}
  \begin{tabular}{lccc}
  & Star & Type & Constellation \\
  \hline
  1st & $\beta$ UMi / Kochab & 1 (single) & UMi \\
  2nd & $\gamma$ Leo / Algieba & 2 (double) & Leo \\
  3rd & $\alpha$ CVn / Cor Caroli & 2 (double) & CVn \\
  \end{tabular}
  \end{center}

  (2 points for each star name, 1 point for each type, 2 points for each
  constellation.)
\end{description}

\solhead{15.4}{Three missing bright stars in the Magelang sky}
% source id: 2015_OBS-PLAN_1
The three missing bright stars are $\alpha$ Cen, $\alpha$ Vir and
$\alpha$ Boo. On the official answer chart (all-sky circle with S at top,
W at left, E at right, N at bottom; the Moon plotted near the centre for
reference) they are marked: $\alpha$ Cen towards the south, $\alpha$ Vir
in the west at mid-altitude, and $\alpha$ Boo in the north-west.

\ioaafig{2015_OBS-PLAN_1-s1.pdf}

\solhead{15.5}{Position of a bright supernova}
% source id: 2015_OBS-PLAN_2
The supernova position is marked on the official answer chart (all-sky
circle with S at top, W at left, E at right, N at bottom): the accepted
position lies a little east of the zenith, between the Moon (plotted just
west of the chart centre) and $\alpha$ Eri, with $\alpha$ Aql and
$\alpha$ Vir also plotted for reference.

\ioaafig{2015_OBS-PLAN_2-s1.pdf}

\solhead{15.6}{Rise and set times of a pointed star}
% source id: 2015_OBS-PLAN_3
The pointed star is Canopus.

\begin{center}
\begin{tabular}{ll}
Rise time & 16:33:57 \\
Set time & 05:51:10 \\
\end{tabular}
\end{center}

Full mark: $\pm 20$ minutes; half mark: $\pm 40$ minutes; quarter mark:
$\pm 60$ minutes.

\solhead{15.7}{Eight historical supernovae}
% source id: 2016_OBS-PLAN_OP1
\textit{No official solution was published for the 2016/2017 observation rounds.}

\solhead{15.8}{Sky of another planet}
% source id: 2016_OBS-PLAN_OP2
\textit{No official solution was published for the 2016/2017 observation rounds.}

\solhead{15.9}{Sky above Keszthely at midnight}
% source id: 2019_OBS-PLAN_1
\begin{description}
\item[(1.1)] The positions of the three novae are marked on the official
  answer map. \hfill(3 points each, maximum 9 points)
  \ioaafig{2019_OBS-PLAN_1-s1.pdf}
\item[(1.2)] Visible Messier globular clusters: M2, M13, M15, M30, M56,
  M71, M72, M92. (Checked with a transparent mask: inner circle $5^\circ$ and
  outer circle $10^\circ$ in diameter; X inside the inner circle 2 points,
  outer circle 1 point; correct Messier number 1 point each.)
  \hfill(maximum 24 points)
  % note: the checking-foil map on page 4 also carries a mask at M79 and
  % marks M75 (at the map edge) with "NO"
\item[(1.3)] October (the projected sky is for October 19, 2019).
  October = 2 points, September or November = 1 point.
  \hfill(maximum 2 points)
\item[(1.4)] Sidereal time is $1^{\mathrm{h}}00^{\mathrm{m}}$. Suggested
  method: $\gamma$ Cassiopeiae is near upper culmination; its RA is
  $0^{\mathrm{h}}57^{\mathrm{m}}$.
  ($1^{\mathrm{h}}00^{\mathrm{m}}$ = 3 points; $\pm 15^{\mathrm{m}}$ = 2
  points; $\pm 30^{\mathrm{m}}$ = 1 point.) \hfill(maximum 3 points)
\item[(1.5)] Sagittarius (Sgr), Capricornus (Cap), Aquarius (Aqr), Pisces
  (Psc), Aries (Ari), Taurus (Tau), Gemini (Gem), Cancer (Cnc).
  \hfill(1 point each, maximum 6 points)
\end{description}

\solhead{15.10}{Planetarium projected image 2}
% source id: 2019_OBS-PLAN_2
\begin{description}
\item[(2.1)] Geographical latitude of the observing site: $25^\circ$,
  Southern hemisphere. The constellation Crux shows the direction of the
  southern celestial pole; the altitude of the pole gives the latitude.
  (Latitude between $24^\circ$--$26^\circ$ = 2 points, within $\pm 3^\circ$
  = 1 point; Southern hemisphere = 1 point.) \hfill(maximum 3 points)
\item[(2.2)] % note: the official solution measures azimuth from South
  % towards West (0--360), unlike the answer-sheet wording
  1st brightest star: Sirius ($\alpha$ CMa), $A = 82^\circ$
  ($74^\circ$--$90^\circ$ = 2 points, $64^\circ$--$74^\circ$ and
  $90^\circ$--$100^\circ$ = 1 point);\\
  2nd brightest star: Canopus ($\alpha$ Car), $A = 42^\circ$
  ($34^\circ$--$50^\circ$ = 2 points, $24^\circ$--$34^\circ$ and
  $50^\circ$--$66^\circ$ = 1 point);\\
  3rd brightest star: Toliman ($\alpha$ Cen), $A = 331^\circ$
  ($324^\circ$--$339^\circ$ = 2 points, $314^\circ$--$324^\circ$ and
  $339^\circ$--$349^\circ$ = 1 point).\\
  (Star name 2 points each; incorrect position in the list $-1$ point.)
  \hfill(maximum 12 points)
\item[(2.3)] The number two comet is in the closest position to the
  ecliptic. (Distances of the comets from the ecliptic: No.\ 1 =
  $20^\circ$, No.\ 2 = $0^\circ$, No.\ 3 = $30^\circ$.) The line
  connecting Regulus and Spica approximately shows the direction of the
  ecliptic. \hfill(2 points)
\item[(2.4)] 1 point for each of: Octans (Oct), Mensa (Men), Chamaeleon
  (Cha), Apus (Aps), Hydrus (Hyi), Volans (Vol), Carina (Car), Musca
  (Mus), Circinus (Cir), Triangulum Australe (TrA), Pavo (Pav), Indus
  (Ind), Tucana (Tuc), Horologium (Hor), Ara (Ara).
  \hfill(maximum 9 points)
\item[(2.5)] 12 hours. $\delta$ Orionis is located approximately on the
  celestial equator ($\delta = -0^\circ 18'$). \hfill(2 points)
\end{description}

\solhead{15.11}{Planetarium projected image 3}
% source id: 2019_OBS-PLAN_3
\begin{description}
\item[(3.1)] Autumn (fall). The Moon, Earth and Sun are located on the same
  line; the Sun is seen from Earth in the constellation Virgo, next to the
  autumnal equinox point. \hfill(2 points)
\item[(3.2)] $598.8$ ($\approx 599$) million km. Juno, the Sun, the Earth
  and the Moon are nearly on the same line, so the distance is
  $3\,\mathrm{au} + 1\,\mathrm{au} + d(\mbox{Earth--Moon}) \approx
  4 \times 149.6 + 0.4 = 598.8$ million km.
  (2 points for $\pm 1$ million km, 1 point for $\pm 2$ million km.)
  \hfill(maximum 2 points)
\item[(3.3)] Sunset: about 3--4 days later (in mean solar days). The length
  of the day on the Moon equals the synodic month, 29.53 days (708.72
  hours), so the Sun moves $360^\circ/29.53 = 12.2^\circ$ per day on the
  lunar sky from east to west. The Sun is $40^\circ$ above the horizon and
  the site is on the lunar equator, so $40^\circ/12.2^\circ$ days
  $= 3.28$ days $= 78.7$ hours are needed for setting. (The extent of the
  solar disc is not considered; that adds about half an hour to complete
  sunset.) (Between 3--4 days: 4 points; 2 or 5 days: 2 points.)
  \hfill(maximum 4 points)\\
  Earth set: never. The Earth--Moon system is in synchronous rotation; the
  observing site is on the near side of the Moon, outside the libration
  zone, so the Earth is continuously observable. \hfill(2 points)
\item[(3.4)] Latitude $0^\circ$, longitude $50^\circ$\,E. If we sat at the
  origin of the selenographic coordinate system, the solar eclipse by the
  Earth would happen at the zenith; since the event is seen almost exactly
  towards the west at a horizontal altitude of $40^\circ$, the site is on
  the lunar equator, about $50^\circ$ east of the origin.
  (Latitude within $\pm 2^\circ$ = 3 points, $\pm 5^\circ$ = 2 points;
  longitude within $\pm 2^\circ$ ($48^\circ$--$52^\circ$) = 3 points,
  $\pm 5^\circ$ = 2 points.) \hfill(maximum 6 points)\\
  The area is Mare Fecunditatis (``Sea of Fecundity'' or ``Sea of
  Fertility''). \hfill(6 points)
\item[(3.5)] About 800 km. Disregarding the latitudinal difference of the
  Apollo-11 landing site from the observing site (it is less than
  $1^\circ$), both sites can be taken on the lunar equator, a great circle
  of circumference $D_{\mathrm{Moon}}\pi = 3474\,\mathrm{km} \times 3.14 =
  10914$ km. This corresponds to $360^\circ$, and the longitudinal
  difference between the sites is $50^\circ - 23.473^\circ = 26.527^\circ$,
  % sic -- the key uses 23.473 deg here although the problem gives the
  % Apollo-11 longitude as 23.4333 deg E
  so the arc length is 803.7 km.
  ($\pm 100$ km: 6 points; $\pm 200$ km: 5 points; $\pm 300$ km: 4 points;
  $\pm 400$ km: 3 points.) \hfill(maximum 6 points)
\end{description}

\solhead{15.12}{Simulated Sky}
% source id: 2021_OBS-PLAN_1
\begin{description}
\item[(1.1)] Line 1: Local Meridian; line 2: Celestial Equator; line 3:
  Ecliptic. \hfill(2+2+2 points)
\item[(1.2)] B. $17^\circ$ \hfill(3 points)
\item[(1.3)] Mizar and Spica. \hfill(4 points)
\item[(1.4)] A. Before sunrise to the east. \hfill(2 points)
\item[(1.5)] C. 8 h \hfill(5 points)
\item[(1.6)] D. 21:00h \hfill(4 points)
\item[(1.7)] A. Ori, Tau, Cet, Eri, Psc, Aqr, Aql, Ser \hfill(3 points)
\item[(1.8)] Star 1: Altair; star 2: Achernar; star 3: Canopus.
  \hfill(6 points)
\item[(1.9)] M15 (Peg): YES; M16 (Aql): YES; M27 (Vul): YES; M6 (Sco):
  YES; M15 (Cas): NO; M27 (CMi): NO. \hfill(6 points)
\item[(1.10)] Star 1: Sham (Sge); star 2: The Persian (Ind); star 3:
  Miaplacidus (Car). \hfill(3 points)
\item[(1.11)] B. Lib, Sco, Oph, Sgr \hfill(4 points)
\item[(1.12)] D. Corona Australis (CrA) \hfill(3 points)
\item[(1.13)] D. M8 -- Lagoon Nebula \hfill(3 points)
\item[(1.14)] E. M102 -- Spindle Galaxy \hfill(5 points)
\item[(1.15)] D. 2.5 (2 points); C. 2 or E. 3 earn 1 point each.
  \hfill(2 points)
\item[(1.16)] Star 1: Gomeisa; star 2: Aludra; star 3: Alzirr; star 4:
  $\alpha$ Mon. \hfill(8 points)
\item[(1.17)] Object 1: M48; object 2: M46; object 3: M50; object 4: M41.
  \hfill(8 points)
\end{description}

\solhead{15.13}{Planetarium round}
% source id: 2023_OBS-PLAN_1
\begin{description}
\item[(a)] Constellation: Aqr; right ascension:
  $22{:}30 \pm 0.5^{\mathrm{h}}$; declination: $-15^\circ \pm 8^\circ$.
  \hfill(1 point for each field, total 3)
\item[(b)] 1 point for each correct DIM/BRIGHT (total 8):

  \begin{center}
  \begin{tabular}{lccc}
  Name & atlas mag. & mag. range & DIM / BRIGHT \\
  \hline
  $\gamma$ Cas (Cih) & 2 & 1.6 -- 3.0 & DIM \\
  $\delta$ Cep & 4 & 3.5 -- 4.4 & DIM \\
  $\mu$ Cep (Erakis) & 4 & 3.4 -- 5.1 & DIM \\
  $\beta$ Per (Algol) & 2 & 2.2 -- 3.4 & DIM \\
  $o$ Cet (Mira) & 3.5 & 2.0 -- 10.1 & BRIGHT \\
  $\chi$ Cyg & 4.5 & 3.3 -- 14.1 & BRIGHT \\
  L$_2$ Pup & 4.5 & 2.6 -- 6 & BRIGHT \\
  $\delta$ Sco (Dschubba) & 2 & 1.6 -- 2.3 & BRIGHT \\
  \end{tabular}
  \end{center}
\item[(c)] Cet (1 pt), Cae (1 pt), Pup (1 pt), Crt (1 pt), Mus (1 pt),
  Ser (2 pt), CrA (1 pt), Equ (1 pt). \hfill(total 9)
\end{description}

\solhead{15.14}{Retrograde Mars}
% source id: 2023_OBS-PLAN_2
% Official solution (2023_PLAN_sol.pdf, Planetarium Round 2 'Retrograde Mars', p.71).
\begin{description}
\item[(a)] 2 points for correct date of opposition and 1 point for each other answer
(total 8)
\begin{center}
\begin{tabular}{|ll|l|}
\hline
i. & the dates of quadrature (when the elongation of Mars is $90^\circ$)
  & 8 Nov 1915 $\pm$ 5 days\\
 & & 15 May 1916 $\pm$ 5 days\\
\hline
ii. & the date of the beginning of retrograde motion & 1 Jan 1916 $\pm$ 5 days\\
 & and the date of the end of retrograde motion & 21 Mar 1916 $\pm$ 5 days\\
\hline
iii. & the date of opposition & 20 Feb 1916 $\pm$ 5 days\\
\hline
iv. & the ecliptic latitude at opposition & $4.5^\circ \pm 1^\circ$\\
\hline
v. & the width in ecliptic longitude of the loop made by the planet
  & $19^\circ \pm 2^\circ$\\
\hline
\end{tabular}
\end{center}

\item[(b)] radius 1.5 au.
% No worked construction is printed; the answer sheet (p.70) shows only the blank
% heliocentric diagram (Earth's orbit with the Sun at centre and the Earth on it) on which
% the geometric construction was to be drawn.

\item[(c)] inclination $1.5^\circ$.
\end{description}

\solhead{15.15}{TRAPPIST-1}
% source id: 2023_OBS-PLAN_3
\begin{description}
\item[(a)]~
\begin{center}
\begin{tabular}{lll}
\toprule
i. & length of the sidereal day [h] & $221\pm5$\\
ii. & orbital period [h] & $221\pm5$ (the same)\\
iii. & length of the `solar' day [h] & infinity\\
iv. & circular orbit & YES\\
v. & obliquity (axial tilt) & $0$ ($\le1^\circ$)\\
\bottomrule
\end{tabular}
\end{center}
\textit{For determining the orbital period -- 3 points, for other quantities
-- 1 point each.}

\item[(b)]~
\begin{center}
\begin{tabular}{lcccc}
\toprule
 & b & c & d & e\\
\midrule
synodic period [h] & $43.6\pm2$ & $78.9\pm2$ & $173.5\pm2$ & $433.8\pm2$\\
maximum elongation [$^\circ$] & $17.5\pm2$ & $24\pm2$ & $37\pm2$ & $49\pm2$\\
\bottomrule
\end{tabular}
\end{center}

\item[(c)]~
\begin{center}
\begin{tabular}{lcccc}
\toprule
 & b & c & d & e\\
\midrule
orbital period [h] & $36.2\pm2$ & $58.1\pm2$ & $97.2\pm2$ & $166.4\pm2$\\
semi-major axis [tau] & $0.30\pm.02$ & $0.41\pm.02$ & $0.60\pm.03$
  & $0.75\pm.04$\\
\bottomrule
\end{tabular}
\end{center}

\item[(d)]~
\begin{center}
\begin{tabular}{ll}
\toprule
Resonance & Pair of planets\\
\midrule
3:2 & e/d and f/e\\
8:5 & c/b\\
5:3 & d/c\\
8:3 & d/b\\
4:1 & e/b\\
6:1 & f/b\\
\bottomrule
\end{tabular}
\end{center}
\textit{Each resonance for 0.5 point, apart from f/e, which is worth 1
point.}
\end{description}

\solhead{15.16}{Planetarium problem 1}
% source id: 2024_OBS-PLAN_1
Note: the total of 20 points is scaled up to 75.

\begin{description}
\item[Question 1] Zodiacal constellations partially or completely above the
  horizon:

  \begin{center}
  \begin{tabular}{ll}
  Pisces & Psc \\
  Aquarius & Aqr \\
  Capricornus & Cap \\
  Sagittarius & Sgr \\
  Ophiuchus & Oph \\
  Scorpius & Sco \\
  Libra & Lib \\
  Virgo & Vir \\
  \end{tabular}
  \end{center}

  (0.5 point for each correct answer; $-0.8$ point for each constellation
  mentioned wrongly.)
\item[Question 2]~
  \begin{description}
  \item[(a)] Azimuth $254^\circ 23'$, altitude $68^\circ 22'$:
    $\alpha$ Sco. \hfill(2 points)
  \item[(b)] Azimuth $113^\circ 24'$, altitude $24^\circ 30'$:
    $\alpha$ PsA. \hfill(2 points)
  \end{description}
\item[Question 3]~
  \begin{description}
  \item[(a)] $\alpha$ Lyr. \hfill(1 point)
  \item[(b)] M8 or NGC 6523, or M20 or NGC 6514 --- any of these answers
    is correct. \hfill(1 point)
  \end{description}
\end{description}

\solhead{15.17}{Planetarium problem 2}
% source id: 2024_OBS-PLAN_2
\begin{description}
\item[Question 1]~
  \begin{description}
  \item[(a)] The Moon is located in \textbf{Leo}. \hfill(2 points)
  \item[(b)] Since the Sun set 1 hour and 26 minutes ago in the projected
    sky, the current month is \textbf{March} (21/03). \hfill(2 points)
  \end{description}
\ioaafig{2024_OBS-PLAN_2-s1.pdf}
\ioaafig{2024_OBS-PLAN_2-s2.pdf}
\item[Question 2] The circumpolar stars are
  \[
  \beta\ \mathrm{Car}, \quad \beta\ \mathrm{Hyi}, \quad
  \beta\ \mathrm{Pav}, \quad \upsilon\ \mathrm{Oct}.
  \]
  (1 point for each correct answer.)
\ioaafig{2024_OBS-PLAN_2-s3.pdf}
\item[Question 3] Bright star very close to the satellite's position:
  \begin{description}
  \item[(a)] At maximum brightness (azimuth $2^\circ 27'$, altitude
    $54^\circ 47'$): $\alpha$ CMi. \hfill(1 point)
  \item[(b)] At disappearance (azimuth $311^\circ 19'$, altitude
    $26^\circ 47'$): $\alpha$ Tau. \hfill(1 point)
  \end{description}
\end{description}

\solhead{15.18}{Hypothesized Planet in the Solar System}
% source id: 2025_OBS-PLAN_OP01
\begin{description}
\item[(OP01.1)]
The orbits of the 4 TNOs are elliptical, which can be noticed from careful observations of
their speed. Hence, to estimate lengths of their semi-major axes it is essential to measure
the time period of one complete revolution of each TNO. The time periods of the planets are
\[
\mathrm{O1{:}\ 40\ sec},\quad \mathrm{O2{:}\ 30\ sec},\quad
\mathrm{O3{:}\ 48\ sec},\quad \mathrm{O4{:}\ 34\ sec}
\]
Hence, the correct ascending order is: $\mathrm{O2} < \mathrm{O4} < \mathrm{O1} < \mathrm{O3}$.
% The 24-row credit table for all possible orderings (marking scheme) is omitted.

\item[(OP01.2)]
The orbits of the TNOs can be drawn by marking a few important constellations and noticing
the fact that some prominent or important stars lie along the TNOs' paths. Also the
intersection points of each TNO trajectory with respect to the Earth's orbit (i.e.\ the
Ecliptic) lie near prominent stars.
\ioaafig{2025_OBS-PLAN_OP01-s1.pdf}

Since this map has an angle preserving (Stereoscopic) projection, measuring the angles
directly using a protractor at the intersection points gives the inclination of the orbits
with respect to the Ecliptic.

The measured values of inclination are:
\begin{itemize}
\item $i_1 = 18^\circ$
\item $i_2 = 14^\circ$
\item $i_3 = 20^\circ$
\item $i_4 = 12^\circ$
\end{itemize}
The average inclination $i_{hp} = 16^\circ$.
% The per-TNO trace/angle marking scheme lines are omitted.
\end{description}

\solhead{15.19}{Observer on a Ringed Planet}
% source id: 2025_OBS-PLAN_OP02
\begin{description}
\item[(OP02.1)]
% The solution is the marked sky map: the positions of the 10 moons are indicated by
% scoring circles on the solution overlay.
\ioaafig{2025_OBS-PLAN_OP02-s1.pdf}
For each moon:
\begin{itemize}
\item Centre of the $\otimes$ within the inner circle \hfill 2.0
\item Centre of the $\otimes$ between the inner and outer circles \hfill 1.0
\end{itemize}
Use 4 correctly placed moons for marking. Every incorrectly placed moon beyond 4 gets
$-1.0$, upto a minimum of 0. % sic: "upto"

\item[(OP02.2)]
The different angular extents of the shadow on the inner and outer edges of the ring are
visible on the dome. These edges may be marked on the sky map (using convenient positions
of the stars nearby).

From the grid of the RA lines on the map the angular sizes are measured to be:
\begin{center}
\begin{tabular}{ll}
Inner edge: & $\approx 75^\circ$\\
Outer edge: & $\approx 33^\circ$\\
\end{tabular}
\end{center}
These angles need to be used in the schematic diagram to locate the rings.

Place a protractor with its baseline along the dotted line, and with its origin at P (the
position of the observer).

Find the line of sight to the inner edge that subtends an angle of $75^\circ/2 = 37.5^\circ$
at P. Extend this line to intersect the tangent ray. This point marks one end of the shadow
on the inner ring. Similarly, find the other ends of the shadow on the inner edge.

Draw a circular arc (with centre at the centre of the planet) passing through these two
points of intersection to show the inner edge of the ring. Repeat for the outer edge.
\ioaafig{2025_OBS-PLAN_OP02-s2.pdf}

\item[(OP02.3)]
Once the edges of the ring are drawn, the width can be calculated by measuring the
difference of their distances from the centre of the planet. Use a ruler to measure these
distances. Also measure the radius of the planet on the diagram with the ruler to find the
calibration.

We find,
\begin{itemize}
\item Distance between the two rings along the dotted line: 3.8 cm
\item Radius of the planet on the diagram: 2 cm equivalent to $6.00\times10^4$ km (given)
\end{itemize}
Therefore,
\[
\text{Width of the ring in km:}
= \frac{3.8\,\mathrm{cm}}{2\,\mathrm{cm}}\times 6.00\times10^4\,\mathrm{km}
= 11.4\times10^4\,\mathrm{km}
\]

Tolerance on values of radii of inner and outer edges:
\begin{center}
\begin{tabular}{|c|c|c|}
\hline
Half credit lower limit & Full credit range & Half credit upper limit\\
\hline
$9.45\times10^4$ km & $10.60\times10^4$ km to $13.40\times10^4$ km & $14.55\times10^4$ km\\
\hline
\end{tabular}
\end{center}
\end{description}

\solhead{15.20}{The Sky over Hue}
% source id: 2026_OBS-PLAN_OP01
\textit{No official solution was published for the 2026 planetarium round.}

\solhead{15.21}{Comet Hale-Bopp}
% source id: 2026_OBS-PLAN_OP02
\textit{No official solution was published for the 2026 planetarium round.}

\solhead{15.22}{The Sky from the Far Side of the Moon}
% source id: 2026_OBS-PLAN_OP03
\textit{No official solution was published for the 2026 planetarium round.}

\solhead{15.23}{Pointer: equator, equinox, Pegasus, Aries, Aldebaran offset}
% source id: 2007_OBS-SKY_1
\textit{No official marking scheme for the Part~I pointer tasks was published; the surviving 2007 key covers only the binocular tasks.}

\solhead{15.24}{Binoculars: Hyades limiting magnitude, M31 shape}
% source id: 2007_OBS-SKY_2
\textbf{(2.1)} The correct value of brightness of $\gamma$-Tauri $= 3.6$.

If the answer is 3.5, 3.6 or 3.7 \hfill(5 points)\\
If the answer is 3.3, 3.4, 3.8 or 3.9 \hfill(3 points)\\
If the answer is none of the above \hfill(0 point)

\textbf{(2.2)} Correct shape and size (between 0.5 to 1 degree) --- 3
points. Correct orientation (within 23 degrees of the picture) --- 2
points.

\ioaafig{2007_OBS-SKY_2-s3.pdf}

\solhead{15.25}{Missing bright stars in Perseus and Andromeda}
% source id: 2009_OBS-SKY_1
The four missing bright stars, marked with circles and identified by their
numbers from Table 1-1, are:

\begin{center}
\begin{tabular}{cll}
Number & Common Name & Bayer Name \\
\hline
1 & Mirfak & Alpha Persei \\
2 & Alpheratz & Alpha Andromedae \\
6 & Algol & Beta Persei \\
7 & Almach & Gamma Andromedae \\
\end{tabular}
\end{center}

\ioaafig{2009_OBS-SKY_1-s1.pdf}

Marking (location of each bright star): small circle with correct number,
$+10$ points; large circle with correct number, $+5$ points; small circle
without identifier number, $+5$ points; large circle without identifier
number, 0 points; small or large circle with incorrect identifier number,
0 points.

\solhead{15.26}{Naked-eye magnitude estimates near Cepheus}
% source id: 2009_OBS-SKY_2
% Official solution: "Solution 2" of the 2009 observation round
% ($IOAA_PAPERS_DIR/observation/2009_OBS_sol.pdf, pp. 87-88).

\textbf{Part 1:}

\begin{center}
Table 2

\begin{tabular}{|c|c|}
\hline
\multicolumn{2}{|c|}{\textbf{Angular Distance}}\\
\hline
\textbf{Stars Name} & \textbf{Angular Distance (Degree)}\\
\hline
1 (Errai) and 2 (Alfirk) & $11^\circ\!:09':10'' \sim 11^\circ$\\
\hline
1 (Errai) and 3 (Alderamin) & $18^\circ\!:36':50'' \sim 19^\circ$\\
\hline
\end{tabular}
\end{center}

\textbf{Part 2:}

\begin{center}
Table 4

\begin{tabular}{|c|c|c|}
\hline
\multicolumn{3}{|c|}{\textbf{Magnitude Estimation}}\\
\hline
\textbf{Number} & \textbf{Star Name} & \textbf{Visible Magnitude}\\
\hline
2 & Alfirk & 3.2\\
\hline
3 & Alderamin & 2.4\\
\hline
\end{tabular}
\end{center}

\solhead{15.27}{Laser pointing: Antares, Vega, Altair, Peacock}
% source id: 2012_OBS-SKY_5
\textit{The official key contains no answer content for this task (a timed practical exercise).}

\solhead{15.28}{Star chart part II}
% source id: 2014_OBS-SKY_2
Official marking scheme for the chart part (50 points):

\begin{description}
\item[(1)] Horizon drawn. \hfill(2p)
\item[(2)] The four lines drawn, 2p for each line. \hfill(8p)
\item[(3)] 0.5p for each cardinal point ($0.5\mathrm{p} \times 4 =
  2\mathrm{p}$); 1.4p for each planet ($1.4\mathrm{p} \times 5 =
  7\mathrm{p}$). \hfill(9p)
\item[(4)] Each star name 0.5p, plus 0.5p for the correct position in the
  list. \hfill(4p)

  \begin{center}
  \begin{tabular}{cl}
  \hline
  1 & $\alpha$ Boo (Arcturus) \\
  2 & $\alpha$ Lyr (Vega) \\
  3 & $\alpha$ Aur (Capella) \\
  4 & $\alpha$ Aql (Altair) \\
  \hline
  \end{tabular}
  \end{center}
\item[(5)] 0.4p for each of the 15 constellations. \hfill(6p)
\item[(6)] Each object 1p. \hfill(5p)
\item[(7)] Sidereal time $18^{\mathrm{h}}$; $\pm 15^{\mathrm{m}}$ full
  points, $\pm 30^{\mathrm{m}}$ half points. \hfill(10p)
\item[(8)] Right ascension $20^{\mathrm{h}} \pm 1^{\mathrm{h}}$, 3p;
  declination $+10^\circ \pm 5^\circ$, 3p. \hfill(6p)
\end{description}

\solhead{15.29}{Nova Sagittarii field sketch}
% source id: 2015_OBS-SKY_1
Official marking scheme:

\begin{description}
\item[(1)] Dec and RA grids made: full score.
\item[(2)] Smallest circle, full score; medium circle, half score; largest
  circle, quarter score (concentric tolerance circles are drawn around each
  of the five stars on the key chart).
\item[(3)] Smallest circle, full score; medium circle, half score; largest
  circle, quarter score (tolerance circles around the Moon's position).
\item[(4)] Smallest circle, full score; medium circle, half score; largest
  circle, quarter score (tolerance circles around each Messier object).
\item[(5)] Galactic plane inside the yellow line: full score. Galactic
  plane crosses the yellow line: half score.
\item[(6)] $\pm 2^\circ$, full score; $\pm 5^\circ$, half score.
\end{description}

\ioaafig{2015_OBS-SKY_1-s1.pdf}

\ioaafig{2015_OBS-SKY_1-s2.pdf}

\ioaafig{2015_OBS-SKY_1-s3.pdf}

The key lists the Messier objects in the field of view (any five could be
marked for task 4):

\begin{center}
\begin{tabular}{ccc}
Messier & Magnitude & size ($'$) \\
\hline
7 & 3.3 & 80 \\
6 & 4.2 & 15 \\
25 & 4.6 & 32 \\
22 & 5.2 & 24 \\
23 & 5.5 & 27 \\
8 & 5.8 & 90 \\
21 & 5.9 & 13 \\
16 & 6 & 35 \\
17 & 6 & 46 \\
20 & 6.3 & 29 \\
62 & 6.6 & 14 \\
20 & 6.9 & 11 \\ % sic -- M20 is listed twice in the key table
18 & 6.9 & 9 \\
55 & 7 & 19 \\
19 & 7.2 & 13 \\
69 & 7.7 & 7 \\
54 & 7.7 & 9 \\
9 & 7.9 & 9 \\
70 & 8.1 & 7 \\
24 & 11 & 5 \\ % sic -- M24's listed magnitude and size as in the key
\end{tabular}
\end{center}

\solhead{15.30}{Mark five stars on the map}
% source id: 2016_OBS-SKY_OM1
\textit{No official solution was published for the 2016/2017 observation rounds.}

\solhead{15.31}{Sky map identification}
% source id: 2016_OBS-SKY_OM2
\textit{No official solution was published for the 2016/2017 observation rounds.}

\solhead{15.32}{Sky map identification}
% source id: 2016_OBS-SKY_OM3
\textit{No official solution was published for the 2016/2017 observation rounds.}

\solhead{15.33}{Sky map identification}
% source id: 2016_OBS-SKY_OM4
\textit{No official solution was published for the 2016/2017 observation rounds.}

\solhead{15.34}{Sky map identification}
% source id: 2016_OBS-SKY_OM5
\textit{No official solution was published for the 2016/2017 observation rounds.}

\solhead{15.35}{Day observation stations (7 problems)}
% source id: 2017_OBS-SKY_1
\textit{No official solution was published for the 2016/2017 observation rounds.}

\solhead{15.36}{Day observation problem 1}
% source id: 2018_OBS-SKY_O1
\begin{description}
\item[(1)] See Figure 1 (key chart). The missing planet is P = Mars. The
  missing stars are $\alpha$ Aql, $\alpha$ Cep, $\varepsilon$ Cas and
  $\alpha$ Ari (anticlockwise). 4 points for each; inside the small circle
  = 4 points, between the small and big circle = 2 points.
\ioaafig{2018_OBS-SKY_O1-s1.pdf}
\item[(2)] See Figure 1: going clockwise around the rim the cardinal
  points are E, N, W, S (E at the upper left, N at the upper right, W at
  the lower right, S at the lower left). 1 point for each, 5 points for
  all; between the curve = 1 point.
\item[(3)] From west to east: Ophiuchus (Oph), Serpens (Ser), Aquila (Aql),
  Aquarius (Aqr), Pisces (Psc), Cetus (Cet), Taurus (Tau), Eridanus (Eri),
  Orion (Ori). 2 points for each, 10 points for 5 or more.
\item[(4)] $18.33^\circ$. 5 points for $16^\circ \sim 21^\circ$, 3 points
  for $13^\circ \sim 15^\circ$ and $22^\circ \sim 24^\circ$, 1 point for
  $10^\circ \sim 12^\circ$ and $25^\circ \sim 27^\circ$.
\end{description}

\solhead{15.37}{Day observation problem 2}
% source id: 2018_OBS-SKY_O2
July 7.

20 points for July 5 $\sim$ July 9, 12 points for July 2 $\sim$ July 4 and
July 10 $\sim$ July 12, 6 points for June 29 $\sim$ July 1 and July 13
$\sim$ July 15.

\solhead{15.38}{Day observation problem 3}
% source id: 2018_OBS-SKY_O3
23:00.

20 points for 22:30 $\sim$ 23:30, 12 points for 21:30 $\sim$ 22:29 and
23:31 -- 0:30, 6 points for 20:30 $\sim$ 21:29 and 0:31 -- 1:30.

\solhead{15.39}{Day observation problem 4}
% source id: 2018_OBS-SKY_O4
M1 -- Taurus (Tau); M104 -- Virgo (Vir) or Corvus (Crv); M27 -- Vulpecula
(Vul); M42 (M43 is also acceptable) -- Orion (Ori).

2.5 points for each Messier catalogue number, 2.5 points for each
constellation name.

\solhead{15.40}{Solar Physics: the solar slap}
% source id: 2021_OBS-SKY_1
\begin{enumerate}
\item \textbf{Data selection} (selecting appropriate images showing the
  evolution of the CME) and identification of CME front.

  \begin{center}
  \includegraphics[width=0.48\linewidth]{2021_OBS-SKY_1-s1-panel1.pdf}\hfill
  \includegraphics[width=0.48\linewidth]{2021_OBS-SKY_1-s1-panel2.pdf}\par\smallskip
  \includegraphics[width=0.48\linewidth]{2021_OBS-SKY_1-s1-panel3.pdf}\hfill
  \includegraphics[width=0.48\linewidth]{2021_OBS-SKY_1-s1-panel4.pdf}\par\smallskip
  \includegraphics[width=0.48\linewidth]{2021_OBS-SKY_1-s1-panel5.pdf}\hfill
  \includegraphics[width=0.48\linewidth]{2021_OBS-SKY_1-s1-panel6.pdf}\par\smallskip
  \includegraphics[width=0.48\linewidth]{2021_OBS-SKY_1-s1-panel7.pdf}
  \end{center}

  A total of 7 time intervals are selected tracing the evolution of the CME:
  \begin{enumerate}
  \item 2021/05/28 22:19:00 UT
  \item 2021/05/28 22:54:11 UT
  \item 2021/05/28 23:30:08 UT
  \item 2021/05/29 23:42:08 UT
  % sic: the official solution lists 2021/05/29 23:42:08 UT here, while its
  % summary table gives 2021/05/28 23:42:08 UT for interval 4
  \item 2021/05/29 01:12:16 UT
  \item 2021/05/29 03:36:08 UT
  \item 2021/05/29 05:00:08 UT
  % sic: the summary table gives 2021/05/29 05:36:08 UT for interval 7, but
  % 05:00:08 UT is consistent with the quoted 401 minutes from onset
  \end{enumerate}

  Summary of the information extracted from the images for all different
  selected intervals:

  {\footnotesize\setlength{\tabcolsep}{2.5pt}
  \begin{center}
  \begin{tabular}{c l c c c c c c}
  \toprule
  Interval & Date and Time & \begin{tabular}{@{}c@{}}Time\\ interval\\ (min)\end{tabular}
    & \begin{tabular}{@{}c@{}}Time from\\ CME onset\\ (min)\end{tabular}
    & \multicolumn{2}{c}{\begin{tabular}{@{}c@{}}Coordinates of\\ the CME front\end{tabular}}
    & \begin{tabular}{@{}c@{}}Distance from\\ disc center\\ (arcsec)\end{tabular}
    & \begin{tabular}{@{}c@{}}Distance from\\ disc center\\ (km)\end{tabular}\\
  & & & & $x''$ & $y''$ & &\\
  \midrule
  1 & 2021/05/28 22:19:00 UT & 0   & 0   & 950   & 144     & 961   & 696617\\
  2 & 2021/05/28 22:54:11 UT & 35  & 35  & 2682  & 174     & 2688  & 1948538\\
  3 & 2021/05/28 23:30:08 UT & 36  & 71  & 4579  & 97      & 4580  & 3320520\\
  4 & 2021/05/28 23:42:08 UT & 47  & 83  & 7441  & 116     & 7442  & 5395380\\
  5 & 2021/05/29 01:12:16 UT & 126 & 173 & 14000 & $-891$  & 14028 & 10170535\\
  6 & 2021/05/29 03:36:08 UT & 191 & 317 & 21048 & $-4570$ & 21538 & 15615349\\
  7 & 2021/05/29 05:36:08 UT & 210 & 401 & 27922 & $-7126$ & 28817 & 20892306\\
  \bottomrule
  \end{tabular}
  \end{center}
  % sic: the "Time interval" column of the official table is not consistent
  % with the differences of the successive timestamps (rows 4--7)

  \begin{center}
  \begin{tabular}{c c c c c c c}
  \toprule
  Interval & \begin{tabular}{@{}c@{}}Distance\\ from onset\\ (km)\end{tabular}
    & \begin{tabular}{@{}c@{}}Cumulative\\ velocity, plane\\ of the sky (km/s)\end{tabular}
    & \begin{tabular}{@{}c@{}}Deprojection of\\ velocity towards\\ ParkerSP (km/s)\end{tabular}
    & \begin{tabular}{@{}c@{}}Distance per\\ time interval\\ (km)\end{tabular}
    & \begin{tabular}{@{}c@{}}Velocity on\\ plane of the\\ sky (km/s)\end{tabular}
    & \begin{tabular}{@{}c@{}}Deprojection of\\ velocity towards\\ ParkerSP (km/s)\end{tabular}\\
  \midrule
  1 & 0        & 0   & 0    & 0       & 0   & 0\\
  2 & 1252198  & 596 & 728  & 1252198 & 596 & 728\\
  3 & 2624180  & 616 & 752  & 1371982 & 635 & 775\\
  4 & 4699040  & 944 & 1152 & 2074861 & 736 & 898\\
  5 & 9474195  & 913 & 1114 & 4775155 & 632 & 771\\
  6 & 14919009 & 784 & 958  & 5444814 & 475 & 580\\
  7 & 20195966 & 839 & 1025 & 5276957 & 419 & 511\\
  \bottomrule
  \end{tabular}
  \end{center}
  }
\item \textbf{Plotting the kinetic evolution of the CME.}

  Plots: distance vs time and velocity vs time.

  \ioaafig{2021_OBS-SKY_1-s2.pdf}
\item \textbf{Calculating the velocity of the CME front reaching ParkerSP.}

  The velocity corresponding to the last interval (up to reaching 30 solar
  radii) is 511\,km/s which, according to the considerations of the problem,
  would be the same velocity of the CME front once it reaches the ParkerSP.

  \textbf{Calculating the time spent by the CME front to reach the ParkerSP
  from its onset.}

  The Sun--ParkerSP distance estimated from Figure 1 is approx.\
  $0.66\,\mathrm{AU} = 9.9\times10^{7}\,\mathrm{km} = 141$ solar radii (given
  the information that 10 solar radii is approx.\ $7\times10^{6}$\,km).

  From the previous analysis, it is known that the CME front reaches a
  distance of 30 solar radii after 401 minutes (6.7 hours) from its onset
  (see Table). The distance that the CME travels at the constant speed of
  511\,km/s is expressed as $141-30=111$ solar radii
  $= 7.8\times 10^{7}$\,km, and thus the time spent in that interval is
  42.4 hours.

  Finally, the total time it takes for the CME front to reach the ParkerSP
  from its onset is computed as: $6.7 + 42.4 = 49.1$ hours.
\item \textbf{True and false statements}
  \begin{description}
  \item[A.] A larger number of selected images would give more precise
    information on the evolution of the CME and the calculated physical
    parameters. TRUE
  \item[B.] The analysis and measurements of the CME evolution should
    consider the differential rotation of the Sun, and therefore the
    calculated velocities will be affected. FALSE
  \item[C.] Any software (numerical) misalignment among the images when
    creating the mosaic will have direct effects on the precision of your
    calculations. TRUE
  \item[D.] The different assumptions made in order to construct the model
    displaying in a heliospheric density map in Figure 1, may affect the
    estimation of the Sun--ParkerSP distance. FALSE
  \item[E.] The interaction of the CME-front with the remnant dust left by
    the 2019 Borisov comet broadens and difuminate the images. This
    intensifies the lack of contrast in the images, increasing the
    uncertainty in determining the CME-front and its propagation. FALSE
    % sic: "difuminate" as in the original
  \end{description}
\item \textbf{Kinetic energy calculations}

  Proton:
  \[
  0.5 \times (1.67262\times10^{-27}~\mathrm{kg}) \times (511~\mathrm{m/s})^2
  = 1.4\times10^{-3}~\mathrm{eV}
  \]
  Alpha particle:
  \[
  0.5 \times (6.64\times10^{-27}~\mathrm{kg}) \times (511~\mathrm{m/s})^2
  = 5.4\times10^{-3}~\mathrm{eV}
  \]
  % sic: the official solution uses 511 m/s (instead of the CME speed of
  % 511 km/s); the quoted eV values follow from the m/s figure.
\end{enumerate}

\solhead{15.41}{Daytime observation O5}
% source id: 2022_OBS-SKY_O5
Naming each lens type: % sic -- ``lense'' in the original
\begin{description}
\item[1.] (C) Gregorian \hfill(3 pt)
\item[2.] (A) Newtonian \hfill(3 pt)
\item[3.] (H) Cassegrain \hfill(3 pt)
\end{description}

\solhead{15.42}{Daytime observation O6}
% source id: 2022_OBS-SKY_O6
Filling missing equatorial constellations:
\begin{description}
\item[3.] Aql (Aquila) \hfill(2 pt)
\item[7.] Tau (Taurus) \hfill(2 pt)
\item[8.] Eri (Eridanus) \hfill(2 pt)
\item[9.] Ori (Orion) \hfill(2 pt)
\item[12.] Hya (Hydra) \hfill(2 pt)
\item[13.] Sex (Sextans) \hfill(2 pt)
\item[14.] Leo (Leo) \hfill(2 pt)
\end{description}

Ecliptic constellations are the following:
\begin{itemize}
\item Aqr (Aquarius) \hfill(1 pt)
\item Leo (Leo) \hfill(1 pt)
\item Oph (Ophiuchus) \hfill(2 pt)
\item Psc (Pisces) \hfill(1 pt)
\item Vir (Virgo) \hfill(1 pt)
\item Tau (Taurus) \hfill(1 pt)
\end{itemize}

\solhead{15.43}{Daytime observation O7}
% source id: 2022_OBS-SKY_O7
% solution source: 2022_TEL_sol.pdf page 2 (Sol: O7, 20 pt)
Each object identification is worth 2.5 points and each constellation 1.5
points.

\begin{center}
\begin{tabular}{ccl}
Image & Messier number & Constellation \\
\hline
I & M81 & UMa (Ursa Major) \\
II & M101 & UMa (Ursa Major) \\
III & M27 & Vul (Vulpecua) \\ % sic -- "Vulpecua" for Vulpecula in the source
IV & M42 & Ori (Orion) \\
V & M73 & Aqr (Aquatius) \\ % sic -- "Aquatius" for Aquarius in the source
\end{tabular}
\end{center}

\solhead{15.44}{Daytime observation O8}
% source id: 2022_OBS-SKY_O8
% solution source: 2022_TEL_sol.pdf page 2 (Sol: O8, 30 pt)
For each correct answer the student gets 5 points.

The official solution shows the original (correct) star sheet and the flawed
sheet with the corrections marked: the three additional stars are marked in
red and the three missing stars in green.

\ioaafig{2022_OBS-SKY_O8-s1.pdf}
\begin{center}\textit{Corrected sheet}\end{center}
\ioaafig[1.0]{2022_OBS-SKY_O8-s1-corrected.pdf}

\solhead{15.45}{Daytime observation O9}
% source id: 2022_OBS-SKY_O9
% solution source: 2022_TEL_sol.pdf page 3 (Sol: O9, 5 pt)
The time difference between the initial and final figures is
$6^{\mathrm{h}}\,20^{\mathrm{m}}$.

\solhead{15.46}{Circumference around Corvi on the star chart}
% source id: 2024_OBS-SKY_1
% solution source: 2024_OBS_sol.pdf page 4 (Question 1, 18 points)
There are 12 stars with $4 \le m \le 6$ inside the circumference of radius
$8^\circ$ around $\alpha$ Corvi; they are numbered 1--12 on the official
solution chart. Each star inside the circumference is worth $+1.5$ points;
$-1.5$ points for each other star marked, except for the star indicated with
an arrow on the solution chart (12 stars, total 18 points).

\ioaafig{2024_OBS-SKY_1-s1.pdf}

\solhead{15.47}{Star chart question 2}
% source id: 2024_OBS-SKY_2
% solution source: 2024_OBS_sol.pdf pages 5-6 (Question 2, 16 points)
\textbf{(2.1)} On the Star Chart, identify and mark the Cardinal Points N, E,
S and W. \hfill(4 points: 1 point for each Cardinal Point)

\textbf{(2.2)} Draw the Celestial Equator, the Local Meridian and the
Ecliptic. \hfill(6 points: 2 points for each line)

\ioaafig{2024_OBS-SKY_2-s1.pdf}

\textbf{(2.3)} Present on the Celeste Chart are the stars Denebola, Izar and
Atria. % sic -- "Celeste Chart" in the source
\hfill(6 points: 2 points for each star, $-2$ points for each star outside
the chart)

\begin{center}
\begin{tabular}{l}
(\ \ ) Procyon ($\alpha$ CMi) \\
(X) Denebola ($\beta$ Leo) \\
(X) Izar ($\epsilon$ Boo) \\
(X) Atria ($\alpha$ TrA) \\
(\ \ ) Tarf ($\beta$ Cancri) \\
\end{tabular}
\end{center}

\solhead{15.48}{Star chart question 3}
% source id: 2024_OBS-SKY_3
% solution source: 2024_OBS_sol.pdf pages 7-8 (Question 3, 11 points)
\textbf{(3.1)} The circled star is Alpha Hydrae ($\alpha$ Hya)
\hfill(4 points)

The star crosses the local meridian at 3\,am:
\begin{itemize}
\item 3\,am $\pm 10$\,min \hfill(3 points)
\item 3\,am $\pm 20$\,min \hfill(2 points)
\item 3\,am $\pm 30$\,min \hfill(1 point)
\end{itemize}

\textbf{(3.2)} $\beta$ Crucis sets below the horizon at:
\begin{itemize}
\item 3:45--4:00 \hfill(4 points)
\item 3:30--4:00 \hfill(3 points)
\item 3:15--4:00 \hfill(2 points)
\item 3:00--4:00 \hfill(1 point)
\end{itemize}

\ioaafig{2024_OBS-SKY_3-s1.pdf}

\solhead{15.49}{Know your Sky}
% source id: 2025_OBS-SKY_OM01
% solution source: 2025_SKYMAP_sol.pdf pages 1-4 (OM01: Solution and Marking Scheme)
\textbf{(OM01.1)} The eight objects are marked with boxes on the official
solution map; every correct star gets 0.5 marks. \hfill(4 marks)

\ioaafig{2025_OBS-SKY_OM01-s1.pdf}

\textbf{(OM01.2)} The six galaxies are marked on the official solution map.
The intersection point of each cross must lie inside the inner tolerance
circle for 1.0 mark (0.5 marks if between the inner and outer circles).
\hfill(6 marks)

\ioaafig{2025_OBS-SKY_OM01-s2.pdf}

\textbf{(OM01.3)} The Ecliptic is drawn on the official solution map with
inner and outer tolerance bounds. Full marks require the curve to pass
within the boundaries of 6 or 5 zodiacal constellations (1.0 mark) and to
lie completely inside the inner bounds (1.0 mark). If the curve is clearly
non-monotonic or not smooth, penalty of $-1.0$ mark. \hfill(2 marks)

\ioaafig{2025_OBS-SKY_OM01-s3.pdf}

\textbf{(OM01.4a)} The date when the eclipse occurred is 1st August. During
the start of the month of August the Sun is in the constellation of Cancer.
The Sun is marked with a cross and the label ``S'' inside the acceptance box
in Cancer. \hfill(1 mark)

\textbf{(OM01.4b)} 28 July was 4 days before the eclipse occurred. The Moon
moves around $13^\circ$ every day from west to east. Hence, the Moon will be
around $52^\circ$ to the east of the Sun, almost along the ecliptic, i.e.\
the Moon will be in the constellation of Taurus.

The phase of the Moon will be crescent (anything less than half is
acceptable). The bright part of the Moon should be towards the Sun.
\hfill(3 marks: crescent drawn 0.5; bright part towards Sun 0.5; drawing
within the box 1.0; elongation of Moon from Sun $\approx 52^\circ$ or
3.5\,h 1.0)

\ioaafig{2025_OBS-SKY_OM01-s4.pdf}

\solhead{15.50}{Know your Grid}
% source id: 2025_OBS-SKY_OM02
% solution source: 2025_SKYMAP_sol.pdf pages 5-11 (OM02: Solution and Marking Scheme)
\textbf{(OM02.1)} It is evident from the constellation lines and their
relative positions that
\[ \text{C1: Virgo (Vir)}, \qquad \text{C2: Corvus (Crv)}. \]
\hfill(1 mark: 0.5 for each constellation)

\textbf{(OM02.2a)} Since the constellation boundaries are parallel to the
grid lines, the grid shown in the figure must be the \textbf{Equatorial
Grid}. Circles of declination are centred around the poles. Hence, the
lines that pass through the red squares are declination ($\delta$) lines.
\hfill(1 mark)

\textbf{(OM02.2b)} Lines emanating from the poles and perpendicular to the
declination circles are RA lines. Hence, the lines that pass through the
blue circles are lines of constant RA ($\alpha$). \hfill(1 mark)

\textbf{(OM02.3)} The constellation of Corvus is towards the South of Virgo.
Hence, the pole of the hemisphere in which Corvus lies is the South pole.
(On the map: the left grid pole is ``S'', the right grid pole is ``N''.)
\hfill(1 mark: 0.5 for each pole)

\textbf{(OM02.4)} The two curves present on the map are the Ecliptic (the
dashed sinusoidal curve, marked with small bars) and the Celestial Equator
(the vertical lines midway between the poles, marked with small circles).
\hfill(2 marks: Ecliptic 1.0; Celestial Equator 1.0, or 0.5 if only one
segment of the Celestial Equator is marked)

\textbf{(OM02.5)} Equinoxes are the intersection points of the equator and
ecliptic. The Autumnal equinox (AE) lies in the constellation of Virgo.
Hence the other intersection point is the Vernal equinox (VE).
\hfill(2 marks: equinoxes at the intersections of CE and Ecliptic 1.0;
both equinoxes correctly identified 1.0)

\textbf{(OM02.6)} From the VE the Sun moves northward, so the Sun's motion
is in the direction of increasing RA. \hfill(1 mark)

\textbf{(OM02.7)} The Vernal Equinox (VE) and Autumnal Equinox (AE) were
identified in part (OM02.5). The Right Ascension of the VE is 0\,h, while
the RA of the AE is 12\,h. From part (OM02.6), the Sun moves northward in
the direction of increasing RA, which allows the RA circles to be marked
accordingly. RA is conventionally measured in hours, minutes, and seconds.
The celestial pole has a declination of $90^\circ$, with northern
declinations positive and southern declinations negative.

The values on the official solution map are: red squares (declination)
$+20^\circ$, $-10^\circ$, $-30^\circ$; blue circles (right ascension)
20\,h, 10\,h, 7\,h.
\hfill(3 marks: 0.5 per correct value with appropriate sign and units;
maximum 1.5 for values incorrect but with appropriate sign and units and
consistent with (OM02.2a) and (OM02.2b))

\textbf{(OM02.8)}
\begin{itemize}
\item In the given list there are 3 zodiacal constellations, Aquarius, Leo
  and Sagittarius, which implies that they should lie on the Ecliptic.
\item The Ecliptic passes through the constellation to the top left which is
  about 2 hours ahead in RA from Virgo, hence it is Sagittarius.
\item From the given list Leo and Orion are the constellations which lie on
  the Celestial Equator. Hence, the constellation on the lower right has to
  be Orion.
\item Out of the remaining constellations, i.e.\ Perseus and Cygnus, the
  constellation of Cygnus lies between RA 20\,h and 22\,h.
\item The remaining box is of Ursa Major, which is not present in the given
  list, so it is marked with a cross ($\times$).
\end{itemize}
\hfill(4 marks: 1.0 for each correct identification, including marking
$\times$ or ``UMa'')

\ioaafig{2025_OBS-SKY_OM02-s1.pdf}

\solhead{15.51}{Know your Time}
% source id: 2025_OBS-SKY_OM03
% solution source: 2025_SKYMAP_sol.pdf pages 12-14 (OM03: Solution and Marking Scheme)
\textbf{(OM03.1)} Draw the celestial equator (the line passing through W and
Mintaka ($\delta$ Ori)). Measure the angle $\phi_X$ between the Celestial
Equator and the vertical at $270^\circ$ azimuth.
\[ \phi_X = 21^\circ\,\mathrm{N} \]
Full credit for $20^\circ$ to $22^\circ$; half credit at the limits
$19^\circ$ and $23^\circ$. \hfill(4 marks: value 1.0; drawing the equator
0.5; equator passing through Mintaka (point C) 1.0; indication of
hemisphere either through ``$+$'', ``N'' or ``North'' 1.5)

\ioaafig{2025_OBS-SKY_OM03-s1.pdf}

\textbf{(OM03.2)} The star Mintaka (which lies very close to the Celestial
Equator, i.e.\ $\delta \approx 0^\circ$) has an approximate altitude at X at
22:00\,hrs on 21 March of $a_{X,\mathrm{21Mar},22:00} = 25^\circ$ from the
West horizon.

Therefore its hour angle at 22:00\,hrs local time on 21 March is given by
\[ H_{\mathrm{21Mar},22:00}
   = \cos^{-1}\!\left[\frac{\sin 25^\circ - \sin 0^\circ \sin 20^\circ}
                           {\cos 0^\circ \cos 20^\circ}\right]
   = 63.27^\circ \]
Therefore, the hour angle of Mintaka at 22:00\,hrs local time on 21 Jan 2026
(306 days later) will be
\[ H_{\mathrm{21Jan},22:00}
   = H_{\mathrm{21Mar},22:00} + 306 \times \frac{360^\circ}{365.2564}
   = 63.27^\circ + 301.60^\circ = 364.87^\circ \equiv 4.87^\circ \]
The hour angle of Mintaka at 18:00\,hrs local time on 21 Jan 2026 will be
\[ H_{\mathrm{21Jan},18:00}
   = H_{\mathrm{21Jan},22:00} - 4 \times \frac{360^\circ}{24}
   = 4.87^\circ - 60^\circ = -55.13^\circ \]
The hour angle of an object at a certain local time on a given day will be
the same at all locations.

Since the hour angle is negative, it indicates that Mintaka is to the east
of the meridian at 18:00\,hrs local time. Therefore, the point P represents
the cardinal point East.

The altitude of Mintaka at 18:00\,hrs local time on 21 Jan 2026 at the
location Y will be (since $\delta \approx 0$)
\[ a_{Y,\mathrm{21Jan},18:00}
   = \sin^{-1}\!\left[\cos H_{\mathrm{21Jan},18:00} \cos\phi_Y\right]
   = \sin^{-1}\!\left[\cos(-55.13^\circ)\cos(-40^\circ)\right]
   \approx 26^\circ \]
Now, the location Y lies in the Southern hemisphere. The Celestial Equator
lies towards the North (i.e.\ tilted towards the left of the cardinal point
E in the given projection) near the Eastern horizon in the Southern
hemisphere. Further, the latitude of Y is $40^\circ$\,S, which gives the
inclination of the Celestial Equator w.r.t.\ the vertical at $90^\circ$
azimuth. The position of Mintaka will be nearly on the Celestial Equator at
an altitude of $\approx 26^\circ$.

The orientation of Orion can be inferred by realizing that the shield of
Orion is ahead of the belt as the sky moves.

\ioaafig{2025_OBS-SKY_OM03-s2.pdf}

Marking scheme (14 marks):
\begin{itemize}
\item ``P'' marked as ``E'' \hfill(2.0)
\item Angle of the line joining Mintaka (C) and cardinal point (P) w.r.t.\
  the vertical between $37^\circ$ and $43^\circ$ \hfill(2.0)
  (1.0 if between $35^\circ$--$37^\circ$ or $43^\circ$--$45^\circ$)
\item All 6 of the A--F points lie to the left of the vertical at $90^\circ$
  azimuth \hfill(2.0)
  (1.0 if only 5 points, 0.5 if only 4 points)
\item Altitude of Mintaka (point C) between $24^\circ$ and $28^\circ$
  \hfill(2.0)
  (1.0 if between $22^\circ$--$24^\circ$ or $28^\circ$--$30^\circ$)
\item Shield (point D) at lesser azimuth than belt (point C) \hfill(2.0)
\item Legs (points A and B) are south of the Celestial Equator \hfill(2.0)
\item Scale: distance B--F is between 1.8\,cm and 2.8\,cm \hfill(2.0)
  (1.0 if between 1.3\,cm--1.8\,cm or 2.8\,cm--3.3\,cm)
\end{itemize}

\solhead{15.52}{CCD imaging and identification of a cluster field}
% source id: 2008_OBS-TEL_1
Model answer pages are provided for each of the five recommended regions:
the sky chart image, an example CCD image with the identified stars
numbered, and the identification table.

\ioaafig{2008_OBS-TEL_1-s1.pdf}
\ioaafig{2008_OBS-TEL_1-s2.pdf}
\ioaafig{2008_OBS-TEL_1-s3.pdf}
\ioaafig{2008_OBS-TEL_1-s4.pdf}
\ioaafig{2008_OBS-TEL_1-s5.pdf}
\ioaafig{2008_OBS-TEL_1-s6.pdf}
\ioaafig{2008_OBS-TEL_1-s7.pdf}
\ioaafig{2008_OBS-TEL_1-s8.pdf}
\ioaafig{2008_OBS-TEL_1-s9.pdf}
\ioaafig{2008_OBS-TEL_1-s10.pdf}

\textbf{M7} ($\alpha = 17^{\mathrm{h}}53.^{\mathrm{m}}3$,
$\delta = -34^\circ 46.'2$). Central equatorial coordinates:
$\alpha = 17^{\mathrm{h}}54^{\mathrm{m}}22.3^{\mathrm{s}}$,
$\delta = -34^\circ 47' 15''$. Exposure time: 5\,sec.

\begin{center}
\begin{tabular}{clllc}
Number & Names or Catalogue number & R.A. & Declination & Magnitude \\
\hline
1 & SAO 209436 & $17^{\mathrm{h}}\,54^{\mathrm{m}}\,14.380^{\mathrm{s}}$ & $-34^\circ 43' 39.535''$ & 6.88 \\
2 & GSC 7386:494 & $17^{\mathrm{h}}\,54^{\mathrm{m}}\,09^{\mathrm{s}}$ & $-34^\circ 44' 31''$ & 11.6 \\
3 & GSC 7386:601 & $17^{\mathrm{h}}\,54^{\mathrm{m}}\,06^{\mathrm{s}}$ & $-34^\circ 44' 44''$ & 12.2 \\
4 & GSC 7386:561 & $17^{\mathrm{h}}\,54^{\mathrm{m}}\,04^{\mathrm{s}}$ & $-34^\circ 44' 06''$ & 8.5 \\
5 & GSC 7386:15 & $17^{\mathrm{h}}\,53^{\mathrm{m}}\,55^{\mathrm{s}}$ & $-34^\circ 48' 20''$ & 11.9 \\
6 & GSC 7386:40 & $17^{\mathrm{h}}\,54^{\mathrm{m}}\,13^{\mathrm{s}}$ & $-34^\circ 48' 18''$ & 11.8 \\
\end{tabular}
\end{center}

\textbf{M8} ($\alpha = 18^{\mathrm{h}}04.^{\mathrm{m}}2$,
$\delta = -24^\circ 22.'1$). Central equatorial coordinates:
$\alpha = 18^{\mathrm{h}}04^{\mathrm{m}}51^{\mathrm{s}}.2$,
$\delta = -24^\circ 23' 04''$. Exposure time: 10\,sec.

\begin{center}
\begin{tabular}{clllc}
Number & Names or Catalogue number & R.A. & Declination & Magnitude \\
\hline
1 & GSC 6842:1751 & $18^{\mathrm{h}}\,04^{\mathrm{m}}\,10^{\mathrm{s}}$ & $-24^\circ 20' 53''$ & 12.5 \\
2 & GSC 6846:166 & $18^{\mathrm{h}}\,04^{\mathrm{m}}\,05^{\mathrm{s}}$ & $-24^\circ 23' 01''$ & 12.1 \\
3 & GSC 6846:211 & $18^{\mathrm{h}}\,03^{\mathrm{m}}\,58^{\mathrm{s}}$ & $-24^\circ 25' 35''$ & 12.0 \\
4 & GSC 6846:718 & $18^{\mathrm{h}}\,04^{\mathrm{m}}\,27^{\mathrm{s}}$ & $-24^\circ 23' 52''$ & 12.4 \\
5 & GSC 6842:1578 & $18^{\mathrm{h}}\,04^{\mathrm{m}}\,23^{\mathrm{s}}$ & $-24^\circ 21' 25''$ & 13.4 \\
\end{tabular}
\end{center}

\textbf{M20} ($\alpha = 18^{\mathrm{h}}02.^{\mathrm{m}}4$,
$\delta = -22^\circ 58.'7$). Central equatorial coordinates:
$\alpha = 18^{\mathrm{h}}02^{\mathrm{m}}56^{\mathrm{s}}.7$,
$\delta = -23^\circ 01' 56''$. Exposure time: 10\,sec.

\begin{center}
\begin{tabular}{clllc}
Number & Names or Catalogue number & R.A. & Declination & Magnitude \\
\hline
1 & GSC 6842:631 & $18^{\mathrm{h}}\,02^{\mathrm{m}}\,18^{\mathrm{s}}$ & $-23^\circ 00' 25''$ & 12.4 \\
2 & GSC 6842:240 & $18^{\mathrm{h}}\,02^{\mathrm{m}}\,19^{\mathrm{s}}$ & $-23^\circ 00' 52''$ & 12.8 \\
\end{tabular}
\end{center}

\textbf{M21} ($\alpha = 18^{\mathrm{h}}04.^{\mathrm{m}}2$,
$\delta = -22^\circ 29.'5$). Central equatorial coordinates:
$\alpha = 18^{\mathrm{h}}04^{\mathrm{m}}46^{\mathrm{s}}.3$,
$\delta = -22^\circ 29' 27''$. Exposure time: 10\,sec.

\begin{center}
\begin{tabular}{clllc}
Number & Names or Catalogue number & R.A. & Declination & Magnitude \\
\hline
1 & GSC 6842:985 & $18^{\mathrm{h}}\,04^{\mathrm{m}}\,09^{\mathrm{s}}$ & $-22^\circ 30' 28''$ & 12.9 \\
2 & GSC 6263:306 & $18^{\mathrm{h}}\,04^{\mathrm{m}}\,01^{\mathrm{s}}$ & $-22^\circ 29' 36''$ & 13.1 \\
\end{tabular}
\end{center}

\textbf{M23} ($\alpha = 17^{\mathrm{h}}57.^{\mathrm{m}}0$,
$\delta = -18^\circ 59.'4$). Central equatorial coordinates:
$\alpha = 17^{\mathrm{h}}57^{\mathrm{m}}32^{\mathrm{s}}.1$,
$\delta = -19^\circ 01' 08''$. Exposure time: 10\,sec.

\begin{center}
\begin{tabular}{clllc}
Number & Names or Catalogue number & R.A. & Declination & Magnitude \\
\hline
1 & GSC 6258:809 & $17^{\mathrm{h}}\,56^{\mathrm{m}}\,58^{\mathrm{s}}$ & $-18^\circ 58' 21''$ & 13.2 \\
2 & GSC 6258:825 & $17^{\mathrm{h}}\,56^{\mathrm{m}}\,59^{\mathrm{s}}$ & $-18^\circ 58' 43''$ & 12.5 \\
\end{tabular}
\end{center}

Marking: student who has the highest number of stars identified ($N_m$)
gets 60\%; other students who get $N$ stars: $N \times 0.60/N_m$.
Estimation of the limiting magnitude: 20\% ($\Delta < 1$ mag: full mark,
$1 < \Delta < 2$ mag: half mark, $\Delta > 3$ mag: 0). Good image quality
(unsaturated, not under-exposed): 20\%.

\solhead{15.53}{Telescope pointing and field identification}
% source id: 2009_OBS-TEL_3
\begin{center}
\begin{tabular}{|l|c|c|c|}
\hline
\multicolumn{4}{|c|}{Coordinate of Selected Star} \\
\hline
Star & Hour Angle & R.A. & Dec. \\
\hline
Deneb (Alpha Cygni) & ST$-$RA & $20^{\mathrm{h}}\,41^{\mathrm{m}}$ & $+45^\circ$ \\
\hline
Alfirk (Beta Cephei) & ST$-$RA & $21^{\mathrm{h}}\,28^{\mathrm{m}}$ & $+71^\circ$ \\
\hline
Algol (Beta Persei) & RA$-$ST & $03^{\mathrm{h}}\,08^{\mathrm{m}}$ & $+41^\circ$ \\
\hline
Capella (Alpha Aurigae) & RA$-$ST & $05^{\mathrm{h}}\,18^{\mathrm{m}}$ & $+46^\circ$ \\
\hline
\end{tabular}
\end{center}

(The answer key leaves the ``Local Time:'' and ``Sidereal Time:'' fields
blank, to be filled at the moment of observation; the hour angle of each
star is expressed as ST$-$RA, or RA$-$ST when the star is east of the
meridian.)

\solhead{15.54}{The Little Dolphin}
% source id: 2011_OBS-TEL_1
% solution source: 2011_NIGHT_sol.pdf page 1 (Night Observation Q1 -- Little
% Dolphin, marking scheme). The official solution is the marking scheme; no
% model drawings are given.
\textbf{On Drawing 1:}

Draw the view of the constellation Delfinus (Del) through the finder
scope: % sic -- "Delfinus" for Delphinus in the source
\begin{itemize}
\item all 4 stars of the parallelogram \hfill(1 point)
\item the `tail' \hfill(1 point)
\item correct scale \hfill(1 point)
\item correct shape \hfill(1 point)
\item additional field stars \hfill(1 point)
\end{itemize}

With an arrow, mark the apparent direction of motion of the stars across the
field of view of the finder scope caused by the rotation of the Earth.
\hfill(2 points)

Label the stars with the Bayer designations given on the map ($\alpha$,
$\beta$, $\gamma$, $\delta$ and $\epsilon$). \hfill(1 point)

Also label the brightest of these 5 stars ``$m_{\max}$''. \hfill(1 point)

Also label the faintest of these 5 stars ``$m_{\min}$''. \hfill(1 point)

\textbf{On Drawing 2:}

Draw the view of the Little Dolphin through the larger telescope:
\begin{itemize}
\item all 4 stars of the parallelogram \hfill(4 points)
\item the `tail' \hfill(1 point)
\item correct scale \hfill(1 point)
\item correct shape \hfill(1 point)
\item additional field stars \hfill(1 point)
\end{itemize}

With an arrow, mark the apparent direction of motion of the stars across the
field of view of the finder scope caused by the rotation of the Earth.
% sic -- the source says "finder scope" here although Drawing 2 is the view
% through the larger telescope
\hfill(2 points)

Label the stars of the Little Dolphin $\alpha'$, $\beta'$, $\gamma'$,
$\delta'$ and $\epsilon'$ such that they match the labels of the stars in
the constellation Delphinus as given on the map. \hfill(3 points)

Label the brightest of these 5 stars ``$m_{\max}$''. \hfill(2 points)

Total: 25 points.

\solhead{15.55}{Determining declination}
% source id: 2011_OBS-TEL_2
% solution source: 2011_NIGHT_sol.pdf page 2 (Night Observation Q2 --
% Measuring Declination, marking scheme). The official solution is the
% marking scheme; no worked numerical answer is given.
Drawing of the field of view (any stars):
\begin{itemize}
\item the picture presents the positions of the stars S1, S2 and Sx
  \hfill(3 points)
\item more stars presented, as on the attached maps \hfill(1 point)
\item North direction marked on the picture of the FoV \hfill(1 point)
\item East direction marked on the picture of the FoV \hfill(1 point)
\end{itemize}

Proper method of observations:
\begin{itemize}
\item if the cross is parallel to N--S and E--W \hfill(0 points)\\
  (and the total score will be calculated from the previous points)
\item cross twisted from the N--S and E--W directions \hfill(2 points)
\item proper position angle of the cross corresponding to the formula used
  for the calculations \hfill(4 points)
\end{itemize}
If the student makes measurements in two or more different position angles,
they receive the 4 points if all of them will be corrected. % sic

\begin{itemize}
\item proper formula for the $\delta_x$ calculation \hfill(6 points)
\item several repetitions of the stopwatch measurement \hfill(3 points)\\
  or: the student gives some subjective estimation of the time measurement
  (from watching the transit of the star under the cross etc.)
  \hfill(2 points)
\item proper calculation of $\delta_x$ \hfill(2 points)
\item proper estimation of the final error \hfill(2 points)
\end{itemize}

Total: 25 points.

\solhead{15.56}{Field of view of the telescope (NGC 6231)}
% source id: 2012_OBS-TEL_1
Answer: 24.0 arc minutes ($0.4^\circ$). We will tolerate up to 10\% error
as correct answer. Error between 10 to 20\% will be considered 70\%
correct.

\solhead{15.57}{Magnitude of a missing star in NGC 6231}
% source id: 2012_OBS-TEL_2
% Official solution: item 2 of "Observational exam -- 1st night -- Solutions"
% ($IOAA_PAPERS_DIR/observation/2012_OBS_sol.pdf, p.1).

Exact answer is magnitude 7.9 but we will assume values between 7.7 -- 8.1 as
correct (0.2 magnitude tolerance). Values between 7.5 -- 7.7 or 8.1 -- 8.3
will be considered as 70\% correct. (5 minutes)

\solhead{15.58}{Colours of the binary epsilon Trianguli Australis}
% source id: 2012_OBS-TEL_3
% Official solution: 2012 observational exam, 1st night, task 3
% ($IOAA_PAPERS_DIR/observation/2012_OBS_sol.pdf, page 1).

\begin{description}
\item[Brighter:] White/blue (\quad) \qquad Yellow (X) \qquad
  Red (\quad) \hfill (K0III)
\item[Dimmer:] White/blue (X) \qquad Yellow (\quad) \qquad
  Red (\quad) \hfill (B9IV)
\end{description}

If the student is not able to point to the star the evaluator will do so, but
student loses half of the points.

\solhead{15.59}{Identify open/globular clusters and nebulae}
% source id: 2012_OBS-TEL_4
Targets pointed by the evaluator will be M22, M8, M7 and M57. Student will
have 30 seconds to evaluate each target.

\begin{description}
\item[Object 1 (M22)] (GC)
\item[Object 2 (M57)] (PN)
\item[Object 3 (M7)] (OC)
\item[Object 4 (M8)] (EN)
\end{description}

\solhead{15.60}{Night observation round}
% source id: 2014_OBS-TEL_1
\textit{No official solution for the 2014 outdoor telescope round was published; the surviving 2014 keys cover only the planetarium and sky-chart rounds.}

\solhead{15.61}{Main telescope: point to M7 and observe}
% source id: 2015_OBS-TEL_1
\begin{enumerate}
\item Checklist (Telescope Assistant): Object Accuracy (Right Object $=$
Full Score).

\item Cardinal points: as marked on the solution chart (E at top, W at
bottom, N at left, S at right). Each correct cardinal point: quarter score.

\item Three missing stars: marked on the solution chart. Region 1 (Full
Score), 2 (Half Score), 3 (Quarter Score). Note: the regions are marked
with the circles (region will be chosen during the IBM).

\ioaafig{2015_OBS-TEL_1-s2.pdf}

\item Missing stars:
\begin{description}
\item[1.] 7.4
\item[2.] 6.4
\item[3.] 7.0
\end{description}
Score range: $\pm 0.2$ = Full Score; $\pm 0.5$ = Half Score;
$\pm 0.8$ = Quarter Score.

\item Using the 15\,mm (eyepiece), the field of view of the instrument is:
$0.75^\circ$
\begin{itemize}
\item Using stopwatch to measure the time for a star to travel from the
  center of the FOV to the edge of the FOV. \hfill(Half Score)
\item Using stopwatch to measure the time for a star to travel from the
  edge of the FOV to the other edge of the FOV. \hfill(Quarter Score)
\item Calculating the FOV,
\[ FOV(\mathrm{deg}) = \frac{\text{Travel time (sec)}}{23^{\mathrm{h}}\,56^{\mathrm{m}}}
   \times \cos(\delta) \times 360\,\mathrm{deg} \]
\hfill(Half Score)
\end{itemize}
\end{enumerate}

\solhead{15.62}{Telescope already set: identify the field}
% source id: 2016_OBS-TEL_OT1
\textit{No official solution was published for the 2016/2017 observation rounds.}

\solhead{15.63}{Telescope observation task}
% source id: 2016_OBS-TEL_OT2
\textit{No official solution was published for the 2016/2017 observation rounds.}

\solhead{15.64}{Telescope observation task}
% source id: 2016_OBS-TEL_OT3
\textit{No official solution was published for the 2016/2017 observation rounds.}

\solhead{15.65}{Night observation with the real sky}
% source id: 2017_OBS-TEL_1
\textit{No official solution was published for the 2016/2017 observation rounds.}

\solhead{15.66}{Eyepiece field of view}
% source id: 2018_OBS-TEL_O5
\[ F \tan(\mathrm{FOV}/2) = f \tan(45^\circ/2) \]
$F = 700$\,mm, $f = 26$\,mm
\[ \mathrm{FOV}/2 = \arctan(0.414 \cdot f/F) = 0.881^\circ \]
\[ \mathrm{FOV} = 2 \times 0.881^\circ = 1.76^\circ \approx 1.8^\circ \]
10 points for $1.6^\circ$--$1.9^\circ$, 5 points for $1.3^\circ$--$1.5^\circ$
and $2.0^\circ$--$2.2^\circ$.

\solhead{15.67}{Text observed through the telescope}
% source id: 2018_OBS-TEL_O6
2 planets: Saturn, Jupiter.\\
2 Messier objects: M15, M74.\\
1 equatorial coordinate: $+30^\circ$ $21^{\mathrm{h}}$.

8 points for each text.

\solhead{15.68}{Telescope task 1}
% source id: 2019_OBS-TEL_1
Evaluation of the alignment and marking is done by the telescope
assistants. \hfill Max.: 5 points

\solhead{15.69}{Telescope task 2}
% source id: 2019_OBS-TEL_2
\textbf{Saturn and Titan.} The official answer depends on the night of
observation; reference values are given for both nights. Titan's position is
checked with a transparent template foil: if Titan is within the template
circle, 2 points; if outside, 0 points.

\ioaafig{2019_OBS-TEL_2-s1.pdf}
\ioaafig{2019_OBS-TEL_2-s2.pdf}
\ioaafig{2019_OBS-TEL_2-s3.pdf}
\ioaafig{2019_OBS-TEL_2-s4.pdf}

\medskip
For 2019.08.04 22:00 UT:
\[ d_{\mathrm{Titan}} = 164'' \]
\begin{center}
\begin{tabular}{ll}
$d$: $148$--$180''$ & 3 points \\
$d$: $139$--$189''$ & 2 points \\
$d$: $123$--$205''$ & 1 point \\
\end{tabular}
\end{center}
\[ \mathrm{PA}_{\mathrm{Titan}} = 109^\circ \]
\begin{center}
\begin{tabular}{ll}
PA: $99$--$119^\circ$ & 4 points \\
PA: $89$--$129^\circ$ & 2 points \\
\end{tabular}
\end{center}

For 2019.08.05 22:00 UT:
\[ d_{\mathrm{Titan}} = 129'' \]
\begin{center}
\begin{tabular}{ll}
$d$: $116$--$131''$ & 3 points \\
$d$: $110$--$148''$ & 2 points \\
$d$: $97$--$161''$ & 1 point \\
\end{tabular}
\end{center}
\[ \mathrm{PA}_{\mathrm{Titan}} = 124^\circ \]
\begin{center}
\begin{tabular}{ll}
PA: $114$--$134^\circ$ & 4 points \\
PA: $104$--$144^\circ$ & 2 points \\
\end{tabular}
\end{center}

\medskip
Scoring:

\begin{center}
\begin{tabular}{lll}
Titan position + name: & & 2 points \\
Titan is outside the circle: & & 0 point \\[4pt]
Ring continuity in front of the disk is proper & & 2 points \\
Ring tilt towards Earth is proper & & 2 points \\
\multicolumn{3}{l}{\quad ring does not exceed northern pole, ring is not too thin} \\
Ring tilt towards sky E--W direction is proper & & 2 points \\[4pt]
PA measurement of Titan: & $\pm 10^\circ$ & 4 points \\
 & $\pm 20^\circ$ & 2 points \\[4pt]
Distance calculation of Titan & $\pm 10$\% & 3 points \\
 & $\pm 15$\% & 2 points \\
 & $\pm 25$\% & 1 point \\
\end{tabular}
\end{center}

The Saturn drawing evaluation scheme uses the sub-Earth latitude
$\Phi = 24.8^\circ$ (ring tilt towards Earth, max 2 points: drawings
equivalent to $22.8$--$26.8^\circ$ get 2p; $20.8^\circ$ gets 1p, S pole
visible; $28.8^\circ$ gets 1p, very thick; thinner or thicker gets 0p) and
the equatorial position angle $\mathrm{PA} = 6.4^\circ$ (ring tilt towards
the sky E--W direction, max 2 points: $6.4^\circ$ or $9.4^\circ$ get 2p;
$3.4^\circ$ or $12.4^\circ$ get 1p; $0^\circ$ or $15.4^\circ$ get 0p).
Inverse, improper, too wide or too narrow rings get 0 points for ring
continuity.

\ioaafig{2019_OBS-TEL_2-s5.pdf}
\ioaafig{2019_OBS-TEL_2-s6.pdf}

\medskip
\textbf{Alternative task: observing Epsilon Lyr.} \hfill TOTAL: 15 points

Checking the object in the FOV, and pointing is evaluated by the telescope
assistant. \hfill(0--3 points)

FOV with 10 mm eyepiece: the components of the close pairs are not resolved.

\ioaafig{2019_OBS-TEL_2-s7.pdf}

FOV drawing: star field and directions:
\begin{center}
\begin{tabular}{ll}
Correct direction \& labelling of North and East relative to the star field: & 2 points \\
Correct drawing of at least 3 stars in the FOV: & 1 point \\
Correct drawing of at least 6 stars in the FOV: & 2 points \\
\end{tabular}
\end{center}

Wide pair distance: $208'' = 3.47'$ \hfill 2 points

\begin{center}
\begin{tabular}{lll}
Distance estimation of wide pair: & $d_{\epsilon 1 - \epsilon 2} = 3.2$--$3.8'$ & 2 points \\
 & $d_{\epsilon 1 - \epsilon 2} = 2.8$--$4.2'$ & 1 point \\
\end{tabular}
\end{center}

Wide pair PA: $172^\circ$ \hfill 3 points

\begin{center}
\begin{tabular}{lll}
Position angle of the wide pair: & $\mathrm{PA}_{\epsilon 2} = 170$--$175^\circ$ & 3 points \\
 & $\mathrm{PA}_{\epsilon 2} = 167$--$177^\circ$ & 2 points \\
 & $\mathrm{PA}_{\epsilon 2} = 162$--$182^\circ$ & 1 point \\
\end{tabular}
\end{center}

Correct estimation of the relative angle of close pairs: \hfill 3 points

Relative angle between the direction of lines fitted onto the two close
pairs: $92^\circ$ (also the $88^\circ$ complementary angle is acceptable ---
referring to this, the evaluation bands are centered to $90^\circ$):
\begin{center}
\begin{tabular}{ll}
$85$--$95^\circ$ & 3 points \\
$80$--$100^\circ$ & 2 points \\
$75$--$105^\circ$ & 1 point \\
\end{tabular}
\end{center}

\solhead{15.70}{Telescope task 3}
% source id: 2019_OBS-TEL_3
Checking the object in the FOV, and pointing is evaluated by the telescope
assistant. \hfill(0--4 points)

The field of view (in the 25 mm eyepiece) drawn on the star chart is checked
against a transparent template foil showing the correct FOV circle around
M57, with dashed circles marking the $\pm 10$\% and $\pm 20$\% accuracy
bands.

\ioaafig{2019_OBS-TEL_3-s1.pdf}
\ioaafig{2019_OBS-TEL_3-s2.pdf}

\begin{center}
\begin{tabular}{ll}
Accuracy within $\pm 10$\%: & 6 points \\
Accuracy within $\pm 20$\%: & 3 points \\
\end{tabular}
\end{center}

\solhead{15.71}{Telescope task 4}
% source id: 2019_OBS-TEL_4
Checking the object in the FOV, and pointing is evaluated by the telescope
assistant. \hfill(0--8 points)

Brightness estimation:
\begin{center}
\begin{tabular}{ll}
Magnitude $\pm 0.2$ mag: & 6 points \\
Magnitude $\pm 0.3$ mag: & 4 points \\
Magnitude $\pm 0.4$ mag: & 2 points \\
\end{tabular}
\end{center}

Time in UTC (if corrected for Summer Time and Time Zone differences):
\hfill 1 point

Reference estimation: according to 3 independent visual amateur observations
during the night.

\solhead{15.72}{Telescope task 5}
% source id: 2019_OBS-TEL_5
Brightness estimation of STAR 2 ($\zeta$ UMi): \textbf{4.3} mag

Brightness estimation of STAR 1 ($\gamma$ UMi): \textbf{3.0} mag

\begin{center}
\begin{tabular}{ll}
Values within $\pm 0.2$ mag: & 2--2 points \\
Values within $\pm 0.3$ mag: & 1--1 point \\
\end{tabular}
\end{center}

Angular distance estimation between $\gamma$ UMi and Polaris:
$\mathbf{19^\circ}$

Values within $\pm 3^\circ$: \hfill 1 point

\solhead{15.73}{Night observation O1}
% source id: 2022_OBS-TEL_O1
Each correct answer is worth 3 pts.

\begin{center}
\begin{tabular}{lc|lc}
Object & Number & Object & Number \\
\hline
Jupiter & --- & $\beta$ Her / Antilicus & --- \\
% sic: the official solution table spells the name "Antilicus"; the question
% paper has "Antilicus" as well (the star is usually called Kornephoros)
Saturn & --- & $\alpha$ Oph / Rasalhague & --- \\
Mars & --- & $\epsilon$ Peg / Enif & 3 \\
$\alpha$ And / Alpheratz & --- & $\alpha$ Per / Mirfak & 1 \\
$\alpha$ Aql / Altair & --- & $\alpha$ Sco / Antares & 2 \\
$\alpha$ Boo / Arcturus & 4 & $\alpha$ Ser / Unukalhai & --- \\
$\alpha$ CrB / Alphecca & --- & $\epsilon$ Sgr / Kaus Australis & --- \\
$\beta$ Dra / Rastaban & --- & $\beta$ UMi / Kochab & 5 \\
\end{tabular}
\end{center}

\solhead{15.74}{Night observation O2}
% source id: 2022_OBS-TEL_O2
One method is to aim the telescope at an object, with easily recognizable
markings (such as a building) and scroll through several ($n$) fields of
view using only a single axis, while recording the difference between angles
($\Delta\alpha$) marked on the axis.

The final formula is:
\[
FOV = \frac{\Delta\alpha}{n}
\]

Full markings will be given for the range values of
$FOV = 0.6-1.2^\circ$, but only if accompanied with a working method.

\solhead{15.75}{Night observation O3}
% source id: 2022_OBS-TEL_O3
\[
F = f \cdot \frac{AFOV}{FOV}
\]
Full markings will be given for the range $f = \textbf{940--1875 mm}$, but
only if calculated using a correct formula and a value of $FOV$ from the
previous question.
% sic: the official solution writes "f" for the telescope focal length in
% the quoted range although the formula calls it F

\solhead{15.76}{Night observation O4}
% source id: 2022_OBS-TEL_O4
One method is to follow the line with a telescope to the end of the building
and remember the point. After that students could use rulers to measure $x$
and $y$ coordinate differences at an arm's length and calculate the angle.

Full markings will be given for the range $\textbf{39.5--47.5}^\circ$.

\solhead{15.77}{Observation round: the alien saucer screens}
% source id: 2023_OBS-TEL_1
\textbf{Observation 1: `Asteroid occultation'.}

Beginning of occultation: 23:03:31\\
End of occultation: 23:03:48

\begin{center}
\begin{tabular}{c|c|c|c}
Mid-time of occultation & $\pm$ error & Duration of occultation & $\pm$ error \\
\hline
23:03:39.5 & 0.4 s & 17 s & 0.2 s \\
\end{tabular}
\end{center}

Marking:
\begin{enumerate}
\item Mid-time $t$ of occultation:
  \begin{itemize}
  \item $23{:}03{:}39.3 \le t < 23{:}03{:}39.7$ $\Longrightarrow$ 4 points
  \item $23{:}03{:}39.0 \le t < 23{:}03{:}39.3$ or
    $23{:}03{:}39.7 < t \le 23{:}03{:}40.0$ $\Longrightarrow$ 3 points
  \item $23{:}03{:}38.0 \le t < 23{:}03{:}39.0$ or
    $23{:}03{:}40.0 < t \le 23{:}03{:}41.0$ $\Longrightarrow$ 1 point
  \item outside this range $\Longrightarrow$ 0 points
  \end{itemize}
\item Error in mid-time $\Delta t$:
  \begin{itemize}
  \item $0.1 < \Delta x \le 0.5$ $\Longrightarrow$ 3 points
  \item $0.5 < \Delta x \le 0.7$ $\Longrightarrow$ 2 points
  \item $0.7 < \Delta x \le 1.0$ $\Longrightarrow$ 1 point
  \item outside this range or missing $\Longrightarrow$ 0 points
  \end{itemize}
  % sic: the official solution uses "Delta x" in the bands for the error in
  % mid-time although the heading defines it as Delta t
\item Duration $x$ of occultation (should be to 1 s.f.):
  \begin{itemize}
  \item $16.8~\mathrm{s} \le x \le 17.2~\mathrm{s}$ $\Longrightarrow$ 4 points
  \item $16~\mathrm{s} \le x < 16.8~\mathrm{s}$ or
    $17.2~\mathrm{s} < x \le 18~\mathrm{s}$ $\Longrightarrow$ 3 points
  \item $15~\mathrm{s} \le x < 16~\mathrm{s}$ or
    $18~\mathrm{s} < x \le 19~\mathrm{s}$ $\Longrightarrow$ 1 point
  \item outside this range $\Longrightarrow$ 0 points
  \end{itemize}
\item Error in duration $\Delta x$:
  \begin{itemize}
  \item $0.05 < \Delta x \le 0.2$ $\Longrightarrow$ 3 points
  \item $0.2 < \Delta x \le 0.4$ $\Longrightarrow$ 2 points
  \item $0.4 < \Delta x \le 1.0$ $\Longrightarrow$ 1 point
  \item outside this range or missing $\Longrightarrow$ 0 points
  \end{itemize}
\item Error of duration lower than and different from error in mid time
  ($\Delta x < \Delta t$) $\Longrightarrow$ 1 point
\end{enumerate}

\medskip
\textbf{Observation 2: `Starlink'.}

The satellites are at an altitude $h = 20^\circ$ and their height above the
surface of the Earth is $H = 400$ km.

1) In the (simplified) solution neglecting the Earth's curvature, the
student will assume the distance to the satellites is equal to:
\[
d_{flat} = \frac{H}{\sin h} = 1170~\mathrm{km}
\]

\ioaafig{2023_OBS-TEL_1-s4.pdf}

2) In the solution taking into account the Earth's curvature, let us draw
the $OSC$ (observer--satellite--center of the Earth) triangle. The angles
are: $\angle COS = 90^\circ + h$, $\angle OSC$ that we shall denote as
$\eta$, and $\angle SCO = 180^\circ - (90^\circ + h) - \eta = 70^\circ - \eta$.

\ioaafig{2023_OBS-TEL_1-s5.pdf}

The law of sines can be applied:
\[
\frac{\sin(90^\circ + h)}{R_\oplus + H} = \frac{\sin\eta}{R_\oplus}
= \frac{\sin(70^\circ - \eta)}{d_{curve}}
\]
We calculate the distance to the satellites:
\[
\eta = \arcsin\left(\frac{R_\oplus}{R_\oplus + H}\sin(90^\circ + h)\right)
= 62^\circ.2
\]
\[
d_{curve} = R_\oplus \frac{\sin(70^\circ - \eta)}{\sin\eta} = 980~\mathrm{km}
\]

Knowing the radius of the orbit is $R_\oplus + H$, we calculate
$v_{orbit} = \sqrt{GM_\oplus/(R_\oplus + H)} = 7.6$ km/s. The observer can
only measure the tangential component of motion $v_t$.

The student will estimate the angular velocity to be
$v_t/d = v_{orbit}\sin h/d = 2.2\cdot 10^{-3}\,[1/\mathrm{s}] = 7.6\,['/\mathrm{s}]$
(for the solution neglecting the Earth's curvature) or
$2.6\cdot 10^{-3}\,[1/\mathrm{s}] = 9.0\,['/\mathrm{s}]$ (for the solution
taking into account the Earth's curvature).

\begin{flushright}
\textbf{5 points for predicting the angular velocity:}\\
5 points if within $8.55 - 9.45\,['/\mathrm{s}]$;\\
4 points if within $8.1 - 9.9\,['/\mathrm{s}]$;\\
3 points if within $7.5 - 10.5\,['/\mathrm{s}]$;\\
0 points otherwise
\end{flushright}

The simulated angular velocity is $8.5\,['/\mathrm{s}]$; the student should
conclude this value is similar to their prediction (or optionally comment it
is slightly lower and might imply e.g.\ the real orbit radius is minimally
larger than predicted, but it is not necessary).

\begin{flushright}
\textbf{4 points for measuring the angular velocity:}\\
4 points if agrees within 5\% ($8.1$ - $8.9$ $['/\mathrm{s}]$);\\
2 points if agrees within 10\% ($7.65$ - $9.35$ $['/\mathrm{s}]$);\\
0 points otherwise\\[4pt]
\textbf{1 point for correct comparison}
\end{flushright}

The student should then measure the time interval between the passes to be
$t = 2$ s.

\begin{flushright}
\textbf{2 points for measuring the time interval:}\\
2 points if agrees within 5\% ($1.9$ - $2.1$ s); 0 points otherwise
\end{flushright}

In such a short time, the path along the orbit can be approximated with a
straight line; the distance passed is $S = v_{orbit}t = 15$ km, which is the
distance between the satellites. (Alternatively, if the student chooses to
use their own $v_t$ measurement, $S = v_t t/\sin h = 14$ km).

\begin{flushright}
\textbf{2 points for calculating the distance}\\
2 points if within 13.8-15.2 km\\
1 point if within 12.5-16.5 km\\
0 points otherwise
\end{flushright}

Finally, the student should mark the observed satellite path on the sky map.

\begin{flushright}
\textbf{1 point for correctly marking the satellite trail}
\end{flushright}

\medskip
\textbf{Observation 3: `Planetary Moons'.}

The displayed planet is Saturn. On the reference solution image the moons
are marked and labelled with their names and magnitudes: Japetus (11.0),
Rhea (9.9), Mimas (13.0), Tethys (10.4), Enceladus (11.8), Dione (10.6) and
Titan (8.5) around the disk of Saturn.

\ioaafig{2023_OBS-TEL_1-s6.pdf}

2 points for each moon: position (1 pt) and name (1 pt).

\medskip
\textbf{Observation 4: `Supernova'.}

\begin{center}
\begin{tabular}{c|c|c}
Right ascension & Declination & est.\ magnitude \\
\hline
$9^h\,55^m\,54^s$ & $+69^\circ\,09'\,11''$ & 10.2 mag \\
\end{tabular}
\end{center}

\begin{enumerate}
\item Declination:
  \begin{itemize}
  \item[(a)] $< \pm 1.5'$ $\Longrightarrow$ 4 points
  \item[(b)] $\le \pm 3'$ $\Longrightarrow$ 2 points
  \item[(c)] $> \pm 3'$ $\Longrightarrow$ 0 points
  \end{itemize}
\item Right ascension:
  \begin{itemize}
  \item[(a)] $< \pm\,22.5$ s $\Longrightarrow$ 4 points
  \item[(b)] $\le \pm\,45$ s $\Longrightarrow$ 2 points
  \item[(c)] $> \pm\,45$ s $\Longrightarrow$ 0 points
  \end{itemize}
\item Magnitude:
  \begin{itemize}
  \item[(a)] $< \pm\,0.4$ mag $\Longrightarrow$ 2 points
  \item[(b)] $\le \pm\,0.8$ mag $\Longrightarrow$ 1 point
  \item[(c)] $> \pm\,0.8$ mag $\Longrightarrow$ 0 points
  \end{itemize}
\end{enumerate}

\solhead{15.78}{Discovery of 'Z'}
% source id: 2025_OBS-TEL_OT01
\textbf{(OT01.1)} The four stars are circled and labelled on the sky map:
S3 ($\beta$ Peg, $23^h3^m47^s$, $28^\circ4'58''$) at top, S1 ($\alpha$ And,
$0^h8^m24^s$, $29^\circ5'16''$) to its lower left on the map, S2
($\alpha$ Peg, $23^h4^m46^s$, $15^\circ12'17''$) at the centre and S4
($\gamma$ Peg, $0^h13^m14^s$, $15^\circ10'59''$) below.

\ioaafig{2025_OBS-TEL_OT01-s2.pdf}

\begin{center}
\begin{tabular}{l r}
\hline
Circle around each correct star & 0.5 \\
Correct identification as S1--S4 of each star & 1.0 \\
\hline
\end{tabular}
\end{center}

\textbf{(OT01.2)}
\begin{itemize}
\item Realize that two stars have almost the same RA and two of them have
  almost same declination.
\item We draw a line through them indicating constant RA and constant
  Declination. (On the official overlay the construction lines are labelled
  $\delta \approx 15^\circ$, $\alpha \approx 23^h$, and their parallels
  through Z are labelled $\alpha \approx 21^h36^m$ and
  $\delta \approx -26^\circ$.)
\item Given the scale of the map and the separation between the 4 stars and
  the new object we can assume that the celestial grid will be almost
  linear.
\item We can thus mark `Z' at its given coordinates
  ($\alpha = 21^h36^m$; $\delta = -26^\circ10'$).
\end{itemize}

\begin{center}
\begin{tabular}{l r}
\hline
Object marked within inner rectangle & 7.0 \\
\quad Object marked between outer and inner rectangles & 4.0 \\
\hline
\end{tabular}
\end{center}

\textbf{(OT01.3)}
\begin{itemize}
\item After marking the new object on the given sky map, we carefully look
  at its position and the star pattern surrounding it.
\item Using the 25 mm focal length eyepiece we try to locate it.
\item By doing some fine adjustments, we try to bring the same pattern in
  the 10 mm focal length eye piece (which has a crosshair).
\end{itemize}

\begin{center}
\begin{tabular}{l r}
\hline
Pointing towards the screen & 1.0 \\
\hline
\multicolumn{2}{l}{If 10 mm objective is used:} \\
% sic: the official marking table says "objective" for the eyepiece
\quad Correct choice of eyepiece & 2.0 \\
\quad Object is anywhere in FOV & 3.0 \\
\quad Any part of the object is touching the origin of crosshair & 6.0 \\
\hline
\multicolumn{2}{l}{If 25 mm objective is used:} \\
\quad Object is anywhere in FOV & 3.0 \\
\hline
\end{tabular}
\end{center}

\solhead{15.79}{Lensing Time Delay}
% source id: 2025_OBS-TEL_OT02
The solution is $t_{\mathrm{DA}} = 1694$ d, $t_{\mathrm{DB}} = 1651$ d and
$t_{\mathrm{DC}} = 2485$ d.

\begin{center}
\begin{tabular}{l r}
\hline
Pointing towards the screen & 1.0 \\
Choice of correct eyepiece & 1.0 \\
\hline
\multicolumn{2}{l}{For each of images A, B, C, D:} \\
\quad Timings of image & 2.0 \\
\qquad At least one reading & 1.0 \\
\qquad Extra reading (0.5 each, maximum 1.0) & 1.0 \\
\hline
\multicolumn{2}{l}{Calculation of time differences:} \\
\quad Calculation each of $t_{\mathrm{DA}}$, $t_{\mathrm{DB}}$ and
  $t_{\mathrm{DC}}$ & 3.0 \\
\qquad full credit for value $-25$ d to value $+25$ d; & \\
\qquad half credit down to value $-50$ d or up to value $+50$ d & \\
\quad Scaling from seconds to days for each interval & 1.0 \\
\hline
Correct order ($t_{\mathrm{DC}} > t_{\mathrm{DA}} > t_{\mathrm{DB}}$) & 3.0 \\
\hline
\end{tabular}
\end{center}

\solhead{15.80}{Saturn}
% source id: 2026_OBS-TEL_OT01
\textbf{(a)} Seen in the field of view using 10\,mm eyepiece. \hfill(10 points)

Seen in the field of view using 25\,mm eyepiece. \hfill(5 points)

\textbf{(b)} Write down the observing time. \hfill(1 point)

Check the results in the table below: within $\pm 5^\circ$ \hfill(4 points)

If out of this range but within $\pm 10^\circ$ \hfill(2 points)

\textbf{(c)} Write down the observing time. \hfill(1 point)

Check the results in the table below: within $\pm 5^\circ$ \hfill(10 points)

If out of this range but within $\pm 10^\circ$ \hfill(5 points)

\textbf{(d)} Write down observing time. \hfill(1 point)

Consider a right spherical triangle whose right angle is at the point of
Saturn's upper culmination. One leg is the segment of Saturn's diurnal path
from that culmination point to its current position; because it lies parallel
to and only $2^\circ$ away from the celestial equator, the problem
approximates it as a great circle passing through the Earth's center like the
celestial equator, thereby indirectly treating its length as the required hour
angle ($HA$). The other leg is the distance from that culmination point to the
zenith, having a length of $(\varphi - \delta)$, while the hypotenuse is
Saturn's zenith distance ($z$) from the previous question. \hfill(9 points)

From this, the formula for a right spherical triangle can be applied:
\[ \cos z = \cos HA \cos(\varphi - \delta_{\mathrm{Saturn}}) \]

Check the results in the table below (available in two formats): within
$\pm 10^{\mathrm{m}}$. \textit{Note: Adjust this value based on the previous
answer of zenith distance.} \hfill(4 points)

If out of this range but within $\pm 20^{\mathrm{m}}$ \hfill(2 points)

The solution values are based on the observing time.
% editorial: in the official table the zenith distance, axis angle and
% hour-angle hour entries are merged cells spanning several observing times;
% here every observing time gets its own row with the merged value repeated.
% Where a merged cell continues across a page break the official table leaves
% its continuation blank, so the following values are inferred from the cell
% begun on the previous page: zenith distance 44 for 21 h 29 m, 25 for 23 h 05 m
% and 23 h 07 m, and 19 for 0 h 06 m to 0 h 19 m (likewise the hour-angle
% hours in the first rows of each page).

\begin{center}
\small
\begin{longtable}{rr c c rr rr}
\toprule
\multicolumn{2}{c}{Observing time} & 1b) Zenith & 1c) Acute angle of rotation & \multicolumn{4}{c}{1d) Hour angle} \\
h & m & distance ($^\circ$) & axis relative to the horizon ($^\circ$) & h & m & h & m \\
\midrule
\endfirsthead
\toprule
\multicolumn{2}{c}{Observing time} & 1b) Zenith & 1c) Acute angle of rotation & \multicolumn{4}{c}{1d) Hour angle} \\
h & m & distance ($^\circ$) & axis relative to the horizon ($^\circ$) & h & m & h & m \\
\midrule
\endhead
\midrule
\multicolumn{8}{r}{\textit{continued on next page}} \\
\endfoot
\bottomrule
\endlastfoot
20 & 0 & 64 & 19 & $-4$ & 11 & 19 & 49 \\
20 & 2 & 63 & 19 & $-4$ & 9 & 19 & 51 \\
20 & 7 & 62 & 20 & $-4$ & 4 & 19 & 56 \\
20 & 11 & 61 & 20 & $-4$ & 0 & 20 & 0 \\
20 & 15 & 60 & 20 & $-3$ & 56 & 20 & 4 \\
20 & 20 & 59 & 20 & $-3$ & 51 & 20 & 9 \\
20 & 24 & 58 & 20 & $-3$ & 47 & 20 & 13 \\
20 & 28 & 57 & 21 & $-3$ & 43 & 20 & 17 \\
20 & 33 & 56 & 21 & $-3$ & 38 & 20 & 22 \\
20 & 37 & 55 & 21 & $-3$ & 34 & 20 & 26 \\
20 & 42 & 54 & 21 & $-3$ & 29 & 20 & 31 \\
20 & 44 & 54 & 22 & $-3$ & 27 & 20 & 33 \\
20 & 46 & 53 & 22 & $-3$ & 25 & 20 & 35 \\
20 & 50 & 52 & 22 & $-3$ & 21 & 20 & 39 \\
20 & 55 & 51 & 22 & $-3$ & 16 & 20 & 44 \\
20 & 58 & 51 & 23 & $-3$ & 13 & 20 & 47 \\
20 & 59 & 50 & 23 & $-3$ & 12 & 20 & 48 \\
21 & 4 & 49 & 23 & $-3$ & 7 & 20 & 53 \\
21 & 8 & 48 & 23 & $-3$ & 3 & 20 & 57 \\
21 & 9 & 48 & 24 & $-3$ & 2 & 20 & 58 \\
21 & 13 & 47 & 24 & $-2$ & 58 & 21 & 2 \\
21 & 17 & 46 & 24 & $-2$ & 54 & 21 & 6 \\
21 & 20 & 46 & 25 & $-2$ & 51 & 21 & 9 \\
21 & 22 & 45 & 25 & $-2$ & 49 & 21 & 11 \\
21 & 26 & 44 & 25 & $-2$ & 45 & 21 & 15 \\
21 & 29 & 44 & 26 & $-2$ & 42 & 21 & 18 \\
21 & 31 & 43 & 26 & $-2$ & 40 & 21 & 20 \\
21 & 35 & 42 & 26 & $-2$ & 36 & 21 & 24 \\
21 & 37 & 42 & 27 & $-2$ & 34 & 21 & 26 \\
21 & 40 & 41 & 27 & $-2$ & 31 & 21 & 29 \\
21 & 45 & 40 & 28 & $-2$ & 26 & 21 & 34 \\
21 & 49 & 39 & 28 & $-2$ & 22 & 21 & 38 \\
21 & 51 & 39 & 29 & $-2$ & 20 & 21 & 40 \\
21 & 54 & 38 & 29 & $-2$ & 17 & 21 & 43 \\
21 & 58 & 38 & 30 & $-2$ & 13 & 21 & 47 \\
21 & 59 & 37 & 30 & $-2$ & 12 & 21 & 48 \\
22 & 3 & 37 & 31 & $-2$ & 8 & 21 & 52 \\
22 & 4 & 36 & 31 & $-2$ & 7 & 21 & 53 \\
22 & 8 & 35 & 31 & $-2$ & 3 & 21 & 57 \\
22 & 9 & 35 & 32 & $-2$ & 2 & 21 & 58 \\
22 & 13 & 34 & 32 & $-1$ & 58 & 22 & 2 \\
22 & 14 & 34 & 33 & $-1$ & 57 & 22 & 3 \\
22 & 18 & 33 & 33 & $-1$ & 53 & 22 & 7 \\
22 & 19 & 33 & 34 & $-1$ & 52 & 22 & 8 \\
22 & 23 & 32 & 35 & $-1$ & 48 & 22 & 12 \\
22 & 27 & 32 & 36 & $-1$ & 44 & 22 & 16 \\
22 & 28 & 31 & 36 & $-1$ & 43 & 22 & 17 \\
22 & 31 & 31 & 37 & $-1$ & 40 & 22 & 20 \\
22 & 34 & 30 & 37 & $-1$ & 37 & 22 & 23 \\
22 & 35 & 30 & 38 & $-1$ & 36 & 22 & 24 \\
22 & 39 & 29 & 39 & $-1$ & 32 & 22 & 28 \\
22 & 42 & 29 & 40 & $-1$ & 29 & 22 & 31 \\
22 & 44 & 28 & 40 & $-1$ & 27 & 22 & 33 \\
22 & 45 & 28 & 41 & $-1$ & 26 & 22 & 34 \\
22 & 49 & 28 & 42 & $-1$ & 22 & 22 & 38 \\
22 & 50 & 27 & 42 & $-1$ & 21 & 22 & 39 \\
22 & 52 & 27 & 43 & $-1$ & 19 & 22 & 41 \\
22 & 54 & 27 & 44 & $-1$ & 17 & 22 & 43 \\
22 & 56 & 26 & 44 & $-1$ & 15 & 22 & 45 \\
22 & 57 & 26 & 45 & $-1$ & 14 & 22 & 46 \\
23 & 0 & 26 & 46 & $-1$ & 11 & 22 & 49 \\
23 & 2 & 25 & 47 & $-1$ & 9 & 22 & 51 \\
23 & 5 & 25 & 48 & $-1$ & 6 & 22 & 54 \\
23 & 7 & 25 & 49 & $-1$ & 4 & 22 & 56 \\
23 & 8 & 24 & 49 & $-1$ & 3 & 22 & 57 \\
23 & 10 & 24 & 50 & $-1$ & 1 & 22 & 59 \\
23 & 12 & 24 & 51 & $-0$ & 59 & 23 & 1 \\
23 & 14 & 24 & 52 & $-0$ & 57 & 23 & 3 \\
23 & 15 & 23 & 52 & $-0$ & 56 & 23 & 4 \\
23 & 16 & 23 & 53 & $-0$ & 55 & 23 & 5 \\
23 & 18 & 23 & 54 & $-0$ & 53 & 23 & 7 \\
23 & 21 & 23 & 55 & $-0$ & 50 & 23 & 10 \\
23 & 22 & 22 & 56 & $-0$ & 49 & 23 & 11 \\
23 & 24 & 22 & 57 & $-0$ & 47 & 23 & 13 \\
23 & 26 & 22 & 58 & $-0$ & 45 & 23 & 15 \\
23 & 28 & 22 & 59 & $-0$ & 43 & 23 & 17 \\
23 & 30 & 21 & 60 & $-0$ & 41 & 23 & 19 \\
23 & 32 & 21 & 61 & $-0$ & 39 & 23 & 21 \\
23 & 34 & 21 & 62 & $-0$ & 37 & 23 & 23 \\
23 & 35 & 21 & 63 & $-0$ & 36 & 23 & 24 \\
23 & 37 & 21 & 64 & $-0$ & 34 & 23 & 26 \\
23 & 39 & 21 & 65 & $-0$ & 32 & 23 & 28 \\
23 & 40 & 20 & 66 & $-0$ & 31 & 23 & 29 \\
23 & 42 & 20 & 67 & $-0$ & 29 & 23 & 31 \\
23 & 43 & 20 & 68 & $-0$ & 28 & 23 & 32 \\
23 & 45 & 20 & 69 & $-0$ & 26 & 23 & 34 \\
23 & 47 & 20 & 70 & $-0$ & 24 & 23 & 36 \\
23 & 48 & 20 & 71 & $-0$ & 23 & 23 & 37 \\
23 & 50 & 20 & 72 & $-0$ & 21 & 23 & 39 \\
23 & 51 & 20 & 73 & $-0$ & 20 & 23 & 40 \\
23 & 53 & 19 & 74 & $-0$ & 18 & 23 & 42 \\
23 & 54 & 19 & 75 & $-0$ & 17 & 23 & 43 \\
23 & 56 & 19 & 76 & $-0$ & 15 & 23 & 45 \\
23 & 57 & 19 & 77 & $-0$ & 14 & 23 & 46 \\
23 & 58 & 19 & 78 & $-0$ & 13 & 23 & 47 \\
0 & 0 & 19 & 79 & $-0$ & 11 & 23 & 49 \\
0 & 1 & 19 & 80 & $-0$ & 10 & 23 & 50 \\
0 & 3 & 19 & 81 & $-0$ & 8 & 23 & 52 \\
0 & 4 & 19 & 82 & $-0$ & 7 & 23 & 53 \\
0 & 6 & 19 & 83 & $-0$ & 5 & 23 & 55 \\
0 & 7 & 19 & 84 & $-0$ & 4 & 23 & 56 \\
0 & 8 & 19 & 85 & $-0$ & 3 & 23 & 57 \\
0 & 10 & 19 & 86 & $-0$ & 1 & 23 & 59 \\
0 & 11 & 19 & 87 & 0 & 0 & 0 & 0 \\
0 & 13 & 19 & 88 & 0 & 2 & 0 & 2 \\
0 & 14 & 19 & 89 & 0 & 3 & 0 & 3 \\
0 & 15 & 19 & 90 & 0 & 4 & 0 & 4 \\
0 & 17 & 19 & 89 & 0 & 6 & 0 & 6 \\
0 & 18 & 19 & 88 & 0 & 7 & 0 & 7 \\
0 & 19 & 19 & 87 & 0 & 8 & 0 & 8 \\
\end{longtable}
\end{center}

\solhead{15.81}{Messier 29}
% source id: 2026_OBS-TEL_OT02
Seen in the field of view using 10\,mm eyepiece. \hfill(15 points)

Seen in the field of view using 25\,mm eyepiece. \hfill(10 points)

Not seen in the field of view. \hfill(0 points)

\solhead{15.82}{Reverberation Mapping}
% source id: 2026_OBS-TEL_OT03
\textbf{(a)} The ring and central dot flashes seen in the field of view.
\hfill(5 points)

Not seen in the field of view. \hfill(0 points)

\textbf{(b)} Recorded 3 to 5 individual measurements of $\tau$. (Only 1--2
readings: 1 point) \hfill(5 points)

$14.5\,\mathrm{s} \le \tau_{\mathrm{mean}} \le 15.5\,\mathrm{s}$ and
$\Delta\tau \le 0.5\,\mathrm{s}$ \hfill(5 points)
\begin{itemize}
\item \textit{If fails the condition above, but satisfies
  $14.0 \le \tau_{\mathrm{mean}} \le 16.0\,\mathrm{s}$ and
  $\Delta\tau \le 1.0\,\mathrm{s}$ awards 3 points.}
\item \textit{If the mean ($\tau_{\mathrm{mean}}$) falls within a valid range
  above but the uncertainty ($\Delta\tau$) is missing, deduct half of the
  points earned.}
\end{itemize}

Out of the above range. \hfill(0 points)

\textbf{(c)} The ring radius:
\[ R = c\tau = 3.88 \times 10^{14}\,\mathrm{m} \]
\hfill(2 points)

The black hole's mass:
\[ F_g = F_c \rightarrow \frac{GmM}{R^2} = \frac{m(\Delta V)^2}{R}
   \rightarrow M = \frac{R(\Delta V)^2}{G} \approx 2.6 \times 10^{7}\,\mathrm{M}_\odot \]
\hfill(3 points)
